\documentclass[a4paper,12pt]{report}
\usepackage[utf8]{inputenc}


% Deutsche Bezeichnungen ohne babel-german
\renewcommand{\contentsname}{Inhaltsverzeichnis}
\renewcommand{\chaptername}{Kapitel}
\renewcommand{\figurename}{Abbildung}
\renewcommand{\tablename}{Tabelle}
\renewcommand{\bibname}{Literaturverzeichnis}
\renewcommand{\indexname}{Stichwortverzeichnis}
\renewcommand{\listfigurename}{Abbildungsverzeichnis}
\renewcommand{\listtablename}{Tabellenverzeichnis}
\renewcommand{\appendixname}{Anhang}
\renewcommand{\abstractname}{Zusammenfassung}
\usepackage{amsmath}
\usepackage{amssymb}
\usepackage{amsthm}
\usepackage{geometry}
\usepackage{graphicx}
\usepackage{enumitem}
\usepackage[protrusion=true,expansion=true,tracking=false,kerning=true,spacing=true]{microtype}
\usepackage[hidelinks]{hyperref}
\usepackage{tikz}
\usepackage[ruled,vlined,linesnumbered,german]{algorithm2e}
\SetAlgorithmName{Algorithmus}{algorithmus}{Algorithmenverzeichnis}

\usetikzlibrary{arrows.meta,positioning,calc}
\hypersetup{colorlinks=true,linkcolor=blue!60!black,urlcolor=blue!60!black,citecolor=blue!60!black}
\geometry{margin=2.5cm}
\setlist[itemize,1]{label=--}

\theoremstyle{definition}
\newtheorem{definition}{Definition}
\newtheorem{beispiel}{Beispiel}
\newtheorem{satz}{Satz}
\newtheorem{lemma}{Lemma}
\newtheorem{korollar}{Korollar}
\newtheorem{beobachtung}{Beobachtung}

\title{\textbf{Lineare Algebra} \\ {\Large Strukturen, Räume und Abbildungen}}
\author{Stephan Epp}
\date{02.\ Juni 2026}

\begin{document}

\renewcommand{\proofname}{Beweis}
\maketitle
\tableofcontents
\newpage

% ============================================================
\chapter{Mengen und Abbildungen}
% ============================================================

\section{Mengen}

Mengen sind das sprachliche Fundament der Mathematik. Cantor fasst eine Menge als
eine Zusammenfassung bestimmter, wohlunterschiedener Objekte zu einem Ganzen auf.
Diese Beschreibung ist intuitiv klar, aber nicht exakt -- der Begriff \emph{Menge}
wird durch den Begriff \emph{Zusammenfassung} erklärt, der selbst undefiniert bleibt.
Wir nehmen den naiven Standpunkt ein: Eine Menge ist definiert, wenn feststeht,
welche Objekte ihr angehören.

\begin{definition}[Menge]
    Eine \textbf{Menge} $M$ ist bestimmt durch ihre \textbf{Elemente}. Ist $a$ ein
    Element von $M$, so schreiben wir $a \in M$. Ist $a$ kein Element, so schreiben
    wir $a \notin M$.

    Grundlegende Konstruktionen:
    \begin{itemize}
        \item[-] \textbf{Teilmenge}: $N \subseteq M$, falls jedes $n \in N$ auch in $M$ liegt.
        \item[-] \textbf{Echte Teilmenge}: $N \subset M$, falls $N \subseteq M$ und
              $N \neq M$.
        \item[-] \textbf{Leere Menge}: $\emptyset$ enthält kein Element; sie ist
              Teilmenge jeder Menge.
        \item[-] \textbf{Potenzmenge}: $\mathcal{P}(M)$ ist die Menge aller
              Teilmengen von $M$.
    \end{itemize}
\end{definition}

\begin{definition}[Mengenoperationen]
    Seien $N_1, N_2 \subseteq M$.
    \begin{itemize}
        \item[-] \textbf{Vereinigung}: $N_1 \cup N_2 = \{m \mid m \in N_1
              \text{ oder } m \in N_2\}$
        \item[-] \textbf{Durchschnitt}: $N_1 \cap N_2 = \{m \mid m \in N_1
              \text{ und } m \in N_2\}$
        \item[-] \textbf{Komplement}: $\overline{N} = \{m \mid m \in M,\, m \notin N\}$
        \item[-] \textbf{Differenz}: $N_1 \setminus N_2 = \{n \mid n \in N_1,\,
              n \notin N_2\}$
        \item[-] \textbf{Kartesisches Produkt}: $M_1 \times \dots \times M_k =
              \{(m_1,\dots,m_k) \mid m_j \in M_j\}$
    \end{itemize}
\end{definition}

\begin{satz}[Rechenregeln für Mengen]
    Seien $A, B, C, N_j$ Teilmengen einer Menge $M$. Dann gilt:
    \begin{itemize}
        \item[-] $A \cup B = B \cup A$ und $A \cap B = B \cap A$ \quad (Kommutativität)
        \item[-] $A \cup (B \cup C) = (A \cup B) \cup C$ \quad (Assoziativität)
        \item[-] $A \cap \bigl(\bigcup_{j \in J} N_j\bigr) = \bigcup_{j \in J}
              (A \cap N_j)$ \quad (Distributivität)
        \item[-] $\overline{\bigcup_{j \in J} N_j} = \bigcap_{j \in J} \overline{N_j}$
              \quad (de Morgansche Regeln)
    \end{itemize}
\end{satz}
\begin{proof}
    Die Kommutativitätsgesetze folgen unmittelbar aus der symmetrischen Definition:
    $A \cup B = \{x \mid x \in A \text{ oder } x \in B\} = B \cup A$.
    Die Distributivität zeigt man durch doppelte Inklusion:
    Ist $x \in A \cap \bigl(\bigcup_j N_j\bigr)$, so gilt $x \in A$ und $x \in N_j$
    für ein $j$, also $x \in A \cap N_j \subseteq \bigcup_j (A \cap N_j)$.
    Die umgekehrte Inklusion ist analog. Die de Morganschen Regeln folgen aus
    $x \notin \bigcup_j N_j \Leftrightarrow x \notin N_j$ für alle $j
    \Leftrightarrow x \in \bigcap_j \overline{N_j}$.
\end{proof}


\begin{definition}[Äquivalenzrelation]
    Eine \textbf{Relation} auf $M$ ist eine Teilmenge $R \subseteq M \times M$.
    Wir schreiben $m \sim m'$ statt $(m, m') \in R$.

    Eine Äquivalenzrelation erfüllt:
    \begin{itemize}
        \item[-] \textbf{Reflexivität}: $m \sim m$ für alle $m \in M$
        \item[-] \textbf{Symmetrie}: Aus $m \sim m'$ folgt $m' \sim m$
        \item[-] \textbf{Transitivität}: Aus $m \sim m'$ und $m' \sim m''$ folgt $m \sim m''$
    \end{itemize}

    Die \textbf{Äquivalenzklasse} von $m$ ist $[m] = \{m' \mid m' \sim m\}$.
    Die Äquivalenzklassen bilden eine \textbf{Partition} von $M$: jedes Element
    liegt in genau einer Klasse.
\end{definition}

\begin{beobachtung}[Partition und Äquivalenz]
    Äquivalenzrelationen und Partitionen beschreiben dasselbe Phänomen aus
    zwei Perspektiven: Eine Äquivalenzrelation \emph{zerlegt} eine Menge in
    disjunkte Klassen; eine Partition \emph{definiert} umgekehrt eine
    Äquivalenzrelation durch $m \sim m'$ genau dann, wenn $m$ und $m'$ in
    derselben Klasse liegen.

    Das kanonische Beispiel: Auf $\mathbb{Z}$ ist $n_1 \equiv n_2 \pmod{m}$
    eine Äquivalenzrelation. Die Klassen sind die Restklassen
    $[r] = r + m\mathbb{Z}$ mit $0 \le r < m$.
\end{beobachtung}

% ============================================================
\section{Abbildungen}
% ============================================================

\begin{definition}[Abbildung]
    Eine \textbf{Abbildung} $f : M \to N$ ordnet jedem $m \in M$ genau ein
    $n = f(m) \in N$ zu. Mit $\mathrm{Ab}(M, N)$ bezeichnen wir die Menge aller
    Abbildungen von $M$ nach $N$.

    Eigenschaften:
    \begin{itemize}
        \item[-] \textbf{Injektiv}: Aus $f(m_1) = f(m_2)$ folgt $m_1 = m_2$ --
              verschiedene Eingaben liefern verschiedene Ausgaben.
        \item[-] \textbf{Surjektiv}: $f(M) = N$ -- jedes $n \in N$ wird getroffen.
        \item[-] \textbf{Bijektiv}: Injektiv und surjektiv -- eine
              Eins-zu-eins-Korrespondenz.
    \end{itemize}
    Das \textbf{Bild} von $U \subseteq M$ ist $f(U) = \{f(u) \mid u \in U\}$,
    das \textbf{Urbild} von $V \subseteq N$ ist
    $f^{-1}(V) = \{m \in M \mid f(m) \in V\}$.
\end{definition}

\begin{satz}[Bijektive Abbildungen]
    Sei $f \in \mathrm{Ab}(M, N)$.
    \begin{itemize}
        \item[-] $f$ ist genau dann injektiv, wenn es ein $g \in \mathrm{Ab}(N, M)$
              gibt mit $g \circ f = \mathrm{id}_M$.
        \item[-] $f$ ist genau dann surjektiv, wenn es ein $h \in \mathrm{Ab}(N, M)$
              gibt mit $f \circ h = \mathrm{id}_N$.
        \item[-] $f$ ist genau dann bijektiv, wenn es ein eindeutiges $g \in \mathrm{Ab}(N, M)$
              gibt mit $g \circ f = \mathrm{id}_M$ und $f \circ g = \mathrm{id}_N$.
              Man nennt $g = f^{-1}$ die \textbf{Inverse} von $f$.
    \end{itemize}
\end{satz}
\begin{proof}
    Ist $f$ bijektiv, so ist $f^{-1}$ wohldefiniert durch $f^{-1}(y) = x$, wobei
    $x$ das eindeutige Urbild von $y$ ist. Dann gilt $f \circ f^{-1} = \mathrm{id}_N$
    und $f^{-1} \circ f = \mathrm{id}_M$ per Definition.
    Sind umgekehrt $g, h$ Links- bzw. Rechtsinverses, so gilt
    $g = g \circ \mathrm{id} = g \circ (f \circ h) = (g \circ f) \circ h
    = \mathrm{id} \circ h = h$,
    also sind Links- und Rechtsinverses gleich.
\end{proof}


\begin{beobachtung}[Abzählbarkeit und Kardinalität]
    Zwei Mengen heißen \textbf{gleichmächtig}, wenn es eine Bijektion zwischen
    ihnen gibt. Eine Menge heißt \textbf{endlich} mit $|M| = n$, falls eine
    Bijektion $M \to \{1, \dots, n\}$ existiert, und \textbf{abzählbar}, falls
    eine Bijektion $M \to \mathbb{N}$ existiert.

    Es gilt: $|\mathbb{N}| = |\mathbb{Z}| = |\mathbb{Q}|$, aber $|\mathbb{R}| > |\mathbb{N}|$.
    Die Menge $\mathbb{R}$ ist nicht abzählbar -- dies zeigt Cantors
    Diagonalargument.
\end{beobachtung}

% ============================================================
\section{Binomialkoeffizienten und Abzählungen}
% ============================================================

\begin{definition}[Binomialkoeffizient]
    Für eine endliche Menge $M$ mit $|M| = m \ge 0$ und $0 \le k \le m$
    bezeichnen wir mit $\binom{m}{k}$ die Anzahl der Teilmengen von $M$ mit
    genau $k$ Elementen. Es gilt:
    \[
        \binom{m}{0} = \binom{m}{m} = 1, \qquad
        \binom{m+1}{k} = \binom{m}{k} + \binom{m}{k-1}, \qquad
        \binom{m}{k} = \frac{m!}{k!\,(m-k)!}.
    \]
\end{definition}

\begin{satz}[Binomischer Lehrsatz]
    Seien $a, b \in \mathbb{R}$ und $m \in \mathbb{N}$. Dann gilt:
    \[
        (a + b)^m = \sum_{j=0}^{m} \binom{m}{j} a^j b^{m-j}.
    \]
\end{satz}
\begin{proof}
    Durch vollständige Induktion nach $n$. Für $n = 0$ gilt $(a+b)^0 = 1 = \binom{0}{0}a^0 b^0$.
    Im Induktionsschritt gilt:
    \[
        (a+b)^{n+1} = (a+b)(a+b)^n = (a+b)\sum_{k=0}^n \binom{n}{k} a^k b^{n-k}.
    \]
    Sortieren nach Potenzen und Anwenden der Rekursion
    $\binom{n}{k} + \binom{n}{k-1} = \binom{n+1}{k}$ (Pascalsches Dreieck)
    liefert das Ergebnis.
\end{proof}


\begin{satz}[Inklusions-Exklusions-Prinzip]
    Seien $M_1, \dots, M_n$ endliche Teilmengen einer Menge $M$. Dann:
    \[
        |M_1 \cup \dots \cup M_n|
        = \sum_{k=1}^{n} (-1)^{k-1}
        \sum_{1 \le j_1 < \dots < j_k \le n} |M_{j_1} \cap \dots \cap M_{j_k}|.
    \]
\end{satz}
\begin{proof}
    Durch Induktion nach $n$. Für $n = 1$ ist nichts zu zeigen. Der Induktionsschritt
    von $n$ auf $n+1$ folgt aus $|A \cup B| = |A| + |B| - |A \cap B|$ (Anfangsfall $n=2$,
    der seinerseits aus dem Zählen jedes Elements folgt) und Anwenden der Formel
    auf $A_1 \cup \cdots \cup A_n$ sowie Ãœbergang zu $B = A_1 \cup \cdots \cup A_n$
    und $A_{n+1}$.
\end{proof}


\begin{satz}[Anzahl surjektiver Abbildungen]
    Seien $M$ und $N$ endliche Mengen mit $|M| = m$ und $|N| = n$. Die Anzahl
    der surjektiven Abbildungen von $M$ auf $N$ ist
    \[
        s(n, m) = \sum_{k=0}^{n-1} (-1)^k \binom{n}{k} (n-k)^m.
    \]
    Die Zahlen $s(n, m)$ heißen \textbf{Stirling-Zahlen}.
\end{satz}
\begin{proof}
    Sei $S(m, n)$ die Anzahl. Nach dem Inklusions-Exklusions-Prinzip gilt:
    \[
        S(m, n) = \sum_{k=0}^n (-1)^k \binom{n}{k} (n-k)^m,
    \]
    denn man zieht von $n^m$ (alle Abbildungen) sukzessive jene ab, die mindestens
    ein Element nicht treffen, addiert jene zurück, die mindestens zwei Elemente
    nicht treffen, usw. Dies ergibt genau die angegebene Formel.
\end{proof}


% ============================================================
\chapter{Vektorräume}
% ============================================================

\section{Gruppen}

\begin{definition}[Gruppe]
    Sei $G$ eine nichtleere Menge. $G$ heißt eine \textbf{Gruppe}, falls:
    \begin{itemize}
        \item[-] Jedem Paar $(a, b) \in G \times G$ ist eindeutig ein Element
              $c = a \circ b \in G$ zugeordnet.
        \item[-] Es gilt das \textbf{Assoziativgesetz}:
              $(a \circ b) \circ c = a \circ (b \circ c)$.
        \item[-] Es gibt ein \textbf{neutrales Element} $e \in G$ mit $ea = a$
              für alle $a \in G$.
        \item[-] Zu jedem $a \in G$ gibt es ein \textbf{inverses Element}
              $a^{-1} \in G$ mit $a^{-1} a = e$.
    \end{itemize}
    Gilt zusätzlich $a \circ b = b \circ a$ für alle $a, b \in G$, so heißt
    $G$ \textbf{abelsch} (kommutativ). Wir schreiben in abelschen Gruppen meist
    $+$ statt $\circ$ und $0$ statt $e$.
\end{definition}

\begin{beobachtung}[Was eine Gruppe leistet]
    Das Gruppenaxiom ist minimalistisch: Es verlangt nur Assoziativität,
    neutrales Element und Inverse -- kein Kommutativgesetz. Dennoch folgen
    daraus bereits starke Aussagen:
    \begin{itemize}
        \item[-] Das neutrale Element ist \emph{eindeutig}.
        \item[-] Das Inverse eines Elements ist \emph{eindeutig}.
        \item[-] $(ab)^{-1} = b^{-1} a^{-1}$ -- die Reihenfolge kehrt sich um.
        \item[-] Die Kürzungsregeln gelten: Aus $ab = ac$ folgt $b = c$.
    \end{itemize}
    Die Struktur ist gerade reich genug, um Symmetrien und Transformationen
    zu beschreiben, aber arm genug, um universell einsetzbar zu sein.
\end{beobachtung}

\begin{definition}[Untergruppe und Nebenklassen]
    Sei $G$ eine Gruppe und $U \subseteq G$ nichtleer. $U$ heißt
    \textbf{Untergruppe} von $G$ ($U \le G$), falls:
    \begin{itemize}
        \item[-] $u_1, u_2 \in U \Rightarrow u_1 u_2 \in U$
        \item[-] $u \in U \Rightarrow u^{-1} \in U$
    \end{itemize}
    Für $g \in G$ heißt $gU = \{gu \mid u \in U\}$ eine \textbf{Linksnebenklasse}
    von $U$ in $G$. Sie sind entweder gleich oder disjunkt und bilden eine
    Partition von $G$.
\end{definition}

\begin{satz}[Lagrange]
    Sei $G$ eine endliche Gruppe und $U \le G$. Dann gilt $|G| = |U| \cdot [G:U]$,
    wobei $[G:U]$ die Anzahl der Nebenklassen ist. Insbesondere ist $|U|$ ein
    Teiler von $|G|$.
\end{satz}
\begin{proof}
    Nach dem Satz über Nebenklassenzerlegungen gilt $G = \bigsqcup_{j \in J} g_j U$
    als disjunkte Vereinigung von Linksnebenklassen. Da die Abbildung
    $u \mapsto g_j u$ eine Bijektion von $U$ auf $g_j U$ ist, gilt
    $|g_j U| = |U|$ für alle $j$. Somit ist $|G| = |J| \cdot |U| = [G:U] \cdot |U|$.
\end{proof}


\begin{beobachtung}[Bedeutung von Lagrange]
    Der Satz von Lagrange ist ein erstes tiefes Ergebnis: Er sagt, dass die
    Größe einer Untergruppe die Größe der Gesamtgruppe teilt -- eine starke
    Einschränkung, die allein aus der Gruppenstruktur folgt. Insbesondere hat
    jedes Element $g$ einer endlichen Gruppe $G$ eine Ordnung $\mathrm{Ord}\, g$,
    die $|G|$ teilt, denn $\langle g \rangle$ ist eine Untergruppe.
\end{beobachtung}

\begin{definition}[Zyklische Gruppe, Ordnung]
    Sei $G$ eine Gruppe und $g \in G$. Die von $g$ erzeugte Untergruppe ist
    $\langle g \rangle = \{g^k \mid k \in \mathbb{Z}\}$.
    Ist $|\langle g \rangle| < \infty$, so heißt diese Zahl die
    \textbf{Ordnung} von $g$, geschrieben $\mathrm{Ord}\, g$.
    Ist $G = \langle g \rangle$ für ein $g \in G$, so heißt $G$ \textbf{zyklisch}.
\end{definition}

\begin{definition}[Eulersche Funktion]
    Für $n \in \mathbb{N}$ ist die \textbf{Eulersche Funktion}
    \[
        \varphi(n) = |\{k \mid k \in \mathbb{Z},\; 1 \le k \le n,\; \gcd(k, n) = 1\}|.
    \]
    Es gilt $\varphi(p) = p - 1$ für Primzahlen $p$.
\end{definition}

% ============================================================
\section{Ringe und Körper}
% ============================================================

\begin{definition}[Ring und Körper]
    Eine Menge $R$ mit zwei Verknüpfungen $+$ und $\cdot$ heißt \textbf{Ring},
    falls:
    \begin{itemize}
        \item[-] $(R, +)$ ist eine abelsche Gruppe mit neutralem Element $0$.
        \item[-] Die Multiplikation ist assoziativ mit Einselement $1 \in R$.
        \item[-] Es gelten die \textbf{Distributivgesetze}:
              $(a + b)c = ac + bc$ und $a(b + c) = ab + ac$.
    \end{itemize}
    Ein Ring heißt \textbf{kommutativ}, wenn $ab = ba$ für alle $a, b \in R$.

    Ein kommutativer Ring $R$, in dem jedes $a \neq 0$ ein multiplikatives
    Inverses $a^{-1}$ besitzt, heißt \textbf{Körper}.
\end{definition}

\begin{beobachtung}[Die Körperhierarchie]
    Die wichtigsten Körper sind:
    \[
        \mathbb{Q} \subset \mathbb{R} \subset \mathbb{C}.
    \]
    Daneben gibt es \textbf{endliche Körper}: Für jede Primzahl $p$ ist
    $\mathbb{Z}_p = \mathbb{Z}/p\mathbb{Z}$ ein Körper mit $p$ Elementen. Allgemein
    hat jeder endliche Körper $q = p^n$ Elemente für eine Primzahl $p$.

    Der Körper $\mathbb{F}_2 = \{0, 1\}$ ist der kleinstmögliche Körper -- er
    ist die algebraische Grundlage der Informatik.
\end{beobachtung}

% ============================================================
\section{Der komplexe Zahlkörper}
% ============================================================

\begin{definition}[Komplexe Zahlen]
    Auf $\mathbb{R} \times \mathbb{R}$ definieren wir:
    \[
        (a_1, b_1) + (a_2, b_2) = (a_1 + a_2,\; b_1 + b_2), \qquad
        (a_1, b_1)(a_2, b_2) = (a_1 a_2 - b_1 b_2,\; a_1 b_2 + a_2 b_1).
    \]
    Die entstehende Struktur heißt \textbf{komplexer Zahlkörper} $\mathbb{C}$.
    Mit $i = (0, 1)$ gilt $i^2 = -1$ und jedes $c \in \mathbb{C}$ hat die
    Gestalt $c = a + bi$ mit $a, b \in \mathbb{R}$.

    Für $c = a + bi$ definieren wir:
    \begin{itemize}
        \item[-] \textbf{Realteil}: $\mathrm{Re}\, c = a$, \quad \textbf{Imaginärteil}: $\mathrm{Im}\, c = b$
        \item[-] \textbf{Betrag}: $|c| = \sqrt{a^2 + b^2}$
        \item[-] \textbf{Konjugierte}: $\bar{c} = a - bi$, \quad es gilt $c \bar{c} = |c|^2$
    \end{itemize}
\end{definition}

\begin{satz}[Geometrische Deutung der Multiplikation]
    Schreibt man $c_j = r_j(\cos \alpha_j + i \sin \alpha_j)$ mit $r_j = |c_j|$,
    so gilt:
    \[
        c_1 c_2 = r_1 r_2 \bigl(\cos(\alpha_1 + \alpha_2) + i \sin(\alpha_1 + \alpha_2)\bigr).
    \]
    Bei der Multiplikation werden also \emph{Längen multipliziert} und
    \emph{Winkel addiert}.
\end{satz}
\begin{proof}
    Seien $z_1 = r_1 e^{i\varphi_1}$ und $z_2 = r_2 e^{i\varphi_2}$.
    Nach der Eulerschen Formel gilt:
    \[
        z_1 z_2 = r_1 r_2 e^{i(\varphi_1 + \varphi_2)}.
    \]
    Dies folgt direkt aus $e^{i\alpha} e^{i\beta} = e^{i(\alpha+\beta)}$,
    was man durch Ausmultiplizieren von $(\cos\alpha + i\sin\alpha)(\cos\beta + i\sin\beta)$
    und Anwenden der Additionstheoreme beweist.
\end{proof}


\begin{beobachtung}[Warum \texorpdfstring{$\mathbb{C}$}{C} vollständig ist]
    Der \textbf{Fundamentalsatz der Algebra} besagt: Jedes Polynom vom
    Grad $n \ge 1$ über $\mathbb{C}$ hat genau $n$ Nullstellen (mit
    Vielfachheit gezählt). Dies unterscheidet $\mathbb{C}$ fundamental von
    $\mathbb{R}$: Das Polynom $x^2 + 1$ hat über $\mathbb{R}$ keine, über
    $\mathbb{C}$ die Nullstellen $\pm i$.

    $\mathbb{C}$ ist \textbf{algebraisch abgeschlossen} -- man kann keinen
    noch größeren Körper durch adjungieren von Nullstellen von Polynomen
    erhalten. Es ist in diesem Sinne der \emph{vollständige} Körper.
\end{beobachtung}

% ============================================================
\section{Das RSA-Verfahren}
% ============================================================

\begin{beobachtung}[Gruppen in der Kryptographie]
    Das RSA-Verfahren ist ein unmittelbarer Ausfluss der Gruppentheorie.
    Bob wählt zwei große Primzahlen $p \neq q$, setzt $n = pq$ und wählt
    $e$ teilerfremd zu $\varphi(n) = (p-1)(q-1)$. Er bestimmt $d$ mit
    $ed \equiv 1 \pmod{\varphi(n)}$.

    \textbf{Ver- und Entschlüsselung}:
    \begin{itemize}
        \item[-] $\mathbf{e} : \mathbb{Z}_n \to \mathbb{Z}_n, \quad [x] \mapsto [x]^e$ \quad (Verschlüsseln)
        \item[-] $\mathbf{d} : \mathbb{Z}_n \to \mathbb{Z}_n, \quad [y] \mapsto [y]^d$ \quad (Entschlüsseln)
    \end{itemize}
    Es gilt $\mathbf{d} \circ \mathbf{e} = \mathrm{id}$, weil $x^{ed} \equiv x \pmod{n}$
    -- eine direkte Folge des Satzes von Euler: $x^{\varphi(n)} \equiv 1 \pmod{n}$
    für $\gcd(x, n) = 1$.

    Die \emph{Stärke} des Verfahrens liegt darin, dass man ohne die
    Primfaktorisierung von $n$ das $d$ nicht aus $e$ und $n$ berechnen kann --
    und große Zahlen mit heutigem Aufwand nicht faktorisierbar sind.
\end{beobachtung}

% ============================================================
\section{Vektorräume und Unterräume}
% ============================================================

\begin{definition}[Vektorraum]
    Sei $K$ ein Körper. Eine Menge $V$ heißt \textbf{$K$-Vektorraum}, falls:
    \begin{itemize}
        \item[-] $(V, +)$ ist eine abelsche Gruppe.
        \item[-] Es gibt eine \textbf{skalare Multiplikation}
              $K \times V \to V$, $(k, v) \mapsto kv$, die für alle $k, l \in K$
              und $u, v \in V$ erfüllt:
        \begin{itemize}
            \item[-] $k(u + v) = ku + kv$, \quad $(k + l)v = kv + lv$
            \item[-] $(kl)v = k(lv)$, \quad $1v = v$
        \end{itemize}
    \end{itemize}
    Die Elemente von $V$ heißen \textbf{Vektoren}, die Elemente von $K$
    \textbf{Skalare}.
\end{definition}

\begin{beispiel}[Grundlegende Vektorräume]
    \begin{itemize}
        \item[-] $K^n$ mit komponentenweiser Addition und Skalarmultiplikation
              ist der Standard-Vektorraum.
        \item[-] Der Raum $\mathrm{Ab}(M, K)$ aller Abbildungen von einer
              Menge $M$ nach $K$ ist ein Vektorraum.
        \item[-] Der Raum $C(\mathbb{R})$ aller stetigen reellwertigen Funktionen
              ist ein $\mathbb{R}$-Vektorraum.
        \item[-] Jeder Körper $L$, der $K$ enthält, ist ein $K$-Vektorraum.
    \end{itemize}
\end{beispiel}

\begin{definition}[Unterraum]
    Eine nichtleere Teilmenge $U \subseteq V$ heißt \textbf{Unterraum}, falls
    $U$ bezüglich Addition und Skalarmultiplikation abgeschlossen ist, d.\,h.:
    \begin{itemize}
        \item[-] $u_1, u_2 \in U \Rightarrow u_1 + u_2 \in U$
        \item[-] $k \in K,\; u \in U \Rightarrow ku \in U$
    \end{itemize}
    Der \textbf{Durchschnitt} beliebig vieler Unterräume ist wieder ein
    Unterraum. Die \textbf{Summe} $U_1 + U_2 = \{u_1 + u_2 \mid u_j \in U_j\}$
    ist der kleinste Unterraum, der $U_1$ und $U_2$ enthält.
\end{definition}

\begin{beobachtung}[Vereinigung von Unterräumen]
    Die Vereinigung $U_1 \cup U_2$ zweier Unterräume ist \emph{im Allgemeinen
    kein Unterraum}. Sie ist genau dann einer, wenn $U_1 \subseteq U_2$ oder
    $U_2 \subseteq U_1$ gilt.

    Dies ist kein formaler Defekt, sondern spiegelt etwas Reales wider:
    Unterräume sind lineare Strukturen -- ihre Vereinigung enthält keine
    Linearkombinationen aus beiden Teilen.
\end{beobachtung}

% ============================================================
\section{Lineare Abhängigkeit, Träger, Dimension}
% ============================================================

\begin{definition}[Linearkombination und Projektion]
    Sei $M \subseteq V$. Die Menge aller endlichen Linearkombinationen
    \[
        \langle M \rangle = \left\{ \sum_{j=1}^{n} k_j m_j \;\Big|\;
        k_j \in K,\; m_j \in M,\; n \in \mathbb{N} \right\}
    \]
    heißt die \textbf{Projektion} von $M$ -- der kleinste Unterraum
    von $V$, der $M$ enthält. Jeder Vektor in $\langle M \rangle$ ist
    ein Bild von $M$ unter linearen Koeffizienten; $M$ wird also auf
    diesen Unterraum \emph{projiziert}.

    $M$ heißt \textbf{Projektionssystem} von $V$, falls $\langle M \rangle = V$,
    d.\,h. jeder Vektor in $V$ durch Projektion aus $M$ erreichbar ist.
\end{definition}

\begin{definition}[Lineare Unabhängigkeit]
    Ein System $S = [v_i \mid i \in I]$ in $V$ heißt \textbf{linear unabhängig},
    falls kein $v_k$ in der Projektion der übrigen liegt:
    \[
        v_k \notin \langle v_i \mid i \neq k \rangle \quad \text{für alle } k \in I.
    \]
    Das bedeutet: Kein Vektor des Systems ist durch Projektion aus den
    anderen erreichbar -- jeder trägt eine echte neue Richtung bei, die
    durch die übrigen nicht abgedeckt wird.

    Äquivalent: $\sum_{j \in J} k_j v_j = 0$ mit endlichem $J \subseteq I$
    impliziert $k_j = 0$ für alle $j \in J$.

    Wird ein Vektor $v_k$ aus einem linear unabhängigen System entfernt,
    so ist seine Richtung durch keine Projektion der verbleibenden Vektoren
    mehr erreichbar -- die Projektion $\langle v_i \mid i \neq k \rangle$
    kollabiert um genau eine Dimension.
\end{definition}

\begin{definition}[Träger und Dimension]
    Ein System $B = [v_i \mid i \in I]$ in $V$ heißt \textbf{Träger} von $V$,
    falls $B$ linear unabhängig ist und $\langle B \rangle = V$ gilt --
    also das vollständige Projektionssystem von $V$ bildet.

    Äquivalent: Jedes $v \in V$ besitzt eine \emph{eindeutige} Darstellung
    $v = \sum_{j \in J} k_j v_j$ mit $k_j \in K$ und endlichem $J \subseteq I$.

    Ist $V$ durch endlich viele Vektoren projektierbar, so hat $V$ eine
    endlichen Träger, und jede Träger von $V$ enthält gleich viele Vektoren.
    Diese Anzahl heißt die \textbf{Dimension} $\dim_K V$.
\end{definition}

\begin{satz}[Existenz von Trägern]
    Sei $V = \langle w_1, \dots, w_m \rangle$ durch endlich viele Vektoren
    projektierbar.
    \begin{itemize}
        \item[-] Jedes linear unabhängige System $[w_1, \dots, w_k]$ lässt sich
              zu einem Träger von $V$ ergänzen.
        \item[-] Jedes Projektionssystem enthält einen Träger als Teilsystem.
        \item[-] Je zwei Träger von $V$ haben gleich viele Elemente.
    \end{itemize}
\end{satz}
\begin{proof}
    Sei $E$ ein endliches Erzeugendensystem von $V$. Wir wählen aus $E$
    ein inklusionsmaximal linear unabhängiges Teilsystem $B$. Dann ist
    $B$ linear unabhängig per Konstruktion. Jedes $v \in E \setminus B$
    liegt in $\langle B \rangle$, denn sonst wäre $B \cup \{v\}$ noch linear
    unabhängig -- Widerspruch zur Maximalität. Da $E$ ein Erzeugendensystem ist,
    folgt $\langle B \rangle = V$, also ist $B$ ein Träger.
\end{proof}


\begin{beobachtung}[Dimension ist die Freiheit des Raumes]
    Die Dimension zählt, wieviele unabhängige Richtungen in einem Raum
    möglich sind. $\dim K^n = n$, $\dim \{0\} = 0$.

    Für Unterräume gilt stets $\dim U \le \dim V$, mit Gleichheit genau
    dann, wenn $U = V$. Die \textbf{Dimensionsformel} für zwei Unterräume
    lautet:
    \[
        \dim(U_1 + U_2) = \dim U_1 + \dim U_2 - \dim(U_1 \cap U_2).
    \]
    Dies ist das vektorräumliche Analogon des Inklusions-Exklusions-Prinzips
    für Mengen.
\end{beobachtung}

\begin{satz}[Steinitzscher Austauschsatz]
    Sei $[v_1, \dots, v_n]$ einen Träger von $V$ und $[w_1, \dots, w_m]$ ein
    linear unabhängiges System in $V$. Dann gilt $m \le n$, und bei geeigneter
    Numerierung ist $[w_1, \dots, w_m, v_{m+1}, \dots, v_n]$ einen Träger von $V$.
\end{satz}
\begin{proof}
    Durch Induktion nach $m$. \textbf{$m = 1$:} Da $w_1 \neq 0$ und $[v_1, \ldots, v_n]$
    einen Träger ist, gilt $w_1 = \sum_i k_i v_i$ mit einem $k_j \neq 0$.
    Nach dem Trägeraustauschlemma ist dann $[w_1, v_2, \ldots, v_n]$ (nach Umnummerierung)
    einen Träger.

    \textbf{$m > 1$:} Nach Induktionsannahme gibt es einen Träger
    $[w_1, \ldots, w_{m-1}, v_m, \ldots, v_n]$ (umnumeriert). Schreibe
    $w_m = \sum_{j<m} k_j w_j + \sum_{i \geq m} d_i v_i$.
    Da $[w_1, \ldots, w_m]$ linear unabhängig ist, muss ein $d_i \neq 0$ sein
    (sonst wäre $w_m$ Linearkombination der $w_j$). Nach Umnumerierung sei
    $d_m \neq 0$; dann kann $v_m$ durch $w_m$ ersetzt werden, und
    $[w_1, \ldots, w_m, v_{m+1}, \ldots, v_n]$ ist einen Träger.
    Wäre $m > n$, so wäre $[w_1, \ldots, w_n]$ einen Träger, also $w_{n+1}
    \in \langle w_1, \ldots, w_n \rangle$ -- Widerspruch zur linearen Unabhängigkeit.
\end{proof}


\begin{beispiel}[Transzendente Zahlen]
    Der Körper $\mathbb{R}$ ist ein $\mathbb{Q}$-Vektorraum von unendlicher
    Dimension. Ein $a \in \mathbb{C}$ heißt \textbf{algebraisch} über $\mathbb{Q}$,
    wenn $[1, a, a^2, \dots]$ linear abhängig ist -- also $a$ Nullstelle eines
    Polynoms mit rationalen Koeffizienten ist. Andernfalls heißt $a$
    \textbf{transzendent}. Die Zahlen $e$ und $\pi$ sind transzendent;
    dies ist eines der tiefsten Ergebnisse der klassischen Analysis.
\end{beispiel}

% ============================================================
\chapter{Lineare Abbildungen und Matrizen}
% ============================================================

\section{Lineare Abbildungen}

\begin{definition}[Lineare Abbildung]
    Seien $V$ und $W$ zwei $K$-Vektorräume. Eine Abbildung
    $f : V \to W$ heißt \textbf{linear} (oder \textbf{$K$-linear}), falls
    für alle $u, v \in V$ und $k \in K$ gilt:
    \[
        f(u + v) = fu + fv, \qquad f(kv) = k \cdot fv.
    \]
    Äquivalent: $f(k_1 v_1 + k_2 v_2) = k_1 f v_1 + k_2 f v_2$ für alle
    $k_i \in K$, $v_i \in V$.
\end{definition}

\begin{beobachtung}[Kern und Bild]
    Für eine lineare Abbildung $f : V \to W$ sind
    \[
        \ker f = \{v \in V \mid fv = 0\} \subseteq V, \qquad
        \mathrm{Im}\, f = fV \subseteq W
    \]
    jeweils Unterräume. Die fundamentale Beziehung lautet:
    \[
        \dim V = \dim \ker f + \dim \mathrm{Im}\, f.
    \]
    Diese \textbf{Dimensionsformel} (Rangsatz) sagt: Was in den Kern geht,
    geht verloren; was übrig bleibt, erzeugt das Bild. Die beiden Hälften
    addieren sich zur Gesamtdimension.

    Insbesondere: $f$ ist injektiv genau dann, wenn $\ker f = \{0\}$, und
    surjektiv genau dann, wenn $\dim \mathrm{Im}\, f = \dim W$.
\end{beobachtung}

\begin{definition}[Isomorphismus]
    Eine bijektive lineare Abbildung $f : V \to W$ heißt
    \textbf{Isomorphismus}. Existiert ein solcher, heißen $V$ und $W$
    \textbf{isomorph}: $V \cong W$.

    Zwei endlichdimensionale $K$-Vektorräume sind genau dann isomorph,
    wenn sie dieselbe Dimension haben. Insbesondere gilt $V \cong K^n$
    für jeden $n$-dimensionalen $K$-Vektorraum.
\end{definition}

% ============================================================
\section{Matrizen}
% ============================================================

\begin{definition}[Matrix und Darstellungsmatrix]
    Eine $(m \times n)$-\textbf{Matrix} $A$ über $K$ ist ein rechteckiges
    Schema von Einträgen $a_{ij} \in K$:
    \[
        A = (a_{ij})_{\substack{1 \le i \le m \\ 1 \le j \le n}}.
    \]
    Ist $f : V \to W$ linear mit gegebenen Träger $[v_1, \dots, v_n]$ von $V$
    und $[w_1, \dots, w_m]$ von $W$, so ist die \textbf{Darstellungsmatrix}
    $A_f = (a_{ij})$ definiert durch:
    \[
        f v_j = \sum_{i=1}^m a_{ij} w_i \qquad (j = 1, \dots, n).
    \]
    Die $j$-te Spalte von $A_f$ enthält also die Koordinaten von $fv_j$
    bezüglich des Trägers von $W$.
\end{definition}

\begin{satz}[Matrizenrechnung]
    Für lineare Abbildungen $f : V \to W$ und $g : W \to X$ mit gegebenen
    Träger gilt:
    \[
        A_{gf} = A_g \cdot A_f, \qquad (A_g \cdot A_f)_{ij}
        = \sum_{k=1}^{m} (A_g)_{ik} (A_f)_{kj}.
    \]
    Die Matrizenmultiplikation entspricht der Komposition linearer Abbildungen.
\end{satz}
\begin{proof}
    Die Assoziativität $(AB)C = A(BC)$ folgt durch direktes Ausrechnen der
    $(i,k)$-Einträge: $\sum_j (\sum_l a_{il} b_{lj}) c_{jk} = \sum_l a_{il}
    (\sum_j b_{lj} c_{jk})$ nach Umordnung endlicher Summen.
    Das Distributivgesetz $A(B+C) = AB + AC$ ist komponentenweise trivial.
    Die Formel $(AB)^t = B^t A^t$ ergibt sich aus $((AB)^t)_{ij} = (AB)_{ji}
    = \sum_k a_{jk} b_{ki} = \sum_k (B^t)_{ik} (A^t)_{kj} = (B^t A^t)_{ij}$.
\end{proof}


\begin{beobachtung}[Warum Zeilen- und Spaltenindex ihre Rollen tauschen]
    Der $j$-te Trägerelement $v_j$ wird durch die $j$-te \emph{Spalte} von $A_f$
    beschrieben. Das liegt daran, dass $f$ von $V$ nach $W$ abbildet:
    Eingaben indexieren Spalten, Ausgaben Zeilen.

    Bei der Matrizenmultiplikation $C = AB$ ist $(C)_{ij} = \sum_k (A)_{ik}(B)_{kj}$
    -- Zeile $i$ von $A$ wird mit Spalte $j$ von $B$ skalar multipliziert.
    Dies entspricht exakt dem Hintereinanderausführen: erst $B$ (Spaltenstruktur
    als Eingabe), dann $A$ (Zeilenstruktur als Ausgabe).
\end{beobachtung}

% ============================================================
\section{Strukturen mit Schwerpunkt: Vektor und Matrix im Vergleich}
% ============================================================

Der Vektor und die Matrix sind die zwei zentralen Datenstrukturen der
linearen Algebra. Sie unterscheiden sich fundamental in ihrer
\textbf{algebraischen Symmetrie} -- und diese Symmetrie entscheidet
unmittelbar darüber, wie effizient mit ihnen gerechnet werden kann.

\begin{beobachtung}[Der Vektor hat einen Schwerpunkt]
    Ein Vektor $v = (v_1, \dots, v_n)^T \in K^n$ ist eine \textbf{geordnete
    Folge von Skalaren mit einer einzigen Indexstruktur}. Diese Einheitlichkeit
    ist sein algebraischer \textbf{Schwerpunkt}: Alle Elemente liegen auf
    derselben strukturellen Ebene.

    Die Folge daraus ist unmittelbar: Addition und Skalarmultiplikation
    wirken \textbf{komponentenweise} -- paarweise, elementweise, ohne
    Wechselwirkung zwischen verschiedenen Einträgen:
    \[
        u + v = (u_1 + v_1,\; u_2 + v_2,\; \dots,\; u_n + v_n)^T, \qquad
        kv = (kv_1,\; kv_2,\; \dots,\; kv_n)^T.
    \]
    Jede Operation ist lokal: Das Ergebnis an Position $i$ hängt nur von
    $u_i$ und $v_i$ ab, nicht von anderen Einträgen. Dies ist die
    direkteste und effizienteste Form algebraischer Berechnung --
    $n$ unabhängige Operationen, parallelisierbar, ohne Übertragsstruktur.

    Der Vektor ist in seinem Schwerpunkt am effizientesten zu lösen.
\end{beobachtung}

\begin{beobachtung}[Die Matrix hat keinen Schwerpunkt]
    Eine Matrix $A = (a_{ij}) \in K^{m \times n}$ besitzt eine
    \textbf{zweidimensionale Indexstruktur}: Jeder Eintrag trägt zwei
    Koordinaten $(i, j)$ -- Zeile und Spalte. Diese Zweidimensionalität
    erzeugt ein strukturelles Ungleichgewicht: Es gibt keine einheitliche
    ``Richtung'', in der alle Operationen gleichartig verlaufen.

    Dies zeigt sich unmittelbar bei der Matrizenmultiplikation: Das
    Ergebnis $(AB)_{ij} = \sum_k a_{ik} b_{kj}$ hängt von einer ganzen
    \emph{Zeile} von $A$ und einer ganzen \emph{Spalte} von $B$ ab --
    nicht nur von einem Paar von Einzeleinträgen. Addition von Matrizen
    ist noch komponentenweise möglich, aber Multiplikation ist es nicht:
    \[
        (A + B)_{ij} = a_{ij} + b_{ij}, \qquad
        (AB)_{ij} = \sum_{k=1}^{n} a_{ik}\, b_{kj}.
    \]
    Die Matrizenmultiplikation ist eine globale Operation: Jeder
    Ausgabeeintrag hängt von $n$ Eingangseinträgen ab. Die Komplexität
    des naiven Algorithmus ist $O(n^3)$ -- kubisch statt linear.

    Die Matrix hat keinen Schwerpunkt und ist damit nicht ohne Weiteres
    elementweise effizient zu lösen.
\end{beobachtung}

\begin{beobachtung}[Das Ziel: Strukturen mit Schwerpunkt finden]
    Das fundamentale Ziel der linearen Algebra und der algorithmischen
    Matrizentheorie lautet: \textbf{Algebraische Strukturen finden, die
    einen Schwerpunkt besitzen} -- und damit im Schwerpunkt am
    effizientesten lösbar sind.

    Dieses Prinzip manifestiert sich auf mehreren Ebenen:
    \begin{itemize}
        \item[-] \textbf{Diagonalmatrizen}: Eine Diagonalmatrix
              $D = \mathrm{diag}(d_1, \dots, d_n)$ hat ihren Schwerpunkt
              in der Hauptdiagonale. Multiplikation und Invertierung
              wirken wieder komponentenweise: $(DA)_{ij} = d_i a_{ij}$,
              $D^{-1} = \mathrm{diag}(d_1^{-1}, \dots, d_n^{-1})$.
              Die zweidimensionale Matrix ist auf eine eindimensionale
              Struktur reduziert.
        \item[-] \textbf{Dia\-go\-nali\-sie\-rung}: Eine diagonalisierbare Matrix
              $A = PDP^{-1}$ wird durch Träger\-wechsel in eine Diagonalmatrix
              überführt. Im Schwerpunkt der Eigenvektor-Träger verhält sich $A$
              wie ein Vektor von Skalaren -- die nicht-schwerpunkt\-artige Struktur
              ist durch die richtige Wahl des Trägers eliminiert.
        \item[-] \textbf{Dünn besetzte Matrizen (Sparse Matrices)}: Hat
              eine Matrix viele Nulleinträge, so konzentriert sich ihre
              wesentliche Struktur auf wenige Einträge -- einen impliziten
              Schwerpunkt. Algorithmen, die diese Sparsity ausnutzen,
              erreichen subkubische Laufzeit.
        \item[-] \textbf{Orthonormalbasen}: In einer Orthonormalträger
              entkoppeln die Koordinaten vollständig. Das Skalarprodukt
              wird zur komponentenweisen Summe $\langle u, v \rangle =
              \sum_j u_j v_j$ -- Vektorrechnung mit Schwerpunkt.
    \end{itemize}

    Das Streben nach Schwerpunkt ist kein rein algorithmisches Anliegen,
    sondern ein strukturelles: Eine Algebra, in der Operationen lokal
    und komponentenweise wirken, ist nicht nur schneller berechenbar,
    sondern auch konzeptuell transparenter. Die Diagonalform ist die
    reinste Verwirklichung dieses Prinzips -- die Matrix, die sich
    wie ein Vektor verhält.
\end{beobachtung}

% ============================================================
\section{Determinanten}
% ============================================================

\begin{definition}[Determinante]
    Die \textbf{Determinante} $\det A \in K$ einer $(n \times n)$-Matrix
    ist die eindeutige Abbildung $\det : K^{n \times n} \to K$, die
    \begin{itemize}
        \item[-] in jeder Zeile \textbf{linear} ist,
        \item[-] beim Vertauschen zweier Zeilen das \textbf{Vorzeichen wechselt},
        \item[-] für die Einheitsmatrix $\det \mathbf{1} = 1$ gilt.
    \end{itemize}
    Explizit: $\det A = \sum_{\sigma \in S_n} \mathrm{sgn}(\sigma) \prod_{i=1}^n a_{i,\sigma(i)}$,
    wobei die Summe über alle Permutationen $\sigma$ der $\{1, \dots, n\}$ läuft.
\end{definition}

\begin{satz}[Eigenschaften der Determinante]
    Seien $A, B \in K^{n \times n}$. Dann gilt:
    \begin{itemize}
        \item[-] $\det(AB) = \det A \cdot \det B$ \quad (\textbf{Multiplikativität})
        \item[-] $\det A^{-1} = (\det A)^{-1}$, falls $A$ invertierbar
        \item[-] $\det A^T = \det A$
        \item[-] $A$ ist genau dann invertierbar, wenn $\det A \neq 0$.
    \end{itemize}
\end{satz}
\begin{proof}
    \textbf{Multilinearität und Alternierendheit} folgen direkt aus der Leibniz-Formel
    $\det A = \sum_{\pi \in S_n} \mathrm{sgn}(\pi) \prod_i a_{i,\pi(i)}$:
    Skalierung einer Zeile um $c$ multipliziert jeden Summanden mit $c$.
    Vertauschen zweier Zeilen entspricht einer Transposition aller Permutationen,
    was das Signum wechselt.
    \textbf{Entwicklung nach Spalte $j$:} $\det A = \sum_i a_{ij} (-1)^{i+j} M_{ij}$
    folgt durch Gruppieren der Leibniz-Summanden nach dem Wert $\pi(j) = i$
    und Erkennen, dass $(-1)^{i+j} M_{ij}$ das Signum des zugehörigen Minors kodiert.
    \textbf{Ähnlichkeit:} $\det(P^{-1}AP) = \det(P^{-1})\det(A)\det(P) = \det A$.
\end{proof}


\begin{beobachtung}[Geometrische Bedeutung]
    Die Determinante misst das \textbf{orientierte Volumen} des von den
    Spaltenvektoren aufgespannten Parallelotops. Für $n = 2$ ist $|\det A|$
    der Flächeninhalt des Parallelogramms mit den Seiten $a_1, a_2$;
    für $n = 3$ das Volumen des Spats.

    Das Vorzeichen gibt die Orientierung an: $\det A > 0$ bedeutet, dass
    die Spalten dieselbe Orientierung wie die Standardträger haben;
    $\det A < 0$ bedeutet eine Spiegelung.
\end{beobachtung}

% ============================================================
\section{Eigenwerte und Diagonalisierung}
% ============================================================

\begin{definition}[Eigenwert und Eigenvektor]
    Sei $A \in K^{n \times n}$. Ein Skalar $\lambda \in K$ heißt
    \textbf{Eigenwert} von $A$, falls es ein $v \neq 0$ gibt mit
    \[
        Av = \lambda v.
    \]
    Ein solches $v$ heißt \textbf{Eigenvektor} zum Eigenwert $\lambda$.
    Die Eigenwerte sind genau die Nullstellen des
    \textbf{charakteristischen Polynoms}:
    \[
        \chi_A(\lambda) = \det(\lambda \mathbf{1} - A).
    \]
\end{definition}

\begin{beobachtung}[Was Eigenwerte beschreiben]
    Ein Eigenvektor ist eine \emph{Richtung, die invariant bleibt}: Die
    lineare Abbildung $A$ streckt oder staucht in dieser Richtung um den
    Faktor $\lambda$, dreht sie aber nicht. Das macht Eigenvektoren so
    nützlich -- sie zeigen die \emph{natürlichen Koordinaten} eines Systems.

    In der Dynamik $x_{t+1} = Ax_t$ bestimmt der betragsgrößte Eigenwert
    das Langzeitverhalten: Ist $|\lambda_{\max}| < 1$, so konvergiert das
    System gegen $0$; ist $|\lambda_{\max}| > 1$, so divergiert es.
    Dieses Prinzip ist identisch mit dem der e-Funktion in der Signaltheorie:
    Der Realteil des Exponenten entscheidet über Wachstum oder Zerfall.
\end{beobachtung}

\begin{definition}[Diagonalisierbarkeit]
    $A$ heißt \textbf{diagonalisierbar} über $K$, falls einen Träger
    $[v_1, \dots, v_n]$ von $K^n$ aus Eigenvektoren von $A$ existiert.
    Äquivalent: Es gibt eine invertierbare Matrix $P$ mit
    \[
        P^{-1} A P = \mathrm{diag}(\lambda_1, \dots, \lambda_n).
    \]
    Hat $A$ genau $n$ verschiedene Eigenwerte, so ist $A$ stets
    diagonalisierbar.
\end{definition}

\begin{satz}[Jordansche Normalform]
    Über einem algebraisch abgeschlossenen Körper (z.\,B. $\mathbb{C}$) ist
    jede Matrix $A$ ähnlich zu einer \textbf{Jordanmatrix}:
    \[
        J = \mathrm{diag}(J_{k_1}(\lambda_1), \dots, J_{k_r}(\lambda_r)), \qquad
        J_k(\lambda) = \begin{pmatrix} \lambda & 1 & & \\ & \ddots & \ddots & \\
        & & \lambda & 1 \\ & & & \lambda \end{pmatrix}.
    \]
    Die Jordanform ist eindeutig (bis auf Reihenfolge der Blöcke) und die
    feinste Normalform, die mit Ähnlichkeitstransformationen erreichbar ist.
\end{satz}
\begin{proof}[Beweisskizze]
    \textbf{Existenz (über $\mathbb{C}$):} Da das charakteristische Polynom
    $\chi_A$ über $\mathbb{C}$ in Linearfaktoren zerfällt, gilt
    $\chi_A = \prod_j (\lambda - \lambda_j)^{m_j}$ mit algebraischen
    Vielfachheiten $m_j$. Nach der Primärzerlegung (Chinesischer Restsatz)
    ist $V = \bigoplus_j \ker(A - \lambda_j E)^{m_j}$.
    Auf jedem Primärraum $\ker(A - \lambda_j E)^{m_j}$ ist $N_j = A - \lambda_j E$
    nilpotent; für nilpotente Abbildungen zeigt man induktiv, dass
    einen Träger existiert, bezüglich derer $N_j$ Jordanform hat.
    \textbf{Eindeutigkeit:} Folgt aus der Eindeutigkeitsaussage über Ränge
    der Potenzen $(A - \lambda E)^k$ (vgl.\ Satz über Eindeutigkeit der Jordanform).
\end{proof}


% ============================================================
\section{Stochastische Matrizen}
% ============================================================

\begin{definition}[Stochastische Matrix]
    Eine Matrix $A \in \mathbb{R}^{n \times n}$ mit $a_{ij} \ge 0$ heißt
    \textbf{(rechts)stochastisch}, falls alle Zeilensummen gleich $1$ sind:
    $\sum_{j=1}^n a_{ij} = 1$ für alle $i$.

    Eine stochastische Matrix beschreibt einen \textbf{Markovprozess}:
    $a_{ij}$ ist die Wahrscheinlichkeit, vom Zustand $i$ in Zustand $j$
    zu wechseln.
\end{definition}

\begin{beobachtung}[Eigenwert 1 und stationäre Verteilung]
    Jede stochastische Matrix hat den Eigenwert $1$: Der Zeilenvektor
    $\mathbf{1}^T = (1, \dots, 1)^T$ ist rechts ein Eigenvektor. Alle
    anderen Eigenwerte haben Betrag $\le 1$ (Satz von Perron-Frobenius).

    Das Langzeitverhalten $\lim_{t \to \infty} A^t$ konvergiert für
    irreduzible aperiodische Matrizen gegen eine Matrix, deren Zeilen
    gleich der \textbf{stationären Verteilung} $\pi$ sind -- dem
    einzigen Eigenvektor zum Eigenwert $1$.

    Dies ist die Matrizenversion des Satzes: Jedes dynamische System
    mit dominantem Eigenwert $\lambda_1$ wird für große $t$ durch den
    zugehörigen Eigenvektor beschrieben.
\end{beobachtung}

% ============================================================
\section{Vektorräume mit Skalarprodukt}
% ============================================================

\begin{definition}[Skalarprodukt]
    Ein \textbf{Skalarprodukt} auf einem $\mathbb{R}$-Vektorraum $V$ ist eine
    Abbildung $\langle \cdot, \cdot \rangle : V \times V \to \mathbb{R}$ mit:
    \begin{itemize}
        \item[-] \textbf{Bilinearität}: Linear in beiden Argumenten.
        \item[-] \textbf{Symmetrie}: $\langle u, v \rangle = \langle v, u \rangle$.
        \item[-] \textbf{Positive Definitheit}: $\langle v, v \rangle > 0$
              für alle $v \neq 0$.
    \end{itemize}
    Das Standardskalarprodukt auf $\mathbb{R}^n$ ist
    $\langle x, y \rangle = \sum_{i=1}^n x_i y_i = x^T y$.
\end{definition}

\begin{definition}[Norm und Orthogonalität]
    Die durch das Skalarprodukt induzierte \textbf{Norm} ist
    $\|v\| = \sqrt{\langle v, v \rangle}$. Zwei Vektoren heißen
    \textbf{orthogonal}, falls $\langle u, v \rangle = 0$.
    Ein System orthogonaler Einheitsvektoren heißt \textbf{Orthonormalträger} (ONB).
\end{definition}

\begin{satz}[Gram-Schmidt und Orthogonalprojektion]
    Jeder endlichdimensionale Unterraum mit Skalarprodukt besitzt eine
    Orthonormalträger (Gram-Schmidt-Verfahren). Für einen Unterraum $U \le V$
    und $v \in V$ ist die \textbf{Orthogonalprojektion} auf $U$ der eindeutige
    Vektor $p \in U$ mit $v - p \perp U$.
\end{satz}
\begin{proof}
    \textbf{Schmidt-Verfahren:} Wir setzen $\tilde{w}_1 = v_1$ und induktiv
    $\tilde{w}_j = v_j - \sum_{k<j} (v_j, w_k) w_k$, dann $w_j = \tilde{w}_j / \|\tilde{w}_j\|$.
    Per Konstruktion gilt $(w_j, w_k) = 0$ für $k < j$ (jeder Summand hebt genau
    einen Anteil heraus), und $\langle w_1, \ldots, w_j \rangle = \langle v_1, \ldots, v_j \rangle$.

    \textbf{Orthogonalprojektion:} Sei $P_U v = \sum_j (v, w_j) w_j$.
    Dann gilt $(v - P_U v, w_j) = (v, w_j) - (v, w_j) = 0$ für alle $j$,
    also $v - P_U v \perp U$. Die Minimalitätseigenschaft folgt aus
    $\|v - u\|^2 = \|v - P_U v\|^2 + \|P_U v - u\|^2 \geq \|v - P_U v\|^2$.
\end{proof}


\begin{beobachtung}[Fourier als Orthogonalprojektion]
    Das System $[\cos jx, \sin jx \mid j = 0, 1, 2, \dots]$ ist ein
    Orthogonalsystem im Raum der quadratintegrierbaren Funktionen auf
    $[0, 2\pi]$. Die Fourier-Koeffizienten sind genau die
    Orthogonalprojektionen auf diese Trägerelementen.

    Das heißt: Die Fourier-Analyse ist in ihrer Essenz \emph{nichts anderes}
    als die Darstellung eines Vektors (einer Funktion) in einer
    Orthonormalträger -- dasselbe Prinzip, das in $\mathbb{R}^n$ einen
    Vektor in Koordinaten zerlegt.
\end{beobachtung}

% ============================================================
\chapter{Normalformen und Struktursätze}
% ============================================================

\section{Minimalpolynom und Cayley--Hamilton}

Das charakteristische Polynom $\chi_A$ beschreibt die Eigenwertstruktur
einer Matrix. Das Minimalpolynom verfeinert dieses Bild: Es ist das
kleinste Polynom, das $A$ auslöscht -- und es entscheidet über
Diagonalisierbarkeit.

\begin{definition}[Minimalpolynom]
    Sei $V$ ein endlichdimensionaler $K$-Vektorraum und $A \in \mathrm{End}(V)$.
    Das \textbf{Annulatorideal} von $A$ ist
    $I(A) = \{h \in K[x] \mid h(A) = 0\}$.
    Da $K[x]$ ein Hauptidealring ist, wird $I(A)$ von einem eindeutig
    bestimmten normierten Polynom erzeugt: dem
    \textbf{Minimalpolynom} $\mu_A = m_A$.

    Es gilt: $h(A) = 0 \Leftrightarrow m_A \mid h$.
    Jeder Eigenwert von $A$ ist eine Nullstelle von $m_A$.
\end{definition}

\begin{satz}[Cayley--Hamilton]
    Sei $\dim V = n < \infty$ und $A \in \mathrm{End}(V)$ mit
    charakteristischem Polynom $\chi_A$. Dann gilt
    \[
        \chi_A(A) = 0.
    \]
    Insbesondere teilt $m_A$ das Polynom $\chi_A$, und $\deg m_A \leq n$.
\end{satz}
\begin{proof}
    Wir betrachten $A$ als Matrix in $(K)_n$. In $(K[x])_n$ bilden wir
    zur Matrix $xE - A$ die adjungierte Matrix $(xE - A)^\sharp$,
    deren Einträge $(n{-}1)$-reihige Unterdeterminanten von $xE - A$ sind.
    Nach dem Entwicklungssatz gilt $(xE - A)(xE - A)^\sharp = \det(xE-A) \cdot E
    = \chi_A \cdot E$.
    Da $(xE - A)^\sharp = C_0 + C_1 x + \cdots + C_{n-1} x^{n-1}$ mit $C_j \in (K)_n$
    und $\chi_A = a_0 + a_1 x + \cdots + x^n$, liefert Koeffizientenvergleich:
    \[
        a_0 E = -AC_0, \quad a_j E = C_{j-1} - AC_j \; (1 \leq j \leq n-1), \quad E = C_{n-1}.
    \]
    Summiert man $\sum_{j=0}^n A^j(a_j E)$ unter Verwendung dieser Relationen,
    so heben sich alle Terme auf: $\chi_A(A) = 0$.
\end{proof}

\begin{beobachtung}[Was Cayley--Hamilton bedeutet]
    Der Satz sagt: Jede $n \times n$-Matrix erfüllt eine Polynomgleichung
    vom Grad $n$ -- nämlich ihr eigenes charakteristisches Polynom.
    Das ist nicht trivial: Naiv würde man nur erwarten, dass eine
    $n^2$-dimensionale Matrizenalgebra eine Relation vom Grad $\leq n^2$ hat.
    Cayley--Hamilton verschärft dies auf Grad $n$.

    Praktische Folgerung: Für jede invertierbare Matrix $A$ kann man $A^{-1}$
    als Polynom in $A$ ausdrücken, denn aus $\chi_A(A) = 0$ folgt
    $A^{-1} = -\frac{1}{a_0}(a_1 E + a_2 A + \cdots + A^{n-1})$.
\end{beobachtung}

\begin{satz}[Minimalpolynom und Diagonalisierbarkeit]
    Sei $\dim V < \infty$ und $A \in \mathrm{End}(V)$. Dann sind äquivalent:
    \begin{itemize}
        \item[-] $A$ ist \textbf{diagonalisierbar} (es gibt einen Träger aus Eigenvektoren).
        \item[-] Das Minimalpolynom $m_A$ zerfällt in paarweise verschiedene
              Linearfaktoren: $m_A = \prod_{j=1}^r (x - a_j)$ mit $a_j \in K$ paarweise verschieden.
    \end{itemize}
\end{satz}
\begin{proof}
    Ist $V = \bigoplus_{j=1}^r \ker(A - a_j E)$, so gilt $\prod_j (A - a_j E) = 0$
    auf ganz $V$, also $m_A \mid \prod_j (x - a_j)$. Da umgekehrt jeder Eigenwert
    eine Nullstelle von $m_A$ ist, folgt $m_A = \prod_j (x - a_j)$ (normiert,
    paarweise verschieden). Ist umgekehrt $m_A = \prod_j (x - a_j)$, so folgt
    aus dem Chinesischen Restsatz für $K[x]$-Moduln:
    $V = \bigoplus_j \ker(A - a_j E)$.
\end{proof}

\begin{beobachtung}[Minimalpolynom als feinstes Strukturmerkmal]
    Das Minimalpolynom und das charakteristische Polynom tragen
    komplementäre Informationen:
    \begin{itemize}
        \item[-] $\chi_A$ gibt die algebraischen Vielfachheiten der Eigenwerte.
        \item[-] $m_A$ gibt die Größen der größten Jordankästchen zu jedem Eigenwert:
              Die Potenz $(x - a)^k$ in $m_A$ ist genau die Größe des größten
              Jordankästchens zu $a$.
    \end{itemize}
    Ein Eigenwert $a$ hat geometrische Vielfachheit $1$ genau dann, wenn
    $(x-a)^{m_a}$ das größte Jordankästchen hat (wobei $m_a$ die algebraische
    Vielfachheit ist). $A$ ist diagonalisierbar genau dann, wenn alle
    Jordankästchen $1 \times 1$ sind -- was $m_A$ quadratfrei macht.
\end{beobachtung}

% ============================================================
\section{Trigonalisierung}
% ============================================================

\begin{satz}[Trigonalisierbarkeit]
    Sei $\dim V < \infty$ und $A \in \mathrm{End}(V)$. Zerfällt $\chi_A$
    vollständig in Linearfaktoren (was über $\mathbb{C}$ stets der Fall ist),
    so gibt es einen Träger von $V$, bezüglich derer $A$ eine
    \textbf{obere Dreiecksmatrix} hat. Die Diagonaleinträge sind die Eigenwerte
    (mit algebraischen Vielfachheiten).
\end{satz}
\begin{proof}
    Durch Induktion nach $n = \dim V$. Zerfällt $\chi_A$ vollständig,
    so hat $A$ mindestens einen Eigenwert $a_1$ mit Eigenvektor $v_1$.
    Ergänze $v_1$ zu einem Träger; bezüglich dieser hat $A$ die Form
    $\bigl(\begin{smallmatrix} a_1 & * \\ 0 & A' \end{smallmatrix}\bigr)$.
    Auf $A'$ mit $\dim = n-1$ wende man die Induktionsannahme an.
\end{proof}

\begin{beobachtung}[Schur-Zerlegung und numerische Praxis]
    Die Trigonalisierung ist die Grundlage der \textbf{Schur-Zerlegung}:
    Über $\mathbb{C}$ gibt es für jede Matrix $A$ eine unitäre Matrix $U$
    mit $U^{-1}AU$ obere Dreiecksmatrix. Dies ist der
    ausgangspunkt für numerische Eigenwertverfahren (QR-Algorithmus).

    Im Gegensatz zur Jordanform ist die Schur-Zerlegung numerisch stabil:
    Die unitäre Transformation erhält die euklidische Norm, was
    Rundungsfehler kontrollierbar macht. Die meisten numerischen
    Eigenwertberechnungen (LAPACK, MATLAB) basieren auf der Schur-Zerlegung,
    nicht auf der Jordanform.
\end{beobachtung}

% ============================================================
\section{Normierte Vektorräume und Konvergenz}
% ============================================================

\begin{definition}[Normierter Raum]
    Eine \textbf{Norm} auf einem Vektorraum $V$ ist eine Abbildung
    $\|\cdot\| : V \to \mathbb{R}_{\ge 0}$ mit:
    \begin{itemize}
        \item[-] $\|v\| = 0 \Leftrightarrow v = 0$
        \item[-] $\|kv\| = |k| \cdot \|v\|$
        \item[-] $\|u + v\| \le \|u\| + \|v\|$ \quad (\textbf{Dreiecksungleichung})
    \end{itemize}
\end{definition}

\begin{satz}[Matrixexponential und dynamische Systeme]
    Für $A \in K^{n \times n}$ ist die \textbf{Matrixexponentialfunktion}
    \[
        e^{At} = \sum_{k=0}^{\infty} \frac{(At)^k}{k!}
    \]
    absolut konvergent für alle $t \in \mathbb{R}$. Die Lösung des linearen
    Differentialgleichungssystems $\dot{x} = Ax$ mit Anfangsbedingung $x(0) = x_0$
    ist genau $x(t) = e^{At} x_0$.
\end{satz}
\begin{proof}
    Für das System $\dot{x} = Ax$ mit $x(0) = x_0$ setze $y(t) = e^{tA} x_0$.
    Dann gilt $\dot{y}(t) = A e^{tA} x_0 = A y(t)$ (nach den Eigenschaften
    des Matrixexponentials) und $y(0) = e^0 x_0 = x_0$.
    \textbf{Eindeutigkeit:} Ist $z(t)$ eine weitere Lösung mit $z(0) = x_0$,
    so gilt $\frac{d}{dt}(e^{-tA} z) = -A e^{-tA} z + e^{-tA} Az = 0$,
    also ist $e^{-tA} z(t)$ konstant $= z(0) = x_0$, d.\,h.\ $z(t) = e^{tA} x_0$.
    Die Stabilität folgt: Alle Eigenwerte von $A$ haben negativen Realteil
    genau dann, wenn $e^{tA} \to 0$ für $t \to \infty$.
\end{proof}


\begin{beobachtung}[Die Matrix-e-Funktion verbindet alles]
    Die Matrixexponentialfunktion $e^{At}$ ist die direkte Verallgemeinerung
    der skalaren e-Funktion $e^{at}$. Ist $A$ diagonalisierbar mit Eigenwerten
    $\lambda_1, \dots, \lambda_n$, so gilt:
    \[
        e^{At} = P \,\mathrm{diag}(e^{\lambda_1 t}, \dots, e^{\lambda_n t})\, P^{-1}.
    \]
    Jede Lösung des Systems ist eine Superposition von $n$ skalaren
    e-Funktionen $e^{\lambda_j t}$ -- genau die Eigenwerte entscheiden
    über Wachstum, Zerfall und Oszillation, wie in der Signaltheorie der
    Realteil $\sigma$ und der Imaginärteil $\omega$ von $s = \sigma + i\omega$.

    Dies ist der tiefste Zusammenhang: \textbf{Lineare Algebra und
    Differentialgleichungen sind zwei Seiten derselben Struktur}.
\end{beobachtung}

% ============================================================

\section{Die Frobenius-Normalform}

\begin{definition}[Begleitmatrix und Frobenius-Normalform]
    Sei $f(x) = x^n + a_{n-1}x^{n-1} + \cdots + a_0 \in K[x]$.
    Die \textbf{Begleitmatrix} (Companion Matrix) von $f$ ist:
    \[
        C(f) = \begin{pmatrix}
            0 & 0 & \cdots & 0 & -a_0 \\
            1 & 0 & \cdots & 0 & -a_1 \\
            0 & 1 & \cdots & 0 & -a_2 \\
            \vdots & & \ddots & & \vdots \\
            0 & 0 & \cdots & 1 & -a_{n-1}
        \end{pmatrix} \in K^{n \times n}.
    \]
    Das charakteristische Polynom von $C(f)$ ist genau $f$.
    Die \textbf{Frobenius-Normalform} einer Matrix $A$ ist eine
    Blockdiagonalmatrix aus Begleitmatrizen:
    $F = \mathrm{diag}(C(f_1), \ldots, C(f_r))$, wobei
    $f_1 \mid f_2 \mid \cdots \mid f_r$ die Invarianzfaktoren von $A$ sind.
\end{definition}

\begin{beobachtung}[Frobenius- vs. Jordan-Normalform]
    Die Frobenius-Normalform existiert über \emph{jedem} Körper $K$
    (ohne Voraussetzung, dass $\chi_A$ in Linearfaktoren zerfällt),
    während die Jordan-Normalform nur über algebraisch abgeschlossenen
    Körpern garantiert ist.

    Ãœber $\mathbb{Q}$: Die Matrix $A = \bigl(\begin{smallmatrix}0&-1\\1&0\end{smallmatrix}\bigr)$
    hat $\chi_A = x^2 + 1$, das über $\mathbb{Q}$ irreduzibel ist.
    Ihre Jordan-Form existiert nur über $\mathbb{C}$, ihre Frobenius-Form
    ist $C(x^2+1) = A$ selbst -- bereits in Frobenius-Normalform.
\end{beobachtung}

\chapter{Strukturprinzipien der linearen Algebra}
% ============================================================

\begin{beobachtung}[Das Schwerpunkt-Prinzip als leitendes Ziel]
    Die lineare Algebra lässt sich aus einer einzigen Leitfrage verstehen:
    \textbf{Wie findet man in einer algebraischen Struktur einen Schwerpunkt,
    an dem sie sich am effizientesten lösen lässt?}

    Strukturen mit Schwerpunkt erlauben komponentenweise, lokale Operationen:
    Addition und Multiplikation wirken auf Einzelelemente, ohne globale
    Abhängigkeiten. Der Vektor ist das Urmodell dieser Eigenschaft.
    Die Matrix dagegen hat keinen natürlichen Schwerpunkt -- ihre
    Multiplikation koppelt ganze Zeilen mit ganzen Spalten.

    Die gesamte Theorie der Normalformen, der Diagonalisierung und der
    Projektionen ist im Kern die systematische Suche nach Transformationen,
    die einer Matrix einen Schwerpunkt geben:
    \begin{itemize}
        \item[-] Die \textbf{Diagonalisierung} gibt der Matrix den Schwerpunkt
              des Vektors -- Eigenwerte als Skalare, komponentenweise anwendbar.
        \item[-] Die \textbf{Projektion} isoliert Unterräume und macht
              die Zerlegung $V = \mathrm{Im}\,P \oplus \ker P$ zum lokalen
              Werkzeug.
        \item[-] Die \textbf{Orthonormalträger} entkoppelt alle Koordinaten
              vollständig -- das höchste Maß an Schwerpunktstruktur in
              einem Vektorraum.
    \end{itemize}
    Das Ziel ist stets: Die globale, verflochtene Struktur der Matrix
    in eine lokale, komponentenweise Struktur überführen -- sie in
    ihrem Schwerpunkt lösen.
\end{beobachtung}

\begin{beobachtung}[Das Programm der linearen Algebra]
    Die lineare Algebra verfolgt ein einziges großes Programm: Sie sucht
    nach der \textbf{kanonischen Form} -- der einfachsten Darstellung einer
    linearen Abbildung, die alle wesentlichen Informationen sichtbar macht.
    \begin{itemize}
        \item[-] \textbf{Diagonalform}: Wenn die Eigenvektoren einen Träger bilden.
              Die Abbildung wirkt nur als Skalierung.
        \item[-] \textbf{Jordanform}: Im allgemeinen Fall über $\mathbb{C}$.
              Zerfällt in Jordanblöcke, die Nilpotenz und Eigenwert vereinen.
        \item[-] \textbf{Normalform}: Für Matrizen mit Skalarprodukt
              (hermitesch, unitär, normal). Die ONB-Diagonalisierbarkeit ist
              hier garantiert.
    \end{itemize}
    In jedem Fall ist das Ziel dasselbe: die Abbildung in einen Träger zu
    transformieren, in der sie maximal transparent wird.
\end{beobachtung}

\begin{beobachtung}[Algebraische Strukturen als Sprache]
    Die Hierarchie Gruppe $\to$ Ring $\to$ Körper $\to$ Vektorraum ist keine
    Taxonomie um ihrer selbst willen. Sie ist eine \textbf{Sprache}, die
    immer feinere Strukturen beschreibt:
    \begin{itemize}
        \item[-] Eine Gruppe beschreibt Symmetrie und Invertierbarkeit.
        \item[-] Ein Ring beschreibt Addition und Multiplikation (ohne Division).
        \item[-] Ein Körper ist ein Ring mit vollständiger Arithmetik.
        \item[-] Ein Vektorraum ist eine abelsche Gruppe mit skalarer Streckung.
    \end{itemize}
    Jede neue Struktur fügt axiomatisch kontrollierte Operationen hinzu --
    nicht mehr als nötig, aber genug, um starke Sätze zu beweisen.
\end{beobachtung}

\begin{beobachtung}[Abbildungen sind das eigentliche Objekt]
    Die tiefste Einsicht der strukturellen Algebra: Nicht Objekte, sondern
    \textbf{Abbildungen} zwischen ihnen tragen die wesentliche Information.
    Ein Vektorraum wird nicht durch seine Elemente, sondern durch seine
    linearen Abbildungen verstanden.

    Die Dimension entscheidet über Isomorphie -- alle $n$-dimensionalen
    Räume über $K$ sind isomorph zu $K^n$. Damit ist die Frage nach dem
    \emph{Raum selbst} gelöst; was bleibt, ist die Frage nach den
    Abbildungen zwischen Räumen: deren Kern, Bild, Normalform.
\end{beobachtung}

\begin{satz}[Zentrales Resultat: Rangsatz]
    Sei $f : V \to W$ eine lineare Abbildung zwischen endlichdimensionalen
    Räumen. Dann gilt:
    \[
        \dim V = \dim \ker f + \dim \mathrm{Im}\, f = \dim \ker f + \mathrm{Rang}\, f.
    \]
    Dieser Satz ist das strukturelle Herzstück der linearen Algebra. Er
    sagt: Jede lineare Abbildung zerlegt ihren Definitionsbereich in zwei
    Anteile -- den Teil, der vernichtet wird (Kern), und den Teil, der
    abgebildet wird (Komplement des Kerns). Die Summe der Dimensionen
    ist invariant.
\end{satz}
\begin{proof}
    Sei $[u_1, \ldots, u_r]$ ein Träger von $\ker f$
    und ergänze zu einem Träger
    und ergänze zu einem Träger $[u_1, \ldots, u_r,
    v_1, \ldots, v_s]$ von $V$.
    \textbf{Erzeugung:} Für beliebiges $f(w) \in \mathrm{Im}\, f$ schreibe
    $w = \sum_i a_i u_i + \sum_j b_j v_j$; dann $f(w) = \sum_j b_j f(v_j)$.

    \textbf{Lineare Unabhängigkeit:} Sei $\sum_j b_j f(v_j) = 0$, also
    $f(\sum_j b_j v_j) = 0$. Dann ist $\sum_j b_j v_j \in \ker f$,
    also $\sum_j b_j v_j = \sum_i c_i u_i$. Da alle $u_i, v_j$ zusammen
    einen Träger bilden, folgt $b_j = 0$ für alle $j$.

    Also $\dim V = r + s = \dim \ker f + \dim \mathrm{Im}\, f$.
\end{proof}


\begin{beobachtung}[Verbindung zur Signaltheorie]
    In der Signaltheorie ist die Matrix $A$ die Systemmatrix, $e^{At}$
    die Systemantwort, und die Eigenwerte $\lambda_j = \sigma_j + i\omega_j$
    die Pole des Systems. Der Rangsatz sagt: Die Dimension des Zustandsraums
    ist die Summe aus der Dimension des nichterregbaren Unterraums (Kern) und
    der Dimension des erregbaren Unterraums (Bild).

    Damit schließt sich der Kreis: Lineare Algebra ist die
    \textbf{Grammatik} aller linearen dynamischen Systeme -- die Sprache,
    in der Signaltheorie, Regelungstechnik und Quantenmechanik formuliert sind.
\end{beobachtung}

% ============================================================
% ============================================================


\section{Das Träger-Prinzip als Leitfaden}

\begin{beobachtung}[Warum Träger mehr sagen als Erzeugendensysteme]
    Ein Erzeugendensystem erlaubt redundante Darstellungen. Erst die
    lineare Unabhängigkeit -- das zweite Merkmal eines Trägers -- sorgt
    für Eindeutigkeit: Jeder Vektor besitzt genau \emph{eine}
    Koordinatendarstellung. Diese Eindeutigkeit ist der Grund, warum
    Träger so mächtig sind: Sie machen die Struktur des Vektorraums
    vollständig transparent und berechnbar.

    Der Begriff des Trägers ersetzt in dieser Arbeit die übliche
    Bezeichnung ``Träger'' und betont damit den projektiven Charakter:
    Ein Träger ist die Menge der Richtungen, auf die ein Raum projiziert
    wird -- die Menge der \emph{unabhängigen Projektionsziele}.
\end{beobachtung}

\begin{beobachtung}[Dimension als invariante Größe]
    Der Steinitzsche Austauschsatz zeigt, dass alle Träger eines
    endlichdimensionalen Vektorraums dieselbe Kardinalität haben.
    Diese Kardinalität -- die Dimension -- ist der einzige invariante
    Parameter eines Vektorraums bis auf Isomorphie:
    Zwei endlichdimensionale Vektorräume über demselben Körper sind
    isomorph genau dann, wenn sie dieselbe Dimension haben.

    Die Dimension ist daher die vollständige algebraische Invariante
    der Vektorraumstruktur. Alle weiteren Strukturen -- Skalarprodukte,
    Normen, lineare Abbildungen -- sind Zusatzstrukturen auf dem
    Vektorraum, die durch ihre Wahl auf einem Träger beschrieben werden.
\end{beobachtung}

\begin{beobachtung}[Lineare Abbildungen als strukturerhaltende Karten]
    Eine lineare Abbildung $f : V \to W$ ist vollständig durch ihre
    Werte auf einem Träger von $V$ bestimmt. Der Rangsatz
    $\dim\ker f + \dim\mathrm{Im}\,f = \dim V$ quantifiziert den
    Informationsgehalt von $f$: Was verloren geht (Kern), was
    erhalten bleibt (Bild).

    Der Homomorphiesatz $\mathrm{Im}\,f \cong V/\ker f$ sagt:
    Die Bilder der Trägervektoren, die nicht im Kern liegen, bilden
    einen Träger des Bildes. Lineare Abbildungen respektieren die
    Trägerstruktur -- das ist ihre definierende Eigenschaft.
\end{beobachtung}

\chapter{Das Schwerpunkt-Prinzip}
% ============================================================

Das vorangegangene Kapitel hat das Schwerpunkt-Prinzip als leitendes
Ziel der linearen Algebra benannt. Dieses Kapitel macht diese Idee
mathematisch präzise: Es beweist formal, dass Strukturen in ihrem
Schwerpunkt hinsichtlich einer Verknüpfung maximal effizient lösbar
sind, und belegt dies durch ausführlich analysierte Beispiele aus
der linearen Algebra und der Algorithmik.

% ============================================================
\section{Grundbegriffe}
% ============================================================

\begin{definition}[Algebraische Struktur und Verknüpfung]
    Eine \textbf{algebraische Struktur} ist ein Tupel $(S, \circ_1, \ldots, \circ_k)$
    bestehend aus einer Menge $S$ und Verknüpfungen
    $\circ_j : S \times S \to S$, die vorgegebene Axiome erfüllen.

    Eine Verknüpfung $\circ$ auf $S$ heißt \textbf{komponentenweise},
    falls es eine Zerlegung $S \cong S_1 \times \cdots \times S_n$
    gibt, so dass
    \[
        (a_1, \ldots, a_n) \circ (b_1, \ldots, b_n)
        = (a_1 \circ_1 b_1,\; \ldots,\; a_n \circ_n b_n)
    \]
    mit Verknüpfungen $\circ_j$ auf $S_j$.
    Die Struktur ist \textbf{im Schwerpunkt}, falls ihre Verknüpfung
    komponentenweise ist.
\end{definition}

\begin{definition}[Schwerpunkt einer Struktur]
    Sei $(S, \circ)$ eine algebraische Struktur. Ein
    \textbf{Schwerpunkt} von $S$ bezüglich $\circ$ ist eine Darstellung
    \[
        S \cong S_1 \times \cdots \times S_n
    \]
    so dass $\circ$ komponentenweise wirkt und jedes $S_j$ bezüglich
    $\circ_j$ atomar (nicht weiter zerlegbar) ist.

    Formal: $S$ ist genau dann in einem Schwerpunkt, wenn ein
    Isomorphismus $\varphi : S \to \prod_j S_j$ existiert mit
    \[
        \varphi(a \circ b) = \varphi(a) \circ' \varphi(b),
    \]
    wobei $\circ'$ komponentenweise ist. Den Isomorphismus $\varphi$
    nennen wir die \textbf{Schwerpunkttransformation}.
\end{definition}

\begin{beobachtung}[Intuition des Schwerpunkts]
    Eine Struktur im Schwerpunkt bedeutet: Jede globale Operation
    zerfällt in $n$ voneinander unabhängige lokale Operationen.
    Der Aufwand reduziert sich von einer globalen Berechnung
    (Aufwand $T_{\mathrm{global}}$) auf $n$ unabhängige lokale
    Berechnungen (Aufwand $\sum_j T_j$). Die Unabhängigkeit der
    Komponenten ist der entscheidende Gewinn.

    \textbf{Beispiel:} Die Addition zweier $n$-di\-men\-sio\-na\-ler Vektoren
    ist komponentenweise:
    \[
        (a_1,\ldots,a_n) + (b_1,\ldots,b_n) = (a_1+b_1, \ldots, a_n+b_n).
    \]
    addiert -- $n$ Additionen, keine globale Abhängigkeit. Dies ist
    der Schwerpunkt der Vektorstruktur.

    \textbf{Gegenbeispiel:} Die Multiplikation zweier $n \times n$-Matrizen
    ist \emph{nicht} komponentenweise: $(AB)_{ij} = \sum_k a_{ik} b_{kj}$
    koppelt alle Einträge einer Zeile von $A$ mit einer Spalte von $B$.
    Eine Matrix ist nicht in ihrem Schwerpunkt -- sie muss erst
    (durch Diagonalisierung) in einen solchen transformiert werden.
\end{beobachtung}

% ============================================================
\section{Schwerpunkt und Effizienz}
% ============================================================

Um den Effizienzbegriff zu präzisieren, führen wir ein allgemeines
Kostenmodell ein.

\begin{definition}[Kostenfunktion und Berechnungsaufwand]
    Sei $(S, \circ)$ eine algebraische Struktur. Eine
    \textbf{Kostenfunktion} ist eine Abbildung $c : S \times S \to \mathbb{R}_{\geq 0}$,
    die jedem Paar $(a, b)$ die Kosten der Berechnung von $a \circ b$ zuordnet.

    Der \textbf{Gesamtaufwand} für die Auswertung von $a \circ b$ auf
    einer Zerlegung $S \cong \prod_{j=1}^n S_j$ ist
    \[
        T_{\mathrm{zerlegt}}(a, b)
        = \sum_{j=1}^n c_j(a_j, b_j) + T_{\mathrm{Transformation}},
    \]
    wobei $c_j$ die Kosten auf $S_j$ und $T_{\mathrm{Transformation}}$
    die Kosten der Schwerpunkttransformation $\varphi$ und ihrer
    Inversen sind.

    Der \textbf{direkte Aufwand} ohne Zerlegung ist
    $T_{\mathrm{direkt}}(a, b) = c(a, b)$.
\end{definition}

\begin{satz}[Schwerpunkt-Prinzip: Formale Aussage]
    \label{satz:schwerpunkt}
    Sei $(S, \circ)$ eine algebraische Struktur mit Schwerpunkttransformation
    $\varphi : S \to \prod_{j=1}^n S_j$.
    Sei die Kostenfunktion $c$ superadditiv in dem Sinne, dass
    \[
        c(a, b) \geq \sum_{j=1}^n c_j(\varphi(a)_j,\, \varphi(b)_j)
        \qquad \text{für alle } a, b \in S.
    \]
    Dann gilt: Sobald $T_{\mathrm{Transformation}} = o(c(a,b))$
    (die Transformationskosten sind asymptotisch vernachlässigbar
    gegenüber den direkten Kosten), ist die Berechnung im
    Schwerpunkt effizienter:
    \[
        T_{\mathrm{zerlegt}}(a, b) < T_{\mathrm{direkt}}(a, b)
        \quad \text{für hinreichend große } n.
    \]
\end{satz}

\begin{proof}
    Nach Voraussetzung gilt:
    \[
        T_{\mathrm{zerlegt}}(a,b)
        = \sum_{j=1}^n c_j(\varphi(a)_j, \varphi(b)_j)
        + T_{\mathrm{Transformation}}.
    \]
    Da $c$ superadditiv ist, gilt
    $\sum_{j} c_j(\varphi(a)_j, \varphi(b)_j) \leq c(a,b)
    = T_{\mathrm{direkt}}(a,b)$.

    Damit folgt:
    \[
        T_{\mathrm{zerlegt}}(a,b)
        \leq T_{\mathrm{direkt}}(a,b) + T_{\mathrm{Transformation}}.
    \]
    Da $T_{\mathrm{Transformation}} = o(c(a,b))$, gibt es für jedes
    $\varepsilon > 0$ ein $N_0$, so dass für $n \geq N_0$ gilt:
    $T_{\mathrm{Transformation}} \leq \varepsilon \cdot c(a,b)$.

    Ist die Ungleichung der Superadditivität strikt
    (d.\,h.\ $c(a,b) > \sum_j c_j(\varphi(a)_j, \varphi(b)_j) + \delta$
    für ein $\delta > 0$), so folgt
    \[
        T_{\mathrm{zerlegt}}(a,b)
        \leq T_{\mathrm{direkt}}(a,b) - \delta
        + T_{\mathrm{Transformation}}
        < T_{\mathrm{direkt}}(a,b)
    \]
    für $T_{\mathrm{Transformation}} < \delta$,
    was für hinreichend große $n$ gilt.
\end{proof}

\begin{beobachtung}[Bedeutung des Schwerpunkt-Satzes]
    Der Satz quantifiziert, wann die Suche nach dem Schwerpunkt
    rechnerisch lohnt: Die Transformationskosten müssen klein sein
    gegenüber dem Gewinn durch Zerlegung. In den folgenden Beispielen
    werden beide Terme explizit berechnet und verglichen.

    Das Schwerpunkt-Prinzip ist kein Selbstzweck, sondern das
    mathematische Rückgrat einer Reihe fundamentaler Algorithmen:
    Schnelle Fourier-Transformation, Divide-and-Conquer-Sortierverfahren,
    Matrixdiagonalisierung, Primärkomponentenzerlegung -- sie alle
    suchen den Schwerpunkt ihrer jeweiligen Struktur.
\end{beobachtung}

% ============================================================
\section{Beispiel: Vektoren}
% ============================================================

\subsection*{Addition: natürlicher Schwerpunkt}

Die Vektoraddition ist der einfachste Fall: Die Struktur $(K^n, +)$ liegt bereits im Schwerpunkt, kein Trägerwechsel ist nötig. Algorithmus~\ref{alg:vecadd} beschreibt dies explizit.


\begin{algorithm}[H]
\DontPrintSemicolon
\caption{Vektoraddition (im Schwerpunkt)}\label{alg:vecadd}
\KwIn{Vektoren $a = (a_1,\ldots,a_n)^t,\; b = (b_1,\ldots,b_n)^t \in K^n$}
\KwOut{$c = a + b \in K^n$}
\For{$j \leftarrow 1$ \KwTo $n$}{
    $c_j \leftarrow a_j + b_j$ \tcp*{unabhängige Skalaraddition}
}
\Return{$c$}\;
\BlankLine
\textbf{Laufzeit:} $T(n) = n$ Additionen \quad (optimal, untere Schranke $= n$)
\end{algorithm}


\begin{satz}[Vektoraddition im Schwerpunkt]
    Sei $V = K^n$ mit komponentenweiser Addition. Dann ist
    die Addition \textbf{bereits im Schwerpunkt}: Die Identität
    $\varphi = \mathrm{id}$ ist die Schwerpunkttransformation,
    und die Berechnungskosten betragen exakt $T(n) = n$ Operationen.

    Genauer: Für $a = (a_1, \ldots, a_n)^t$ und $b = (b_1, \ldots, b_n)^t$
    gilt $a + b = (a_1 + b_1, \ldots, a_n + b_n)^t$, und jede Komponente
    $a_j + b_j$ ist unabhängig berechenbar.
\end{satz}

\begin{proof}
    Wir zeigen, dass keine Zerlegung die Additionskosten unterschreiten kann.

    \textbf{Untere Schranke:} Jede Komponente $(a+b)_j = a_j + b_j$
    hängt von den Eingaben $a_j$ und $b_j$ ab und von keiner anderen.
    Jede korrekte Additionsberechnung muss daher mindestens $n$ atomare
    Additionen durchführen -- eine pro Komponente.
    Also $T(n) \geq n$.

    \textbf{Obere Schranke:} Die komponentenweise Addition berechnet
    $a + b$ in genau $n$ atomaren Operationen.
    Also $T(n) \leq n$.

    Damit $T(n) = n$, und der Algorithmus ist optimal.
    Die Transformationskosten sind $T_{\mathrm{Transformation}} = 0$,
    da kein Trägerwechsel nötig ist: Der Vektor ist bereits im Schwerpunkt.
\end{proof}

\begin{beobachtung}[Warum der Vektor schon im Schwerpunkt ist]
    Die komponentenweise Addition des Vektors ist das Paradebeispiel
    einer Struktur im Schwerpunkt: Jede Komponente $j$ ist eine
    unabhängige Kopie des Skalarkörpers $(K, +)$. Es gibt keine
    Kopplung zwischen verschiedenen Komponenten.

    Formal: Die Zerlegung $K^n \cong K \times K \times \cdots \times K$
    ist ein Ringisomorphismus bezüglich komponentenweiser Addition.
    Die Atome $S_j = K$ sind irreduzibel (kein weiterer Zerlegungsschritt
    möglich ohne den Körper aufzubrechen).
\end{beobachtung}

\subsection*{Multiplikation: Schwerpunkt durch Diagonalisierung}

Matrizenmultiplikation ist \emph{nicht} komponentenweise. Algorithmus~\ref{alg:matvec} zeigt das direkte Verfahren; Algorithmus~\ref{alg:schwerpunkt} das Schwerpunkt-Verfahren nach Diagonalisierung.


\begin{algorithm}[H]
\DontPrintSemicolon
\caption{Matrix-Vektor-Produkt (direkt)}\label{alg:matvec}
\KwIn{$A \in K^{n \times n}$, $v \in K^n$}
\KwOut{$w = Av \in K^n$}
\For{$i \leftarrow 1$ \KwTo $n$}{
    $w_i \leftarrow 0$\;
    \For{$j \leftarrow 1$ \KwTo $n$}{
        $w_i \leftarrow w_i + a_{ij} \cdot v_j$ \tcp*{globale Kopplung aller Zeilen}
    }
}
\Return{$w$}\;
\BlankLine
\textbf{Laufzeit:} $T(n) = n^2$ Multiplikationen $+\; n(n-1)$ Additionen $= \Theta(n^2)$
\end{algorithm}

\begin{algorithm}[H]
\DontPrintSemicolon
\caption{Matrix-Vektor-Produkt im Schwerpunkt}\label{alg:schwerpunkt}
\KwIn{$P, D = \mathrm{diag}(\lambda_1,\ldots,\lambda_n), P^{-1}$ (vorberechnet); Vektor $v \in K^n$}
\KwOut{$w = Av = PDP^{-1}v$}
\tcp{Schritt 1: Koordinatenwechsel in den Schwerpunkt}
$\tilde{v} \leftarrow P^{-1} v$ \tcp*{$\Theta(n^2)$ Operationen}
\BlankLine
\tcp{Schritt 2: Komponentenweise Skalierung (Schwerpunkt!)}
\For{$j \leftarrow 1$ \KwTo $n$}{
    $\tilde{w}_j \leftarrow \lambda_j \cdot \tilde{v}_j$ \tcp*{$n$ unabhängige Multiplikationen}
}
\BlankLine
\tcp{Schritt 3: Rücktransformation}
$w \leftarrow P\, \tilde{w}$ \tcp*{$\Theta(n^2)$ Operationen}
\Return{$w$}\;
\BlankLine
\textbf{Laufzeit pro Anwendung:} $\Theta(n^2) + \Theta(n) + \Theta(n^2) = \Theta(n^2)$\;
\textbf{Laufzeit für $A^k v$:} Schritt 2 mit $\lambda_j^k$: $\Theta(n^2)$ unabhängig von $k$
\end{algorithm}


\begin{satz}[Vektormultiplikation im Schwerpunkt nach Diagonalisierung]
    \label{satz:diag-schwerpunkt}
    Sei $A \in K^{n \times n}$ diagonalisierbar mit Eigenwerten
    $\lambda_1, \ldots, \lambda_n$ und Eigenvektor-Träger
    $P = [v_1 \mid \cdots \mid v_n]$.

    \textbf{Direkter Aufwand} (Matrix-Vektor-Produkt ohne Schwerpunkt):
    \[
        T_{\mathrm{direkt}}(n) = 2n^2 - n \;=\; \Theta(n^2) \text{ Operationen}.
    \]

    \textbf{Aufwand im Schwerpunkt} (nach Diagonalisierung,
    für $m$ Anwendungen auf denselben $A$):
    \[
        T_{\mathrm{Schwerpunkt}}(n, m)
        = \underbrace{2n^2}_{\text{Vorber. } P^{-1}}
        + \underbrace{n \cdot m}_{\text{Eigenwert-Skalierung}}
        + \underbrace{2n^2 \cdot m}_{\text{Rücktransf. } P}
        = \Theta(n^2) + m \cdot \Theta(n).
    \]

    Für $m \gg n$ ist der Schwerpunkt-Algorithmus asymptotisch
    effizienter: $T_{\mathrm{Schwerpunkt}} = \Theta(n^2 + mn)$
    gegenüber $T_{\mathrm{direkt}} = \Theta(mn^2)$.
\end{satz}

\begin{proof}
    \textbf{Direkter Aufwand:} Das Matrix-Vektor-Produkt $Av$ mit
    $A \in K^{n \times n}$ und $v \in K^n$ erfordert für jeden
    der $n$ Einträge $(Av)_i = \sum_{j=1}^n a_{ij} v_j$ genau
    $n$ Multiplikationen und $n-1$ Additionen. Über alle $n$ Einträge:
    $n \cdot n + n \cdot (n-1) = 2n^2 - n$ Operationen.
    Also $T_{\mathrm{direkt}}(n) = 2n^2 - n = \Theta(n^2)$.

    \textbf{Schwerpunkt-Algorithmus:} Sei $A = P D P^{-1}$ die
    Diagonalzerlegung mit $D = \mathrm{diag}(\lambda_1, \ldots, \lambda_n)$.
    Für die Berechnung von $Av$ geht man in drei Schritten vor:

    \begin{enumerate}
        \item \textbf{Koordinatenwechsel}: $\tilde{v} = P^{-1} v$.
              Aufwand: $2n^2 - n = \Theta(n^2)$ Operationen
              (ein Matrix-Vektor-Produkt).
        \item \textbf{Skalierung im Schwerpunkt}: $\tilde{w}_j = \lambda_j \tilde{v}_j$
              für $j = 1, \ldots, n$.
              Aufwand: $n$ Multiplikationen $= \Theta(n)$ Operationen.
              \emph{Dies ist der Schwerpunkt: komponentenweise, unabhängig.}
        \item \textbf{Rücktransformation}: $w = P \tilde{w}$.
              Aufwand: $2n^2 - n = \Theta(n^2)$ Operationen.
    \end{enumerate}

    Einmaliger Gesamtaufwand: $(2n^2 - n) + n + (2n^2 - n) = 4n^2 - n = \Theta(n^2)$.

    Für $m$ Wiederholungen mit demselben $A$ (aber verschiedenen $v$):
    Die Diagonalisierung $A = PDP^{-1}$ und die Speicherung von
    $P$, $D$, $P^{-1}$ erfordert einmalig $\Theta(n^3)$ Vorberechnungsaufwand
    (Eigenwertzerlegung). Danach kostet jede Anwendung auf neues $v$:
    \begin{itemize}
        \item[-] $P^{-1}v$: $\Theta(n^2)$
        \item[-] Skalierung $D\tilde{v}$: $\Theta(n)$
        \item[-] $P\tilde{w}$: $\Theta(n^2)$
    \end{itemize}
    Gesamtkosten für $m$ Anwendungen:
    $\Theta(n^3) + m \cdot \Theta(n^2) = \Theta(n^3 + mn^2)$.

    Ohne Diagonalisierung: $m \cdot \Theta(n^2) = \Theta(mn^2)$.

    Der Bruch-even-Punkt liegt bei $m = \Theta(n)$: Für $m \geq cn$
    (mit einer Konstante $c > 0$) gilt
    $\Theta(n^3 + mn^2) = \Theta(mn^2) = \Theta(n^2) \cdot m$,
    also ist der Vorfaktor derselbe. Der Schwerpunkt-Vorteil liegt
    in Schritt 2: Statt $\Theta(n^2)$ pro Anwendung nur $\Theta(n)$.

    \textbf{Potenzberechnung} $A^k v$ macht den Gewinn drastisch:
    \begin{itemize}
        \item[-] Direkt: $k$ Matrix-Vektor-Produkte $= \Theta(kn^2)$.
        \item[-] Im Schwerpunkt: $A^k v = P D^k P^{-1} v$
              $D^k = \mathrm{diag}(\lambda_1^k, \ldots, \lambda_n^k)$.
              Vorberechnungsaufwand: $\Theta(n^2)$ einmalig.
              Anwendung: $\Theta(n^2)$ für $P^{-1}v$, $\Theta(n)$
              für $D^k \tilde{v}$, $\Theta(n^2)$ für $Pw$.
              Gesamt: $\Theta(n^2)$ unabhängig von $k$.
    \end{itemize}
    Die Potenzberechnung $A^k v$ kostet im Schwerpunkt $\Theta(n^2)$
    für beliebiges $k$ -- direktes Rechnen kostet $\Theta(kn^2)$.
    Der Schwerpunkt reduziert die Abhängigkeit von $k$ auf eine
    Konstante: \emph{maximale Effizienz im Schwerpunkt}.
\end{proof}

\begin{beobachtung}[Schwerpunkt-Prinzip bei der Matrixexponentialfunktion]
    Die Effizienz des Schwerpunkts zeigt sich am deutlichsten bei
    der Matrixexponentialfunktion $e^{tA}$: Direkt als Potenzreihe
    berechnet, sind $N$ Matrixpotenzen $A^k$ nötig -- Aufwand
    $\Theta(Nn^3)$. Im Schwerpunkt (Eigenvektorzerlegung) gilt
    $e^{tA} = P \cdot \mathrm{diag}(e^{t\lambda_j}) \cdot P^{-1}$:
    Die $n$ Exponentialwerte $e^{t\lambda_j}$ werden skalarmäßig
    berechnet, und die Transformation kostet $\Theta(n^2)$.
    Gesamtaufwand: $\Theta(n^2)$ -- unabhängig von der Entwicklungstiefe $N$.
\end{beobachtung}

% ============================================================
\section{Beispiel: Merge-Sort}
% ============================================================

\begin{definition}[Sortierproblem und Vergleichsstruktur]
    Sei $L = (a_1, \ldots, a_n)$ eine Liste von Elementen einer
    total geordneten Menge $(M, \leq)$. Das \textbf{Sortierproblem}
    ist die Aufgabe, eine Permutation $\pi \in S_n$ zu finden mit
    $a_{\pi(1)} \leq a_{\pi(2)} \leq \cdots \leq a_{\pi(n)}$.

    Die \textbf{Verknüpfung} dieser Struktur ist der
    \textbf{Vergleich} $\leq$: Je zwei Elemente können verglichen werden.
    Die \textbf{Kosten} sind die Anzahl der Vergleiche.
\end{definition}

\begin{beobachtung}[Warum Sortieren eine Schwerpunktsuche ist]
    Eine ungeordnete Liste hat keinen Schwerpunkt bezüglich des
    Vergleichs: Jedes Element kann an jeder Position stehen,
    und die Ordnung der Elemente ist global -- sie hängt von
    allen Paaren ab. Die sortierte Liste ist der Schwerpunkt:
    Jedes Element ist an seiner ``natürlichen'' Position,
    und Ordnungsfragen sind lokal beantwortbar (Binärsuche in $O(\log n)$).

    Merge-Sort sucht den Schwerpunkt durch \textbf{Divide and Conquer}:
    Teile die Liste rekursiv in zwei Hälften (Zerlegung), sortiere
    jede Hälfte unabhängig (lokale Lösung im Schwerpunkt),
    füge die sortierten Hälften zusammen (Rücktransformation).
\end{beobachtung}

Die Algorithmen~\ref{alg:merge} und \ref{alg:mergesort} geben Merge und Merge-Sort in Pseudocode wieder. Algorithmus~\ref{alg:binsearch} zeigt die Binärsuche, die nach der Schwerpunkttransformation (Sortierung) möglich wird.


\begin{algorithm}[H]
\DontPrintSemicolon
\caption{Merge}\label{alg:merge}
\KwIn{Sortierte Listen $L_1 = (b_1 \leq \cdots \leq b_p)$, $L_2 = (c_1 \leq \cdots \leq c_q)$}
\KwOut{Sortierte Liste $L$ der Länge $p + q$}
$i \leftarrow 1;\quad j \leftarrow 1;\quad L \leftarrow ()$\;
\While{$i \leq p$ \textbf{und} $j \leq q$}{
    \eIf{$b_i \leq c_j$}{
        Hänge $b_i$ an $L$ an;\quad $i \leftarrow i + 1$\;
    }{
        Hänge $c_j$ an $L$ an;\quad $j \leftarrow j + 1$\;
    }
}
Hänge verbleibende Elemente von $L_1$ bzw.\ $L_2$ an $L$ an\;
\Return{$L$}\;
\BlankLine
\textbf{Vergleiche:} $\leq p + q - 1$ \quad \textbf{Laufzeit:} $M(n) = \Theta(n)$ mit $n = p + q$
\end{algorithm}

\begin{algorithm}[H]
\DontPrintSemicolon
\caption{Merge-Sort}\label{alg:mergesort}
\KwIn{Liste $L = (a_1, \ldots, a_n)$}
\KwOut{Sortierte Liste $L'$}
\If{$n \leq 1$}{
    \Return{$L$}\;
}
$m \leftarrow \lfloor n/2 \rfloor$\;
$L_1 \leftarrow \mathrm{MergeSort}(a_1, \ldots, a_m)$ \tcp*{linke Hälfte rekursiv sortieren}
$L_2 \leftarrow \mathrm{MergeSort}(a_{m+1}, \ldots, a_n)$ \tcp*{rechte Hälfte rekursiv sortieren}
\Return{$\mathrm{Merge}(L_1, L_2)$} \tcp*{Zusammenführen: $\Theta(n)$ Vergleiche}
\BlankLine
\textbf{Rekurrenz:} $T(n) = 2T(n/2) + \Theta(n)$\quad $\Rightarrow$\quad $T(n) = \Theta(n \log n)$
\end{algorithm}

\begin{algorithm}[H]
\DontPrintSemicolon
\caption{Binärsuche}\label{alg:binsearch}
\KwIn{Sortierte Liste $L = (a_1 \leq \cdots \leq a_n)$, Suchwert $x$}
\KwOut{Index $k$ mit $a_k = x$, oder \textbf{nicht gefunden}}
$\ell \leftarrow 1;\quad r \leftarrow n$\;
\While{$\ell \leq r$}{
    $m \leftarrow \lfloor(\ell + r)/2\rfloor$\;
    \uIf{$a_m = x$}{\Return{$m$}\;}
    \uElseIf{$a_m < x$}{$\ell \leftarrow m + 1$\;}
    \Else{$r \leftarrow m - 1$\;}
}
\Return{\textbf{nicht gefunden}}\;
\BlankLine
\textbf{Laufzeit:} $\Theta(\log n)$ Vergleiche \quad (vs.\ $\Theta(n)$ ohne Schwerpunkt)
\end{algorithm}

\begin{satz}[Merge-Sort: formale Laufzeitanalyse]
    \label{satz:mergesort}
    Der \textbf{Merge-Sort-Algorithmus} auf einer Liste der Länge $n$
    hat Zeitkomplexität
    \[
        T(n) = \Theta(n \log n).
    \]
    Dies ist \textbf{asymptotisch optimal} unter allen
    vergleichsbasierten Sortieralgorithmen.
\end{satz}

\begin{proof}
    \textbf{Teil 1: Obere Schranke $T(n) = O(n \log n)$.}

    Der Merge-Sort-Algorithmus lässt sich durch folgende Rekurrenz
    beschreiben:
    \[
        T(n) = 2 T\!\left(\tfrac{n}{2}\right) + M(n),
        \qquad T(1) = 0,
    \]
    wobei $M(n)$ die Kosten des \textbf{Merge-Schritts}
    (Zusammenführen zweier sortierter Hälften der Länge $n/2$)
    bezeichnet.

    \textbf{Analyse des Merge-Schritts:}
    Seien $L_1 = (b_1 \leq b_2 \leq \cdots \leq b_{n/2})$ und
    $L_2 = (c_1 \leq c_2 \leq \cdots \leq c_{n/2})$ zwei sortierte
    Listen. Der Merge-Algorithmus hält zwei Zeiger $i, j$ auf $L_1$
    bzw. $L_2$ und wählt bei jedem Schritt das kleinere Element:

    \begin{itemize}
        \item[-] Vergleiche $b_i$ mit $c_j$.
        \item[-] Füge das kleinere in die Ausgabeliste ein und
              rücke den entsprechenden Zeiger vor.
        \item[-] Wiederhole, bis einer der Zeiger das Ende erreicht.
    \end{itemize}

    Jeder Schritt fügt genau ein Element in die Ausgabeliste ein
    und benötigt genau einen Vergleich. Die Ausgabeliste hat Länge $n$,
    also ist $M(n) \leq n - 1 < n$ Vergleiche.
    (Im günstigsten Fall: $n/2$ Vergleiche, wenn alle Elemente von
    $L_1$ kleiner als alle von $L_2$ sind.)
    Also $M(n) = \Theta(n)$.

    \textbf{Auflösung der Rekurrenz} (Master-Theorem / Rekursionsbaum):
    Die Rekurrenz $T(n) = 2T(n/2) + \Theta(n)$ hat die Lösung
    $T(n) = \Theta(n \log n)$.

    \emph{Beweis durch Rekursionsbaum:}
    Der Rekursionsbaum hat Tiefe $\log_2 n$. Auf Ebene $k$
    ($0 \leq k \leq \log_2 n$) gibt es $2^k$ Teilprobleme der
    Größe $n/2^k$. Der Merge-Aufwand auf Ebene $k$ beträgt
    $2^k \cdot \Theta(n/2^k) = \Theta(n)$ pro Ebene.
    Ãœber alle $\log_2 n + 1$ Ebenen ergibt sich:
    \[
        T(n) = \sum_{k=0}^{\log_2 n} \Theta(n)
        = (\log_2 n + 1) \cdot \Theta(n)
        = \Theta(n \log n).
    \]

    \emph{Formale Induktion:} Behaupte $T(n) \leq c \cdot n \log_2 n$
    für eine Konstante $c > 0$ und alle $n \geq 2$.

    \emph{Induktionsanfang:} $T(2) = 2T(1) + M(2) = 0 + 1 = 1
    \leq c \cdot 2 \log_2 2 = 2c$ für $c \geq 1/2$.

    \emph{Induktionsschritt:} Sei die Behauptung für alle $m < n$ wahr.
    Dann:
    \begin{align*}
        T(n) &= 2T(n/2) + M(n) \\
             &\leq 2 \cdot c \cdot \tfrac{n}{2} \log_2 \tfrac{n}{2} + n \\
             &= c \cdot n (\log_2 n - 1) + n \\
             &= c \cdot n \log_2 n - cn + n \\
             &= c \cdot n \log_2 n - (c-1)n \\
             &\leq c \cdot n \log_2 n \qquad \text{für } c \geq 1.
    \end{align*}
    Damit gilt $T(n) \leq c \cdot n \log_2 n = O(n \log n)$.

    \textbf{Teil 2: Untere Schranke $T(n) = \Omega(n \log n)$
    für vergleichsbasierte Sortierverfahren.}

    \emph{Entscheidungsbaum-Argument:}
    Jeder vergleichsbasierte Sortieralgorithmus kann als
    \textbf{binärer Entscheidungsbaum} modelliert werden:
    Jeder innere Knoten ist ein Vergleich $a_i \leq a_j$ (ja/nein),
    jedes Blatt ist eine vollständig bestimmte Permutation der Eingabe.

    \begin{itemize}
        \item[-] Die Anzahl der möglichen Ausgaben (Permutationen
              von $n$ Elementen) ist $n!$.
        \item[-] Ein binärer Baum der Tiefe $d$ hat höchstens $2^d$ Blätter.
        \item[-] Um alle $n!$ Permutationen unterscheiden zu können,
              muss gelten: $2^d \geq n!$, also $d \geq \log_2(n!)$.
    \end{itemize}

    Nach der \textbf{Stirlingschen Formel}:
    $\log_2(n!) = \sum_{k=1}^n \log_2 k \geq \int_1^n \log_2 x\, dx
    = \frac{n \log_2 n - n}{\ln 2} = \Omega(n \log n)$.

    Genauer: $\log_2(n!) \geq \log_2\bigl(\frac{n}{2}\bigr)^{n/2}
    = \frac{n}{2} \log_2 \frac{n}{2} = \frac{n}{2}(\log_2 n - 1)
    = \Omega(n \log n)$.

    Also muss jeder vergleichsbasierte Sortieralgorithmus im
    schlechtesten Fall mindestens $\Omega(n \log n)$ Vergleiche
    durchführen.

    \textbf{Zusammenfassung:} Merge-Sort erreicht $T(n) = \Theta(n \log n)$
    und ist damit asymptotisch optimal unter allen
    vergleichsbasierten Sortierverfahren.
\end{proof}

\begin{beobachtung}[Merge-Sort als Schwerpunkt-Transformation]
    Merge-Sort übersetzt das Schwerpunkt-Prinzip direkt in einen
    Algorithmus: Eine unsortierte Liste der Länge $n$ hat
    \emph{keinen Schwerpunkt} -- der minimale Vergleichsaufwand für
    eine Suchanfrage ist $\Theta(n)$ (lineare Suche).

    Nach der Sortierung (Schwerpunkttransformation, Kosten $\Theta(n \log n)$)
    hat die Liste ihren Schwerpunkt: Eine Suchanfrage kostet nur noch
    $\Theta(\log n)$ (Binärsuche). Wird die Liste $m$-mal abgefragt, so ist
    \begin{align*}
        T_{\mathrm{unsortiert}}(n, m) &= m \cdot \Theta(n) = \Theta(mn), \\
        T_{\mathrm{sortiert}}(n, m)   &= \Theta(n \log n) + m \cdot \Theta(\log n).
    \end{align*}
    Der Break-even-Punkt liegt bei
    $\Theta(n \log n) + m \Theta(\log n) = \Theta(mn)$,
    also $m = \Theta(n)$: Für $m \geq c \cdot n$ lohnt sich die
    Schwerpunkttransformation.
\end{beobachtung}

\begin{beobachtung}[Effizienzgewinne im Vergleich]
    Beide Beispiele bestätigen das Schwer\-punkt-Prinzip quanti\-tativ:

    \begin{center}
    \begin{tabular}{lll}
        \hline
        Struktur & Direkte Kosten & Im Schwerpunkt \\
        \hline
        Vektoraddition ($n$-dim.) & $\Theta(n)$ & $\Theta(n)$ (bereits Schwerpunkt) \\
        Matrix-Vektor-Produkt ($k$-fach) & $\Theta(kn^2)$ & $\Theta(n^2 + kn)$ \\
        Matrixpotenz $A^k v$ & $\Theta(kn^2)$ & $\Theta(n^2)$ unabh.\ von $k$ \\
        Matrixexponential $e^{tA}$ & $\Theta(Nn^3)$ & $\Theta(n^2)$ \\
        Sortieren (einmalig) & $\Theta(n \log n)$ & $\Theta(n \log n)$ \\
        Suche (nach Sortierung) & $\Theta(n)$ & $\Theta(\log n)$ \\
        $m$ Suchen (nach Sortierung) & $\Theta(mn)$ & $\Theta(n\log n + m\log n)$ \\
        \hline
    \end{tabular}
    \end{center}

    In allen Fällen gilt: Der Aufwand der Schwerpunkttransformation
    amortisiert sich ab einer kritischen Wiederholungszahl.
    Der fundamentale Gewinn liegt in der \textbf{Reduktion globaler
    auf lokale Operationen} -- das Kernprinzip des Schwerpunkts.
\end{beobachtung}

% ============================================================


\subsection*{Schnelle Fourier-Transformation als Schwerpunkt-Algorithmus}

\begin{algorithm}[H]
\DontPrintSemicolon
\caption{Schnelle Fourier-Transformation (Cooley-Tukey, $n = 2^k$)}\label{alg:fft}
\KwIn{Vektor $a = (a_0, \ldots, a_{n-1}) \in \mathbb{C}^n$, $n = 2^k$}
\KwOut{DFT $\hat{a}$ mit $\hat{a}_j = \sum_{l=0}^{n-1} a_l \omega^{jl}$, $\omega = e^{2\pi i/n}$}
\If{$n = 1$}{\Return{$a$}\;}
$a^{[0]} \leftarrow (a_0, a_2, \ldots, a_{n-2})$ \tcp*{gerade Indizes}
$a^{[1]} \leftarrow (a_1, a_3, \ldots, a_{n-1})$ \tcp*{ungerade Indizes}
$\hat{a}^{[0]} \leftarrow \mathrm{FFT}(a^{[0]})$\;
$\hat{a}^{[1]} \leftarrow \mathrm{FFT}(a^{[1]})$\;
\For{$j \leftarrow 0$ \KwTo $n/2 - 1$}{
    $\hat{a}_j       \leftarrow \hat{a}^{[0]}_j + \omega^j \hat{a}^{[1]}_j$\;
    $\hat{a}_{j+n/2} \leftarrow \hat{a}^{[0]}_j - \omega^j \hat{a}^{[1]}_j$\;
}
\Return{$\hat{a}$}\;
\BlankLine
\textbf{Rekurrenz:} $T(n) = 2T(n/2) + \Theta(n)$\quad $\Rightarrow$\quad $T(n) = \Theta(n\log n)$\;
\textbf{Vgl.\ naive DFT:} $\Theta(n^2)$ \quad \textbf{Faktor:} $n / \log n$ schneller
\end{algorithm}

\begin{beobachtung}[FFT als Schwerpunkt-Transformation für Faltung]
    Die DFT transformiert das Faltungsprodukt $(a * b)_k = \sum_j a_j b_{k-j}$
    in ein komponentenweises Produkt: $\widehat{a * b} = \hat{a} \cdot \hat{b}$.
    Dies ist genau das Schwerpunkt-Prinzip: Im Frequenzbereich ist die
    Faltung (aufwändig, $\Theta(n^2)$) zu einer komponentenweisen
    Multiplikation (trivial, $\Theta(n)$) geworden.

    Die gesamte Faltungsberechnung via FFT:
    $\Theta(n\log n)$ für $\hat{a}$ und $\hat{b}$, $\Theta(n)$ für $\hat{a}\cdot\hat{b}$,
    $\Theta(n\log n)$ für die inverse FFT. Gesamt: $\Theta(n\log n)$ statt $\Theta(n^2)$.
\end{beobachtung}

\section{Schwerpunkt in weiteren Strukturen}
% ============================================================

\begin{beobachtung}[Schwerpunkt in der Fourier-Analysis]
    Die \textbf{Diskrete Fourier-Trans\-formation} (DFT) ist die
    Schwerpunkttransformation für den Vektorraum $\mathbb{C}^n$ mit
    der Faltungsoperation $(a * b)_k = \sum_j a_j b_{k-j}$.
    Faltung ist \emph{nicht} komponentenweise -- analog zur Matrixmultiplikation.

    Die DFT $\hat{a} = F_n a$ mit $F_n = (e^{2\pi i jk/n})_{j,k=0}^{n-1}$
    transformiert in den \textbf{Frequenzbereich}: Dort ist die Faltung
    komponentenweise: $\widehat{a * b} = \hat{a} \cdot \hat{b}$
    (punktweises Produkt). Dies ist der Schwerpunkt der Faltungsstruktur.

    Die \textbf{Schnelle Fourier-Transformation} (FFT, Cooley--Tukey, 1965)
    berechnet $F_n a$ in $\Theta(n \log n)$ Operationen statt $\Theta(n^2)$,
    indem sie die Divide-and-Conquer-Struktur des Schwerpunkts nutzt --
    exakt wie Merge-Sort, aber für die Faltungsstruktur.
\end{beobachtung}

\begin{beobachtung}[Schwerpunkt und der Hauptsatz über Moduln]
    Der Hauptsatz über endlich erzeugte Moduln über Hauptidealringen
    (Kapitel 19) ist der allgemeinste Ausdruck des Schwerpunkt-Prinzips
    in der Algebra: Jeder endlich erzeugte Modul $M$ über einem
    Hauptidealring $R$ lässt sich in seinen Schwerpunkt transformieren:
    \[
        M \cong R^k \oplus \bigoplus_j R/Rp_j^{e_j}.
    \]
    Diese Zerlegung ist der Schwerpunkt von $M$ bezüglich der
    $R$-Modulstruktur: Addition und skalare Multiplikation wirken
    komponentenweise auf den Direktsummanden.

    Die Jordansche Normalform (Spezialfall $R = K[x]$) und die
    Klassifikation endlicher abelscher Gruppen (Spezialfall $R = \mathbb{Z}$)
    sind beides Instanzen desselben Schwerpunkt-Prinzips in einer
    einzigen algebraischen Aussage.
\end{beobachtung}

\chapter{Endliche Körper, Rekursionen und Faktorräume}
% ============================================================

\section{Endliche Körper}

Lange waren endliche Körper fast ausschließlich ein Hilfsmittel der
abstrakten Algebra. Ihre handfeste Relevanz tritt erst durch die
Codierungstheorie zutage -- die Lehre der Sicherung von Daten gegen
zufällige Störungen. Diese Brücke motiviert, die Theorie endlicher
Körper hier knapp aber vollständig aufzubauen.

\begin{satz}[Charakteristik und Ordnung]
    Sei $K$ ein endlicher Körper mit $|K| = q$ und $\mathrm{Char}\, K = p$
    (einer Primzahl). Dann gilt $p \mid q$. Für jedes $a \in K$ ist ferner
    $a^q = a$. Insbesondere gilt $a^{q-1} = 1$ für alle $a \neq 0$.
\end{satz}
\begin{proof}
    Die multiplikative Gruppe $K^* = K \setminus \{0\}$ hat Ordnung $q - 1$.
    Nach dem Satz von Lagrange gilt $a^{q-1} = 1$ für alle $a \in K^*$,
    also $a^q = a$. Für $a = 0$ ist $0^q = 0$ trivialerweise. Da
    $\mathrm{Char}\, K = p$ und $K$ ein $\mathbb{F}_p$-Vektorraum der
    Dimension $n$ ist, folgt $q = p^n$; insbesondere $p \mid q$.
\end{proof}


\begin{beobachtung}[Warum \texorpdfstring{$a^q = a$}{a^q = a} gilt]
    Der Satz folgt aus dem Satz von Lagrange angewandt auf die multiplikative
    Gruppe $K^* = K \setminus \{0\}$: Diese Gruppe hat Ordnung $q - 1$,
    also gilt $a^{q-1} = 1$ für $a \neq 0$ und damit $a^q = a$.
    Für $a = 0$ ist die Aussage trivial.

    Dies ist die direkte Verallgemeinerung des kleinen Fermatschen Satzes
    $a^p \equiv a \pmod{p}$ auf beliebige endliche Körper -- die Gruppenstruktur
    macht den Satz zu einem Spezialfall einer universellen Aussage.
\end{beobachtung}

\begin{satz}[Frobenius-Automorphismus]
    Sei $K$ ein endlicher Körper der Charakteristik $p$. Die Abbildung
    $\alpha : K \to K$, $\alpha(k) = k^p$, ist ein \textbf{Automorphismus}
    von $K$ -- der \textbf{Frobenius-Automorphismus}.
\end{satz}
\begin{proof}
    Die Multiplikativität $\alpha(k_1 k_2) = (k_1 k_2)^p = k_1^p k_2^p
    = \alpha(k_1)\alpha(k_2)$ ist klar.
    Die Additivität folgt aus dem Binomischen Lehrsatz in Charakteristik $p$:
    \[
        (k_1 + k_2)^p = \sum_{j=0}^p \binom{p}{j} k_1^j k_2^{p-j}
        = k_1^p + k_2^p,
    \]
    da $p \mid \binom{p}{j}$ für $0 < j < p$.
    $\alpha$ ist injektiv, denn $\ker \alpha = 0$ (kein Nichtnull-Element
    hat $k^p = 0$ in einem Körper). Da $K$ endlich ist, ist $\alpha$ daher
    bijektiv -- also ein Automorphismus.
\end{proof}


\begin{beobachtung}[Die Kraft des Frobenius-Automorphismus]
    Der Frobenius ist keine Kuriosität, sondern ein Strukturprinzip:
    Er ist der einzige nicht-triviale Automorphismus, der sich aus der
    Charakteristik allein ergibt, und er erzeugt die gesamte Galois-Gruppe
    des Körpers. Die Eigenschaft $\alpha(k_1 + k_2) = (k_1 + k_2)^p =
    k_1^p + k_2^p$ -- die scheinbar nichttriviale Binomialformel -- gilt
    genau wegen $\mathrm{Char}\, K = p$: Alle gemischten Terme
    $\binom{p}{j}$ mit $0 < j < p$ sind durch $p$ teilbar und fallen heraus.

    Dieses Phänomen, dass Charakteristik $p$ die Potenzierung zu einem
    Ringhomomorphismus macht, ist der algebraische Kern vieler tiefer
    Sätze über endliche Körper.
\end{beobachtung}

\begin{satz}[Konstruktion von Erweiterungskörpern]
    Sei $K$ ein Körper und $f(x) = x^2 + ax + b$ ein Polynom über $K$
    ohne Nullstellen in $K$. Auf $L = K \times K$ wird durch
    \[
        (x_1, y_1)(x_2, y_2) = (x_1 x_2 - b y_1 y_2,\; x_1 y_2 + x_2 y_1 - a y_1 y_2)
    \]
    ein Körper $L$ mit $K$-isomorphem Unterkörper $\{(x, 0) \mid x \in K\}$
    definiert. Ist $|K| = q$, so hat $L$ die Ordnung $q^2$.
\end{satz}
\begin{proof}
    Es ist zu zeigen, dass $L = K \times K$ mit der angegebenen Multiplikation
    ein Körper ist. Die Distributiv- und Assoziativgesetze rechnet man nach.
    Das neutrale Element der Multiplikation ist $(1, 0)$.
    Das Inverse von $(x, y) \neq (0, 0)$: Da $f(z) = z^2 + az + b$ keine
    Nullstelle in $K$ hat, ist $N(x,y) = x^2 - bxy - ay^2 \neq 0$ für alle
    $(x,y) \neq (0,0)$ (denn $N(x,y) = 0$ würde eine Nullstelle von $f$ liefern).
    Also ist $(x,y)^{-1} = N(x,y)^{-1}(x - ay, -y)$.
    Die Einbettung $k \mapsto (k,0)$ ist ein Körperhomomorphismus.
    Da $|L| = |K|^2$, gilt die Ordnungsaussage.
\end{proof}


\begin{beobachtung}[Die Struktur endlicher Körper]
    Die vollständige Wahrheit über endliche Körper, die hier ohne Beweis
    festgehalten sei:
    \begin{itemize}
        \item[-] Jeder endliche Körper der Charakteristik $p$ hat $p^n$ Elemente
              für ein $n \in \mathbb{N}$.
        \item[-] Zu jeder Primzahlpotenz $p^n$ gibt es bis auf Isomorphie genau
              einen Körper $\mathbb{F}_{p^n}$ mit $p^n$ Elementen.
        \item[-] Die multiplikative Gruppe $\mathbb{F}_{p^n}^*$ ist zyklisch der
              Ordnung $p^n - 1$.
    \end{itemize}
    Dies bedeutet insbesondere: Die endlichen Körper sind vollständig
    klassifiziert -- eine seltene Situation in der Algebra, wo vollständige
    Klassifikationen die Ausnahme sind.
\end{beobachtung}

% ============================================================
\section{Rekursionsgleichungen}
% ============================================================

Zahlreiche Fragen aus Mathematik, Informatik und Ingenieurwissenschaften
führen auf lineare Rekursionsgleichungen. Die Vektorraumtheorie liefert
ein vollständiges und elegantes Lösungsverfahren.

\begin{definition}[Rekursionsgleichung]
    Sei $K$ ein Körper, $F = \{(x_j) \mid x_j \in K\}$ der $K$-Vektorraum
    aller Folgen über $K$, und seien $a_0, \dots, a_{n-1} \in K$ gegeben.
    Eine \textbf{Rekursionsgleichung} der Ordnung $n$ ist:
    \[
        (R) \quad x_{k+n} = a_0 x_k + a_1 x_{k+1} + \cdots + a_{n-1} x_{k+n-1}
        \qquad (k \geq 0).
    \]
    Das zugehörige \textbf{charakteristische Polynom} ist
    $f(x) = x^n - a_0 - a_1 x - \cdots - a_{n-1} x^{n-1}$.
\end{definition}

\begin{satz}[Lösungsraum einer Rekursionsgleichung]
    Die Lösungen von $(R)$ bilden einen Unterraum $L \leq F$ der Dimension $n$.
    Zu gegebenen Anfangswerten $(x_0, \dots, x_{n-1})$ gibt es genau eine
    Lösung.

    Sind $b_1, \dots, b_m \in K$ paarweise verschieden mit $f(b_i) = 0$,
    so sind die Folgen $(b_i^j)_{j \geq 0}$ linear unabhängige Elemente von $L$.
    Hat $f$ genau $n$ verschiedene Nullstellen, so bilden diese Folgen eine
    Träger -- jede Lösung hat die Form
    \[
        x_j = \sum_{i=1}^n c_i b_i^j \qquad (c_i \in K).
    \]
    Ist $b$ eine mehrfache Nullstelle mit $f(b) = f'(b) = 0$, so ist
    auch $(j \cdot b^{j-1})_{j \geq 0}$ eine von $(b^j)$ linear unabhängige
    Lösung.
\end{satz}
\begin{proof}
    \textbf{Unterraum:} Sind $(x_j)$ und $(y_j)$ Lösungen von $(R)$, so auch
    jede Linearkombination -- die Rekursion ist linear. Damit ist $L$ ein Unterraum.
    \textbf{Dimension $n$:} Die Abbildung $L \to K^n$, $(x_j) \mapsto (x_0, \ldots, x_{n-1})$,
    ist linear und injektiv (die Anfangswerte bestimmen die Folge eindeutig per Rekursion).
    Sie ist surjektiv, da man zu jedem Anfangswert-Tupel eine eindeutige Lösung konstruiert.
    \textbf{Nullstellen erzeugen Lösungen:} Ist $f(b) = 0$, so verifiziert man direkt,
    dass $(b^j)$ die Rekursion erfüllt: $b^{j+n} = b^j \cdot b^n = b^j(a_0 + \cdots
    + a_{n-1} b^{n-1})$. Verschiedene Nullstellen $b_1, \ldots, b_m$ liefern linear
    unabhängige Lösungen (Vandermonde-Determinante $\det(b_i^j) \neq 0$).
\end{proof}


\begin{beispiel}[Fibonacci-Zahlen und Markov-Grenzwert]
    \begin{itemize}
        \item[-] Die älteste Rekursionsgleichung $x_{j+2} = x_j + x_{j+1}$
              hat das charakteristische Polynom $f(x) = x^2 - x - 1$ mit
              den beiden Nullstellen $\frac{1 \pm \sqrt{5}}{2}$. Die allgemeine
              Lösung ist
              \[
                  x_j = c_1 \left(\frac{1+\sqrt{5}}{2}\right)^j
                  + c_2 \left(\frac{1-\sqrt{5}}{2}\right)^j.
              \]
              Die Fibonacci-Zahlen $(1, 1, 2, 3, 5, 8, 13, \dots)$ sind
              der Spezialfall $x_0 = x_1 = 1$; sie wachsen asymptotisch
              wie $\left(\frac{1+\sqrt{5}}{2}\right)^j$.
        \item[-] Für $x_{j+2} = \frac{1}{2}(x_j + x_{j+1})$ erhält man
              $\lim_{j\to\infty} x_j = \frac{1}{3}(x_0 + 2x_1)$ -- der
              Grenzwert ist eine gewichtete Linearkombination der Anfangswerte.
    \end{itemize}
\end{beispiel}

\begin{beobachtung}[Rekursionsgleichungen als Kryptographie-Baustein]
    Über dem Körper $K = \mathbb{F}_2 = \{0, 1\}$ dienen Rekursionsgleichungen
    als Grundlage von \textbf{Stromchiffren}: Alice verschlüsselt eine
    Nachricht $x = (x_j) \in F$ durch Addition mit einer Schlüsselfolge
    $y = (y_j)$, die als Lösung einer Rekursionsgleichung erzeugt wird.
    Bob, der die Rekursionskoeffizienten und die Anfangswerte kennt,
    entschlüsselt durch $x = (x + y) + y$.

    Die Sicherheit hängt von der Periode der Folge ab -- bei einem
    zyklischen Generator $b \in \mathbb{F}_{p^n}^*$ ist die maximale
    Periode $p^n - 1$. Hier schließt sich die Theorie endlicher Körper
    nahtlos an die Vektorraum-Theorie der Rekursionen an.
\end{beobachtung}

% ============================================================
\section{Der Faktorraum}
% ============================================================

\begin{definition}[Faktorraum]
    Sei $V$ ein $K$-Vektorraum und $U \leq V$. Die Menge der
    \textbf{Nebenklassen}
    \[
        V/U = \{v + U \mid v \in V\}, \qquad v + U = \{v + u \mid u \in U\},
    \]
    wird durch
    \[
        (v_1 + U) + (v_2 + U) = v_1 + v_2 + U, \qquad k(v + U) = kv + U
    \]
    zu einem $K$-Vektorraum, dem \textbf{Faktorraum} $V/U$. Das neutrale
    Element ist $U = 0 + U$.

    Für endlichdimensionales $V$ gilt: $\dim V/U = \dim V - \dim U$.
\end{definition}

\begin{satz}[Dimensionsformel für Faktorräume]
    Seien $U \leq W \leq V$ Unterräume mit $\dim V/U < \infty$. Dann gilt:
    \[
        \dim V/U = \dim V/W + \dim W/U.
    \]
    Sind $[w_i + U]$ einen Träger von $W/U$ und $[v_j + W]$ einen Träger
    von $V/W$, so bilden $[w_i + U,\, v_j + U]$ zusammen einen Träger
    von $V/U$.
\end{satz}
\begin{proof}
    Sei $[w_i + U]_{i=1}^r$ einen Träger von $W/U$ und $[v_j + W]_{j=1}^s$
    einen Träger von $V/W$. Wir zeigen, dass $[w_1 + U, \ldots, w_r + U,
    v_1 + U, \ldots, v_s + U]$ einen Träger von $V/U$ ist.
    \textbf{Erzeugung:} Sei $v + U \in V/U$. Dann $v + W = \sum_j b_j(v_j + W)$,
    also $v - \sum_j b_j v_j \in W$. Dieser Vektor ist Linearkombination
    der $w_i$ modulo $U$, was die Erzeugungseigenschaft liefert.
    \textbf{Lineare Unabhängigkeit:} Sei $\sum_i a_i(w_i + U) + \sum_j b_j(v_j + U) = U$.
    Dann liegt $\sum_j b_j v_j \in W$, also $\sum_j b_j(v_j + W) = W$, und
    da $[v_j + W]$ linear unabhängig ist, folgt $b_j = 0$. Dann liegt
    $\sum_i a_i w_i \in U$, also $a_i = 0$.
    Somit $\dim V/U = r + s = \dim W/U + \dim V/W$.
\end{proof}


\begin{beobachtung}[Was der Faktorraum leistet]
    Der Faktorraum ist zunächst eine formale Konstruktion, erweist sich
    aber als universelles Werkzeug. Er beschreibt präzise, was in einer
    linearen Abbildung \emph{vergessen} wird: Ist $f : V \to W$ linear,
    so ist $V/\ker f \cong \mathrm{Im}\, f$ -- der Homomorphiesatz.

    Konzeptuell: Der Faktorraum $V/U$ zeigt dieselbe Struktur wie $V$,
    aber mit $U$ auf Null kollabiert. Er ist der \emph{Blickwinkel}, in dem
    $U$ transparent ist -- und aus dem sich $V$ modulo der Unterraumstruktur
    von $U$ beschreiben lässt.

    Ein konkretes Beispiel: Die Restklassenringe $\mathbb{Z}/n\mathbb{Z}$
    der Zahlentheorie sind nichts anderes als Faktorräume über $\mathbb{Z}$
    als $\mathbb{Z}$-Modul -- dasselbe Konstruktionsprinzip in einem
    allgemeineren Kontext.
\end{beobachtung}

\begin{definition}[Hyperebene]
    Ein Unterraum $U \leq V$ heißt \textbf{Hyperebene}, falls
    $\dim V/U = 1$ gilt, d.\,h. $U$ hat Kodimension $1$ in $V$.
\end{definition}

\begin{satz}[Hyperebenen als Kerne von Funktionalen]
    Sei $V$ ein $K$-Vektorraum. Dann gilt:
    \begin{itemize}
        \item[-] Ist $0 \neq f \in \mathrm{Hom}_K(V, K)$, so ist $\ker f$
              eine Hyperebene in $V$.
        \item[-] Ist $U \leq V$ eine Hyperebene, so gibt es ein
              $0 \neq f \in \mathrm{Hom}_K(V, K)$ mit $\ker f = U$.
        \item[-] Jeder Unterraum $U$ der Kodimension $k$ ist der
              Durchschnitt von $k$ Hyperebenen.
    \end{itemize}
\end{satz}
\begin{proof}
    \textbf{(a)} Ist $f \neq 0$, so $\mathrm{Im}\, f = K$, also $\dim \mathrm{Im}\, f = 1$.
    Nach dem Rangsatz gilt $\dim V/\ker f = \dim \mathrm{Im}\, f = 1$.
    \textbf{(b)} Ist $\dim V/U = 1$, wähle $w \notin U$; dann $V = U \oplus \langle w \rangle$.
    Definiere $f(u + aw) = a$. Dies ist linear mit $\ker f = U$.
    \textbf{(c)} Nach der Dimensionsformel für Faktorräume ist jeder Unterraum der
    Kodimension $k$ Durchschnitt von $k$ Unterräumen der Kodimension $1$,
    und für jeden davon liefert (b) ein entsprechendes Funktional.
\end{proof}


\begin{beobachtung}[Faktorräume und der Homomorphiesatz]
    Der Faktorraum trägt den Schlüssel zum Homomorphiesatz, dem
    fundamentalen Struktursatz linearer Abbildungen: Jede lineare
    Abbildung $f : V \to W$ faktorisiert als
    \[
        V \xrightarrow{B} V/\ker f \xrightarrow{C} W,
    \]
    wobei $B$ der kanonische Epimorphismus und $C$ ein Monomorphismus ist
    mit $\mathrm{Im}\, C = \mathrm{Im}\, f$. Insbesondere:
    $\mathrm{Im}\, f \cong V / \ker f$.

    Dies ist der vektorräumliche Ausdruck des Rangsatzes: Was verloren geht
    (Kern), und was übrig bleibt (Bild), addieren sich zur Gesamtdimension.
    Der Faktorraum macht diesen Sachverhalt zur exakten Gleichung
    $\dim V = \dim \ker f + \dim \mathrm{Im}\, f$.
\end{beobachtung}

% ============================================================
\chapter{Spur, Projektionen und Gleichungssysteme}
% ============================================================

\section{Das Rechnen mit linearen Abbildungen}

\begin{definition}[Vektorraumstruktur auf \texorpdfstring{$\mathrm{Hom}_K(V,W)$}{Hom_K(V,W)}]
    Die Menge $\mathrm{Hom}_K(V, W)$ aller linearen Abbildungen von $V$
    nach $W$ ist selbst ein $K$-Vektorraum: Für $A, B \in \mathrm{Hom}_K(V,W)$
    und $k \in K$ setzt man
    \[
        (A + B)v = Av + Bv, \qquad (kA)v = k(Av).
    \]
    Ist $\dim V = m$ und $\dim W = n$, so gilt
    $\dim \mathrm{Hom}_K(V, W) = mn$.

    Für Endomorphismen $A, B \in \mathrm{End}_K(V)$ ist das Produkt
    $AB$ durch Komposition definiert: $(AB)v = A(Bv)$. Damit wird
    $\mathrm{End}_K(V)$ zu einem Ring.
\end{definition}

\begin{beobachtung}[Warum \texorpdfstring{$\dim \mathrm{Hom}_K(V,W) = mn$}{ Hom_K(V,W) = mn}]
    Jede lineare Abbildung ist durch ihre Werte auf eines Trägers vollständig
    und frei bestimmt: Wählt man einen Träger $[v_1, \dots, v_m]$ von $V$
    und einen Träger $[w_1, \dots, w_n]$ von $W$, so entspricht jedes
    $A \in \mathrm{Hom}_K(V,W)$ eindeutig einer $(n \times m)$-Matrix --
    einem Element von $K^{n \times m}$. Die Dimension $mn$ ist daher
    keine Ãœberraschung, sondern die unvermeidliche Konsequenz der Freiheit
    linearer Abbildungen auf Träger.
\end{beobachtung}

\begin{satz}[Homomorphiesatz]
    Sei $A \in \mathrm{Hom}_K(V, W)$. Es gibt einen Epimorphismus
    $B : V \twoheadrightarrow V/\ker A$ und einen Monomorphismus
    $C : V/\ker A \hookrightarrow W$ mit $A = CB$ und
    $\mathrm{Im}\, C = \mathrm{Im}\, A$. Insbesondere ist
    $\mathrm{Im}\, A \cong V / \ker A$.

    Für endlichdimensionales $V$ folgt daraus:
    \[
        \dim \ker A + \dim \mathrm{Im}\, A = \dim V.
    \]
\end{satz}
\begin{proof}
    Definiere $C : V/\ker A \to W$ durch $C(v + \ker A) = Av$.
    Dies ist wohldefiniert: Ist $v + \ker A = v' + \ker A$, so $v - v' \in \ker A$,
    also $A(v - v') = 0$, also $Av = Av'$.
    $C$ ist linear und injektiv: Aus $C(v + \ker A) = 0$ folgt $Av = 0$,
    also $v \in \ker A$, also $v + \ker A = \ker A$. Das Bild von $C$
    ist $\mathrm{Im}\, A$. Der kanonische Epimorphismus $B : v \mapsto v + \ker A$
    erfüllt $A = CB$ per Konstruktion.
    Der Rangsatz folgt: $\dim V = \dim \ker A + \dim(V/\ker A) = \dim \ker A + \dim \mathrm{Im}\, A$.
\end{proof}


% ============================================================
\section{Die Spur}
% ============================================================

\begin{definition}[Spur]
    Sei $A = (a_{ij}) \in K^{n \times n}$. Die \textbf{Spur} von $A$ ist
    die Summe der Diagonaleinträge:
    \[
        \mathrm{tr}\, A = \sum_{i=1}^n a_{ii}.
    \]
\end{definition}

\begin{satz}[Eigenschaften der Spur]
    Seien $A, B \in K^{n \times n}$ und $k \in K$. Dann gilt:
    \begin{itemize}
        \item[-] $\mathrm{tr}(A + B) = \mathrm{tr}\, A + \mathrm{tr}\, B$
              und $\mathrm{tr}(kA) = k \cdot \mathrm{tr}\, A$
              \quad (Linearität)
        \item[-] $\mathrm{tr}(AB) = \mathrm{tr}(BA)$
              \quad (Zyklizität)
        \item[-] $\mathrm{tr}\, A^T = \mathrm{tr}\, A$
        \item[-] Für ähnliche Matrizen $A' = P^{-1}AP$ gilt
              $\mathrm{tr}\, A' = \mathrm{tr}\, A$.
    \end{itemize}
\end{satz}
\begin{proof}
    \textbf{Linearität} ist komponentenweise klar.
    \textbf{Zyklizität:} $(AB)_{ii} = \sum_k a_{ik} b_{ki}$, also
    $\mathrm{tr}(AB) = \sum_i \sum_k a_{ik} b_{ki} = \sum_k \sum_i b_{ki} a_{ik}
    = \mathrm{tr}(BA)$.
    \textbf{Transponierte:} $(A^t)_{ii} = a_{ii}$, also $\mathrm{tr}\, A^t = \mathrm{tr}\, A$.
    \textbf{Ähnlichkeit:} $\mathrm{tr}(P^{-1}AP) = \mathrm{tr}(AP \cdot P^{-1})
    = \mathrm{tr}(A)$ nach der Zyklizität.
\end{proof}


\begin{beobachtung}[Spur als Invariante]
    Die Zyklizitätseigenschaft $\mathrm{tr}(AB) = \mathrm{tr}(BA)$ ist
    der Kern der Bedeutung der Spur: Sie zeigt, dass die Spur
    \emph{trägerunabhängig} ist -- invariant unter Ähnlichkeitstransformationen.
    Damit ist die Spur eine wohlbestimmte Eigenschaft des Endomorphismus
    $f$ selbst, nicht nur seiner Darstellungsmatrix.

    Zusammen mit der Determinante bildet die Spur das kompakteste
    Kennzeichen einer linearen Abbildung: Für eine $(2 \times 2)$-Matrix ist
    das charakteristische Polynom vollständig durch $\mathrm{tr}\, A$ und
    $\det A$ bestimmt:
    $\chi_A(\lambda) = \lambda^2 - (\mathrm{tr}\, A)\lambda + \det A$.

    In der Physik: Die Spur ist die Summe der Eigenwerte -- in der
    Quantenmechanik die Gesamtwahrscheinlichkeit (Spur des Dichteoperators
    ist stets $1$), in der Relativitätstheorie die Viererspur des
    Energie-Impuls-Tensors.
\end{beobachtung}

% ============================================================
\section{Projektionen und direkte Zerlegungen}
% ============================================================

\begin{definition}[Projektion]
    Ein Endomorphismus $P \in \mathrm{End}_K(V)$ heißt
    \textbf{Projektion}, falls $P^2 = P$.
\end{definition}

\begin{satz}[Direkte Zerlegung durch Projektionen]
    Sei $P \in \mathrm{End}_K(V)$ eine Projektion. Dann gilt:
    \[
        V = \mathrm{Im}\, P \oplus \ker P.
    \]
    Umgekehrt definiert jede direkte Zerlegung $V = U_1 \oplus U_2$ eine
    Projektion $P$ mit $\mathrm{Im}\, P = U_1$ und $\ker P = U_2$ durch
    $P(u_1 + u_2) = u_1$.

    Die Abbildung $Q = \mathrm{id} - P$ ist die komplementäre Projektion
    auf $U_2$ längs $U_1$, und es gilt $PQ = QP = 0$.
\end{satz}
\begin{proof}
    Sei $P^2 = P$. Für $v \in V$ gilt $v = Pv + (v - Pv)$, wobei $Pv \in \mathrm{Im}\, P$
    und $P(v - Pv) = Pv - P^2v = Pv - Pv = 0$, also $v - Pv \in \ker P$.
    Damit $V = \mathrm{Im}\, P + \ker P$.
    Die Summe ist direkt: Ist $Pv = w \in \mathrm{Im}\, P \cap \ker P$,
    so $w = Pw' $ für ein $w'$ und $Pw = P^2w' = Pw' = w$, aber auch $Pw = 0$,
    also $w = 0$.
    Umgekehrt: Für $V = U_1 \oplus U_2$ definiere $P(u_1 + u_2) = u_1$.
    Dann $P^2(u_1 + u_2) = P(u_1) = u_1 = P(u_1 + u_2)$, also $P^2 = P$.
    $Q = \mathrm{id} - P$ erfüllt $Q^2 = Q$ und $PQ = P(\mathrm{id} - P)
    = P - P^2 = 0$.
\end{proof}


\begin{beobachtung}[Die direkte Summe als Schwerpunkt-Zerlegung]
    Eine Projektion $P$ mit $P^2 = P$ gibt dem Vektorraum exakt die
    Struktur, die das Schwerpunkt-Prinzip fordert: $V$ zerfällt in zwei
    unabhängige Hälften $U_1 = \mathrm{Im}\, P$ und $U_2 = \ker P$,
    auf denen $P$ mit maximaler Lokalität wirkt -- als Identität auf
    $U_1$ und als Null auf $U_2$.

    In Signaltheorie und Regelungstechnik entspricht dies der Zerlegung
    eines Zustandsraums in erregbaren und nicht-erregbaren Anteil.
    Im algebraischen Kontext: Direkte Zerlegungen sind die Vektorraum-Analog
    der disjunkten Zerlegungen von Mengen -- vollständig, ohne Überlappung,
    ohne Rest.
\end{beobachtung}

\begin{beobachtung}[Projektionen und Idempotente]
    Die Gleichung $P^2 = P$ ist die algebraische Definition von
    \textbf{Idempotenz}. Sie sagt: Einmal projizieren reicht --
    ein zweites Anwenden ändert nichts mehr.

    Dies ist genau die Eigenschaft, die eine Projektion von einer
    allgemeinen Abbildung unterscheidet: Jeder Vektor in $\mathrm{Im}\, P$
    hat seinen Schwerpunkt bereits erreicht. Die Idempotenz ist die
    algebraische Formulierung des Satzes: \emph{Projektion auf einen
    Unterraum ist Fixpunkt genau auf diesem Unterraum}.
\end{beobachtung}

% ============================================================
\section{Elementare Umformungen und lineare Gleichungs\-systeme}
% ============================================================

\begin{definition}[Elementare Zeilenumformungen]
    Sei $A \in K^{m \times n}$. Die folgenden Operationen auf den Zeilen
    von $A$ heißen \textbf{elementare Zeilenumformungen}:
    \begin{itemize}
        \item[-] \textbf{Zeilenvertauschung}: Vertausche Zeile $i$ und Zeile $j$.
        \item[-] \textbf{Skalierung}: Multipliziere Zeile $i$ mit $k \neq 0$.
        \item[-] \textbf{Addition}: Addiere das $k$-fache von Zeile $j$ zu Zeile $i$ ($i \neq j$).
    \end{itemize}
    Jede elementare Umformung entspricht der Linksmultiplikation mit einer
    invertierbaren \textbf{Elementarmatrix}. Sie verändern den Rang von $A$
    nicht.
\end{definition}

\begin{satz}[Gaußscher Algorithmus und Treppenform]
    Jede Matrix $A \in K^{m \times n}$ lässt sich durch elementare
    Zeilenumformungen in \textbf{Zeilenstufenform} (Treppenform) bringen:
    eine Matrix, in der die führende Eins jeder Zeile rechts von der
    führenden Eins der darüber liegenden Zeile steht, und alle Einträge
    unterhalb einer führenden Eins gleich Null sind.

    Der \textbf{Rang} von $A$ ist die Anzahl der Nicht-Nullzeilen in der
    Treppenform -- invariant unter elementaren Umformungen.
\end{satz}
\begin{proof}
    Wir zeigen durch Induktion, dass jede Matrix durch elementare Zeilenumformungen
    in Treppenform gebracht werden kann.
    Ist $A = 0$, fertig. Andernfalls sei $a_{ij}$ der erste Nichtnulleintrag (Pivotspalte $j$).
    Durch Zeilenvertauschung bringe man diesen in Zeile $1$; durch Skalierung
    mache man $a_{1j} = 1$; durch Subtraktion passender Vielfache von Zeile $1$
    setze man alle anderen Einträge in Spalte $j$ auf $0$. Das Ergebnis hat
    Treppengestalt in den ersten Einträgen; Induktion auf die Untermatrix ab
    Zeile $2$ liefert den Rest.
    Der Rang als Anzahl der Nichtnullzeilen ist invariant, da elementare
    Umformungen den Zeilenraum erhalten.
\end{proof}


\begin{algorithm}[H]
\DontPrintSemicolon
\caption{Gaußscher Algorithmus}\label{alg:gauss}
\KwIn{Erweiterte Matrix $(A \mid b) \in K^{m \times (n+1)}$}
\KwOut{Zeilenstufenform $(A' \mid b')$, $\mathrm{Rang}\,A$}
$r \leftarrow 0;\quad j \leftarrow 1$\;
\While{$r < m$ \textbf{und} $j \leq n$}{
    Finde $i \geq r+1$ mit $a_{ij} \neq 0$\;
    \eIf{\emph{kein solches $i$}}{
        $j \leftarrow j + 1$\;
    }{
        Tausche Zeile $r+1 \leftrightarrow$ Zeile $i$\;
        $r \leftarrow r + 1$\;
        \For{$l \leftarrow r+1$ \KwTo $m$}{
            $\lambda \leftarrow a_{lj} / a_{rj}$\;
            Zeile $l \;\leftarrow\;$ Zeile $l - \lambda \cdot$ Zeile $r$\;
        }
        $j \leftarrow j + 1$\;
    }
}
\Return{$(A' \mid b')$,\; $r = \mathrm{Rang}\,A$}\;
\BlankLine
\textbf{Aufwand:} $\Theta(n^3)$ Operationen (quadratisches $A$)
\end{algorithm}


\begin{satz}[Lineare Gleichungssysteme]
    Sei $Ax = b$ ein lineares Gleichungssystem mit $A \in K^{m \times n}$
    und $b \in K^m$. Dann gilt:
    \begin{itemize}
        \item[-] Das System ist \textbf{lösbar} genau dann, wenn
              $\mathrm{Rang}\, A = \mathrm{Rang}\,(A \mid b)$.
        \item[-] Die Lösungsmenge ist ein \textbf{affiner Unterraum}: eine
              Nebenklasse $x_0 + \ker A$ einer partikulären Lösung $x_0$.
        \item[-] Der Lösungsraum hat Dimension $n - \mathrm{Rang}\, A$.
    \end{itemize}
\end{satz}
\begin{proof}
    Das System $Ax = b$ ist lösbar genau dann, wenn $b \in \mathrm{Im}\, A$,
    d.\,h.\ wenn $b$ eine Linearkombination der Spalten von $A$ ist.
    Dies ist äquivalent zu $\mathrm{Rang}(A) = \mathrm{Rang}(A \mid b)$,
    da das Hinzufügen von $b$ als Spalte den Rang genau dann erhöht,
    wenn $b \notin \mathrm{Im}\, A$.
    Ist $x_0$ eine partikuläre Lösung, so ist die allgemeine Lösung
    $x_0 + \ker A$, denn $Ax = b \Leftrightarrow A(x - x_0) = 0
    \Leftrightarrow x - x_0 \in \ker A$.
    Die Dimension von $\ker A$ beträgt $n - \mathrm{Rang}\, A$ nach dem Rangsatz.
\end{proof}


\begin{beobachtung}[Lineare Gleichungssysteme als Kern-Bild-Zerlegung]
    Das lineare Gleichungssystem $Ax = b$ ist nichts anderes als die
    Frage: Liegt $b$ im Bild von $A$? Die Antwort -- ja oder nein --
    entscheidet sich am Rang.

    Die Struktur der Lösungsmenge spiegelt den Rangsatz: Die
    Dimension des Lösungsraums $n - \mathrm{Rang}\, A$ ist genau
    $\dim \ker A$. Jede Lösung ist eine partikuläre Lösung plus
    eine beliebige homogene Lösung -- die affine Struktur entsteht
    durch Verschiebung des Kerns in die Richtung von $b$.

    Der Gaußsche Algorithmus macht dieses abstrakte Bild algorithmisch
    greifbar: Er transformiert $A$ durch Trägerwechsel in eine Form, in der
    Kern und Bild direkt ablesbar sind. Die Treppenform ist die
    rechnerische Verwirklichung des Rangsatzes.
\end{beobachtung}

% ============================================================
% ============================================================
\chapter{Gruppenhomomorphismen, Permutationen und Determinanten}
% ============================================================

\section{Gruppenhomomorphismen und Normalteiler}

Der Ãœbergang von Gruppen zu linearen Abbildungen ist kein Zufall:
Beide beschreiben strukturerhaltende Abbildungen zwischen algebraischen
Objekten. Die Gruppentheorie verallgemeinert dieses Prinzip auf
beliebige algebraische Strukturen.

\begin{definition}[Gruppenhomomorphismus und Normalteiler]
    Seien $G$ und $H$ Gruppen. Eine Abbildung $\alpha : G \to H$ heißt
    \textbf{Homomorphismus}, falls $\alpha(g_1 g_2) = \alpha(g_1)\alpha(g_2)$
    für alle $g_1, g_2 \in G$.

    Eine Untergruppe $N \leq G$ heißt \textbf{Normalteiler} ($N \trianglelefteq G$),
    falls $g^{-1}ng \in N$ für alle $n \in N$, $g \in G$. Der \textbf{Kern} eines
    Homomorphismus ist stets ein Normalteiler; sein \textbf{Bild} ist eine
    Untergruppe von $H$.
\end{definition}

\begin{satz}[Faktorgruppe und Homomorphiesatz]
    Sei $N \trianglelefteq G$. Die Menge $G/N = \{gN \mid g \in G\}$ wird
    durch $g_1 N \cdot g_2 N = g_1 g_2 N$ zu einer Gruppe -- der
    \textbf{Faktorgruppe}. Die kanonische Abbildung $\tau : G \to G/N$,
    $\tau g = gN$, ist ein Epimorphismus mit $\ker \tau = N$.

    Für jeden Homomorphismus $\alpha : G \to H$ gilt:
    \[
        G / \ker \alpha \;\cong\; \mathrm{Bild}\, \alpha.
    \]
\end{satz}
\begin{proof}
    \textbf{Wohldefiniertheit der Gruppenoperation:} Sei $g_1 N = g_1' N$ und
    $g_2 N = g_2' N$, also $g_1' = g_1 n_1$ und $g_2' = g_2 n_2$ mit $n_i \in N$.
    Dann $g_1' g_2' = g_1 n_1 g_2 n_2 = g_1 g_2 (g_2^{-1} n_1 g_2) n_2$.
    Da $N \trianglelefteq G$, ist $g_2^{-1} n_1 g_2 \in N$, also
    $(g_1' g_2') N = (g_1 g_2) N$.
    \textbf{Homomorphiesatz:} Die Abbildung $\bar{\alpha} : G/\ker\alpha \to H$,
    $\bar{\alpha}(g \ker\alpha) = \alpha(g)$, ist wohldefiniert und injektiv
    (da $\alpha(g) = e \Leftrightarrow g \in \ker\alpha$) mit Bild $= \mathrm{Bild}\,\alpha$.
\end{proof}

\begin{beobachtung}[Normalteiler als Maß der Symmetrie]
    Nicht jede Untergruppe ist ein Normalteiler. Die Bedingung
    $g^{-1}Ng = N$ für alle $g \in G$ ist eine Symmetriebedingung:
    $N$ verhält sich gleichartig gegenüber allen Konjugationen.
    In kommutativen Gruppen sind alle Untergruppen automatisch Normalteiler.

    Der Homomorphiesatz ist das gruppentheoretische Gegenstück des
    Rangsatzes für Vektorräume: Der Kern misst, was verloren geht;
    das Bild, was erhalten bleibt. Die Faktorgruppe $G/\ker\alpha$
    beschreibt exakt den Teil von $G$, der in $H$ injiziert wird.
\end{beobachtung}

\begin{definition}[Einfache Gruppen und Kompositionsreihen]
    Eine Gruppe $G \neq \{e\}$ heißt \textbf{einfach}, falls $G$ keine
    echten Normalteiler außer $\{e\}$ und $G$ selbst hat.

    Eine \textbf{Kompositionsreihe} von $G$ ist eine absteigende Kette
    $G = G_0 \supsetneq G_1 \supsetneq \cdots \supsetneq G_r = \{e\}$,
    in der jedes $G_{j+1}$ ein Normalteiler von $G_j$ ist und
    $G_j/G_{j+1}$ einfach ist. Die Quotienten $G_j/G_{j+1}$ heißen
    \textbf{Kompositionsfaktoren}.
\end{definition}

\begin{satz}[Jordan--Hölder]
    Je zwei Kompositionsreihen einer endlichen Gruppe $G$ haben dieselbe
    Länge $r$ und bis auf Reihenfolge dieselben Kompositionsfaktoren.
\end{satz}
\begin{proof}
    Durch Induktion nach $|G|$. Seien $G_\bullet$ und $H_\bullet$ zwei
    Kompositionsreihen. Ist $G_1 = H_1$, so wendet man die Induktionsannahme
    auf $G_1$ an. Andernfalls gilt $G_1 H_1 = G$ (da $G_1 H_1 \trianglelefteq G$
    und $G_1 \subsetneq G_1 H_1$, also $G_1 H_1 = G$ wegen Einfachheit von $G/G_1$).
    Die Kompositionsfaktoren von $G_1 \cap H_1$ zählt man in beiden
    Reihen auf dieselbe Weise -- durch den Isomorphiesatz
    $G_1/(G_1 \cap H_1) \cong G/H_1$ und $H_1/(G_1 \cap H_1) \cong G/G_1$.
\end{proof}

\begin{beobachtung}[Jordan--Hölder und die Klassifikation einfacher Gruppen]
    Der Satz von Jordan--Hölder ist das gruppentheoretische Analogon
    der eindeutigen Primfaktorzerlegung: Einfache Gruppen sind die
    ``Primfaktoren'' endlicher Gruppen.

    Das $20.$ Jahrhundert vollbrachte die vollständige Klassifikation
    aller endlichen einfachen Gruppen -- eine der größten kollektiven
    mathematischen Leistungen (über 10\,000 Seiten Beweis, 100+ Autoren,
    abgeschlossen $\approx 2004$). Die Klassen sind:
    \begin{itemize}
        \item[-] Zyklische Gruppen $\mathbb{Z}_p$ (Primzahl $p$)
        \item[-] Alternierende Gruppen $A_n$ ($n \geq 5$)
        \item[-] 16 Familien von Gruppen vom Lie-Typ
        \item[-] 26 sporadische Gruppen (darunter $M_{24}$, die Automorphismengruppe
              des Golay-Codes)
    \end{itemize}
\end{beobachtung}

% ============================================================
\section{Permutationen und das Signum}
% ============================================================

\begin{definition}[Symmetrische Gruppe und Zyklen]
    Die Menge $S_n$ aller bijektiven Abbildungen von $\{1, \dots, n\}$
    auf sich bildet unter Komposition die \textbf{symmetrische Gruppe}
    mit $|S_n| = n!$. Ihre Elemente heißen \textbf{Permutationen}.

    Ein \textbf{$k$-Zykel} $(a_1, a_2, \dots, a_k)$ bildet $a_i \mapsto a_{i+1}$
    (Index modulo $k$) und fixiert alle übrigen Elemente.
    \textbf{Transpositionen} sind 2-Zyklen $(i,j)$.

    Jede Permutation besitzt eine eindeutige \textbf{Zyklenzerlegung}
    in disjunkte Zyklen, und jede Permutation ist ein Produkt von
    Transpositionen.
\end{definition}

\begin{satz}[Signum]
    Für $n > 1$ gibt es einen eindeutigen Epimorphismus
    $\mathrm{sgn} : S_n \to \{1, -1\}$ mit $\mathrm{sgn}(\tau) = -1$
    für alle Transpositionen $\tau$. Explizit:
    \[
        \mathrm{sgn}(\pi) = \prod_{\{i,j\},\, i < j}
        \mathrm{sgn}(\pi j - \pi i),
    \]
    wobei das Produkt über alle Paare $i < j$ läuft.
    Eine Permutation heißt \textbf{gerade} (bzw. \textbf{ungerade}),
    falls $\mathrm{sgn}(\pi) = 1$ (bzw. $-1$).
    Der Kern $A_n = \ker(\mathrm{sgn})$ heißt \textbf{alternierende Gruppe}.
\end{satz}
\begin{proof}
    \textbf{Wohldefiniertheit:} Jede Permutation ist ein Produkt von Transpositionen.
    Wir zeigen, dass die Parität der Anzahl der Transpositionen in jeder Zerlegung
    dieselbe ist. Dies folgt daraus, dass das Polynom
    $\Delta(x_1,\ldots,x_n) = \prod_{i<j}(x_i - x_j)$ unter einer Transposition
    $(k,l)$ das Vorzeichen wechselt (da genau ein Faktor $x_k - x_l$ sein Vorzeichen
    ändert und alle anderen Faktoren höchstens permutiert werden).
    Daher ist $\mathrm{sgn}(\pi) = \pm 1$ je nachdem, ob $\pi(\Delta) = +\Delta$
    oder $\pi(\Delta) = -\Delta$.
    \textbf{Homomorphismus:} $\pi(\sigma(\Delta)) = (\pi\sigma)(\Delta)$, also
    $\mathrm{sgn}(\pi)\mathrm{sgn}(\sigma) = \mathrm{sgn}(\pi\sigma)$.
\end{proof}

\begin{beobachtung}[Signum und Determinante]
    Das Signum ist die Brücke zur Determinante: Die Leibniz-Formel
    \[
        \det A = \sum_{\pi \in S_n} \mathrm{sgn}(\pi) \prod_{i=1}^n a_{i,\pi(i)}
    \]
    ist genau das gewichtete Signum über alle Permutationen.
    Das Vorzeichen $\mathrm{sgn}(\pi)$ entscheidet, ob der Beitrag
    der Permutation $\pi$ addiert oder subtrahiert wird.

    Die Tiefe des Signum-Begriffs zeigt sich darin, dass jede
    Abbildung $f : S_n \to K^*$ entweder konstant $1$ ist oder
    mit dem Signum übereinstimmt (sofern $\mathrm{Char}\, K \neq 2$) --
    es gibt also genau \emph{einen} nichttrivialen Gruppencharakter von $S_n$.
\end{beobachtung}

\begin{beobachtung}[Die Ordnungen von $A_n$ und Einfachheit]
    Da $A_n = \ker(\mathrm{sgn})$ der Index-2-Normalteiler von $S_n$ ist,
    gilt $|A_n| = n!/2$. Für $n \geq 5$ ist $A_n$ einfach -- dies
    ist der gruppentheoretische Kern der Unlösbarkeit der allgemeinen
    Gleichung $n$-ten Grades durch Radikale für $n \geq 5$: Die
    Galoisgruppe der allgemeinen Quintic enthält $A_5$, und weil
    $A_5$ einfach und nichtabelsch ist, gibt es keine
    Kompositionsreihe mit abelschen Faktoren, die eine Radikallösung
    ermöglichen würde.
\end{beobachtung}

\chapter{Polynomringe, Ideale und die Jordansche Normalform}
% ============================================================

\section{Polynomringe}

\begin{definition}[Polynomring]
    Sei $R$ ein Ring. Der \textbf{Polynomring} $R[x]$ ist die Menge
    aller Folgen $(a_0, a_1, \dots)$ mit $a_j \in R$ und nur endlich
    vielen $a_j \neq 0$, versehen mit Addition und Multiplikation:
    \[
        (a_j) + (b_j) = (a_j + b_j), \qquad
        (a_j)(b_j) = (c_k) \text{ mit } c_k = \sum_{j=0}^k a_j b_{k-j}.
    \]
    Für einen Körper $K$ ist $K[x]$ eine kommutative $K$-Algebra
    mit $K$-Träger $[1, x, x^2, \dots]$, wobei $x = (0,1,0,\dots)$.
    Der \textbf{Grad} $\deg f$ eines Polynoms $f \neq 0$ ist der
    Index des höchsten Koeffizienten $\neq 0$.
\end{definition}

\begin{satz}[Division mit Rest]
    Seien $f, g \in K[x]$ mit $g \neq 0$. Dann gibt es eindeutige
    $q, r \in K[x]$ mit
    \[
        f = q \cdot g + r \qquad \text{und} \qquad \deg r < \deg g.
    \]
    Insbesondere: $a \in K$ ist Nullstelle von $f$ genau dann, wenn
    $(x - a) \mid f$.
\end{satz}
\begin{proof}
    Wir beweisen Existenz und Eindeutigkeit.
    \textbf{Existenz:} Sei $f \neq 0$. Unter allen Darstellungen $f = qg + r$ mit
    $r = 0$ oder $\deg r < \deg g$ wähle diejenige mit kleinstmöglichem $\deg r$.
    Wäre $\deg r \geq \deg g$, so subtrahiere $(r_{\mathrm{lc}}/g_{\mathrm{lc}}) x^{\deg r - \deg g} \cdot g$
    von $r$ -- das ergibt einen Widerspruch zur Minimalität.
    \textbf{Eindeutigkeit:} Aus $qg + r = q'g + r'$ folgt $(q-q')g = r' - r$.
    Wäre $q \neq q'$, so $\deg(r' - r) = \deg((q-q')g) \geq \deg g$ -- Widerspruch
    zu $\deg r, \deg r' < \deg g$.
\end{proof}


\begin{beobachtung}[Warum der Polynomring so mächtig ist]
    Der Polynomring $K[x]$ ist ein \textbf{euklidischer Ring}: Der
    Divisionsalgorithmus mit Rest funktioniert genauso wie für ganze
    Zahlen. Daraus folgt: $K[x]$ ist ein Hauptidealring -- jedes
    Ideal hat die Form $f K[x]$ für ein einziges Polynom $f$.

    Diese Analogie zwischen $\mathbb{Z}$ und $K[x]$ ist kein Zufall,
    sondern Ausdruck einer gemeinsamen algebraischen Struktur. Der
    Begriff des \emph{größten gemeinsamen Teilers}, der Euklidische
    Algorithmus, die Primfaktorzerlegung -- all das gilt in $K[x]$
    genau wie in $\mathbb{Z}$.
\end{beobachtung}

% ============================================================
\section{Ringe, Ideale und der Chinesische Restsatz}
% ============================================================

\begin{definition}[Ideal und Faktorring]
    Sei $R$ ein Ring. Eine Teilmenge $A \subseteq R$ heißt \textbf{Ideal},
    falls $A$ eine Untergruppe der additiven Gruppe ist und
    $r_1 a r_2 \in A$ für alle $a \in A$, $r_1, r_2 \in R$.

    Die Menge $R/A = \{r + A \mid r \in R\}$ wird durch
    $(r_1 + A)(r_2 + A) = r_1 r_2 + A$ zum \textbf{Faktorring}.
    Ein Ringhomomorphismus $\alpha : R \to S$ hat Kern $\ker \alpha$
    als Ideal in $R$, und es gilt $R/\ker\alpha \cong \mathrm{Bild}\,\alpha$.
\end{definition}

\begin{satz}[Chinesischer Restsatz]
    Sei $R$ ein kommutativer Ring und $A_1, \dots, A_n$ Ideale mit
    $A_i + A_j = R$ für alle $i \neq j$ (paarweise koprim).
    Dann gilt für jede Wahl von $r_1, \dots, r_n \in R$:
    \[
        \exists\, r \in R : \quad r \equiv r_j \pmod{A_j}
        \quad (j = 1, \dots, n).
    \]
    Insbesondere: $R / (A_1 \cdots A_n) \cong R/A_1 \times \cdots \times R/A_n$.
\end{satz}
\begin{proof}
    \textbf{Lösbarkeit:} Es genügt, für jedes $j$ ein $e_j \in R$ zu konstruieren
    mit $e_j \equiv 1 \pmod{A_j}$ und $e_j \equiv 0 \pmod{A_i}$ für $i \neq j$.
    Da $A_j + A_i = R$ für $i \neq j$, gilt $\prod_{i \neq j} (A_j + A_i) \supseteq A_j^{n-1} + \ldots$,
    und mit paarweiser Koprimheit folgt $A_j + \prod_{i \neq j} A_i = R$,
    also $\prod_{i \neq j} A_i \not\subseteq A_j$: Es gibt $e_j \in \prod_{i \neq j} A_i$
    mit $e_j \equiv 1 \pmod{A_j}$. Die Lösung ist dann $r = \sum_j r_j e_j$.
    \textbf{Isomorphismus:} Die kanonische Abbildung $\varphi : R \to \prod_j R/A_j$
    ist nach obigem surjektiv mit Kern $\bigcap_j A_j = \prod_j A_j$.
\end{proof}


\begin{beobachtung}[Der Chinesische Restsatz als Zerlegungsprinzip]
    Der Chinesische Restsatz ist ein universelles Zerlegungsprinzip:
    Er sagt, dass paarweise koprime Bedingungen unabhängig voneinander
    lösbar sind -- und dass die gemeinsame Lösung stets existiert.
    Für $R = \mathbb{Z}$ ist dies der klassische Satz: Aus
    $x \equiv a_i \pmod{n_i}$ mit $\mathrm{ggT}(n_i, n_j) = 1$ folgt
    stets eine eindeutige Lösung modulo $n_1 \cdots n_n$.

    In der Linearen Algebra erscheint er als Zerlegung
    des Vektorraums in Primär\-kom\-po\-nen\-ten:
    des Vektorraums in Primär\-komponenten: Wenn das Minimalpolynom
    $\mu_A = \prod_i p_i^{a_i}$ in paarweise teilerfremde Faktoren
    zerfällt, dann zerfällt der Vektorraum
    in die direkte Summe der $p_i$-Primär\-räume.
\end{beobachtung}
% ============================================================
\section{Die Jordansche Normalform}
% ============================================================

\begin{satz}[Jordansche Normalform]
    Sei $K$ algebraisch abgeschlossen (z.\,B. $K = \mathbb{C}$),
    $V$ ein $n$-dimensionaler $K$-Vektorraum und $A \in \mathrm{End}(V)$.
    Dann gibt es einen Träger von $V$, bezüglich derer $A$ die
    \textbf{Jordansche Normalform}
    \[
        J = \mathrm{diag}(J_{m_1}(\lambda_1), \dots, J_{m_r}(\lambda_r))
    \]
    hat, mit den \textbf{Jordankästchen}
    \[
        J_m(\lambda) = \begin{pmatrix}
            \lambda & 1 & & \\
            & \lambda & \ddots & \\
            & & \ddots & 1 \\
            & & & \lambda
        \end{pmatrix} \in K^{m \times m}.
    \]
    Die Jordanform ist eindeutig bis auf Reihenfolge der Kästchen.
\end{satz}
\begin{proof}[Beweisskizze (Kapitel 9)]
    Über $\mathbb{C}$ zerfällt $\chi_A$ vollständig. Durch den Chinesischen
    Restsatz zerfällt $V$ als $K[x]$-Modul in Primärräume
    $V = \bigoplus_j V_{\lambda_j}$ mit $V_{\lambda_j} = \ker(A-\lambda_j E)^{m_j}$.
    Auf jedem $V_{\lambda_j}$ ist $N = A - \lambda_j E$ nilpotent.
    Für eine nilpotente Abbildung $N$ mit $N^m = 0$, $N^{m-1} \neq 0$
    wähle $v$ mit $N^{m-1}v \neq 0$ und bilde die Kette
    $[v, Nv, \ldots, N^{m-1}v]$; diese ist linear unabhängig.
    Per Induktion zerlegt sich $V_{\lambda_j}$ vollständig in solche Ketten,
    was der Jordanform entspricht. Eindeutigkeit folgt aus Kapitel 11 (Rangsatz für Potenzen).
\end{proof}


\begin{beobachtung}[Bedeutung der Jordankästchen]
    Ein Jordankästchen $J_m(\lambda)$ beschreibt einen Endomorphismus,
    der $m$-dimensional ist und der Gleichung $(A - \lambda E)^m = 0$
    genügt, aber nicht $(A - \lambda E)^{m-1} = 0$ -- $\lambda$ ist
    ein Eigenwert der \emph{algebraischen} Vielfachheit $m$, aber
    der \emph{geometrischen} Vielfachheit $1$.

    Der Teil $\lambda E$ ist die diagonale Streckung; der Teil
    $N = J_m(\lambda) - \lambda E$ ist nilpotent: $N^m = 0$, $N^{m-1} \neq 0$.
    Jedes Jordankästchen ist die Summe aus einem Skalar und einem
    nilpotenten Anteil -- die Jordan-Zerlegung $A = S + N$
    in den \emph{halbeinfachen} Teil $S$ (diagonalisierbar) und den
    \emph{nilpotenten} Teil $N$.
\end{beobachtung}

\begin{beobachtung}[Jordanform und Matrixexponential]
    Die Jordanform klärt die Struktur des Matrixexponentials vollständig.
    Für ein Jordankästchen $J_m(\lambda)$ gilt:
    \[
        e^{J_m(\lambda)t} = e^{\lambda t}
        \begin{pmatrix}
            1 & t & \frac{t^2}{2!} & \cdots & \frac{t^{m-1}}{(m-1)!} \\
              & 1 & t & \cdots & \vdots \\
              & & \ddots & \ddots & \frac{t^2}{2!} \\
              & & & 1 & t \\
              & & & & 1
        \end{pmatrix}.
    \]
    Jede Lösung des Systems $\dot{x} = Ax$ ist also eine Linearkombination
    von Termen der Form $t^k e^{\lambda_j t}$ -- die Eigenwerte bestimmen
    die Exponenten, die Kästchengröße die polynomialen Vorfaktoren.
    Dies erklärt, warum mehrfache Eigenwerte in dynamischen Systemen
    zu polynomiellem Wachstum führen.
\end{beobachtung}

\begin{satz}[Eindeutigkeit der Jordanform]
    Die Kästchengrößen sind durch $A$ eindeutig bestimmt: Die Anzahl
    $N_\lambda(k)$ der Jordankästchen zu $\lambda$ der Größe $\geq k$ ist
    \[
        N_\lambda(k) = \mathrm{Rang}(A - \lambda E)^{k-1} -
        \mathrm{Rang}(A - \lambda E)^k.
    \]
    Die Jordanform ist daher aus den Rängen der Potenzen $(A - \lambda E)^k$
    vollständig und algorithmisch bestimmbar.
\end{satz}
\begin{proof}
    Zu zeigen ist, dass die Anzahl der Jordankästchen zu $\lambda$ der Größe $\geq k$
    durch $N_\lambda(k) = \mathrm{Rang}(A - \lambda E)^{k-1} - \mathrm{Rang}(A - \lambda E)^k$
    gegeben ist. Für ein einzelnes Jordankästchen $J_m(\lambda)$ gilt:
    $(J_m(\lambda) - \lambda E)^r = N^r$ mit nilpotentem $N$, und
    $\mathrm{Rang}(N^r) = \max(m - r, 0)$. Für die gesamte Jordanform addieren sich
    die Ränge über alle Kästchen, und die Differenz $\mathrm{Rang}(A-\lambda E)^{k-1}
    - \mathrm{Rang}(A-\lambda E)^k$ zählt genau die Kästchen der Größe $\geq k$.
    Da diese Größen aus $A$ berechnet werden können, ist die Jordanform eindeutig.
\end{proof}


% ============================================================

% ── Erweiterung ──────────────────────────────────────────────────────────
\section{Determinante und Orientierung in der Geometrie}

\begin{beobachtung}[Determinante als Orientierungsindikator]
    Sei $A \in \mathbb{R}^{n \times n}$. Das Vorzeichen von $\det A$
    kodiert die \textbf{Orientierung}: Ist $\det A > 0$, so erhält
    $A$ die Orientierung des Raums (wie die Identität); ist
    $\det A < 0$, so kehrt $A$ die Orientierung um (wie eine
    Spiegelung).

    In $\mathbb{R}^2$: $\det(v_1, v_2) > 0$ bedeutet, dass $(v_1, v_2)$
    ein \emph{rechtshändiges} System bildet (Linksdreh von $v_1$ nach $v_2$);
    $\det < 0$ bedeutet Linkshändigkeit.

    In $\mathbb{R}^3$: Die Determinante ist das \textbf{Spatprodukt}
    $\det(v_1, v_2, v_3) = v_1 \cdot (v_2 \times v_3)$, das vorzeichenbehaftete
    Volumen des von $v_1, v_2, v_3$ aufgespannten Parallelepipeds.
\end{beobachtung}

\begin{beobachtung}[Leibniz-Formel und Permutationsstruktur]
    Die Leibniz-Formel $\det A = \sum_{\pi \in S_n} \mathrm{sgn}(\pi)
    \prod_i a_{i,\pi(i)}$ hat eine schöne kombinatorische Deutung:
    Sie summiert über alle $n!$ Permutationen der Spaltenindizes.
    Jede Permutation wählt genau einen Eintrag aus jeder Zeile,
    das Signum entscheidet über das Vorzeichen des Beitrags.

    Für $n = 2$: $\det\bigl(\begin{smallmatrix}a & b \\ c & d\end{smallmatrix}\bigr)
    = ad - bc$. Die zwei Permutationen $(1\mapsto 1, 2\mapsto 2)$ und
    $(1\mapsto 2, 2\mapsto 1)$ mit Signum $+1$ und $-1$.

    Für $n = 3$: Sechs Permutationen, drei mit $\mathrm{sgn} = +1$
    (Sarrus-Regel: drei Hauptdiagonalprodukte) und drei mit $-1$
    (drei Nebendiagonalprodukte).
\end{beobachtung}

\begin{beobachtung}[Determinante und das charakteristische Polynom]
    Das charakteristische Polynom $\chi_A(\lambda) = \det(\lambda E - A)$
    ist eine Determinante der $\lambda$-parametrisierten Matrix
    $\lambda E - A$. Die Koeffizienten sind Polynome in den
    Einträgen von $A$:
    \begin{itemize}
        \item[-] Der führende Term ist $\lambda^n$.
        \item[-] Der Koeffizient von $\lambda^{n-1}$ ist $-\mathrm{tr}\,A$.
        \item[-] Der konstante Term ist $(-1)^n \det A$.
    \end{itemize}
    Dies zeigt: Die Spur ist der negative Koeffizient des zweithöchsten
    Terms, die Determinante (bis auf Vorzeichen) der konstante Term von
    $\chi_A$. Beide sind Spezialfälle der \textbf{Koeffizienten
    des charakteristischen Polynoms}, die ihrerseits symmetrische
    Polynome in den Eigenwerten sind.
\end{beobachtung}


\section{Entwicklungssatz, Rekursion und Rechenpraktik}

\begin{beobachtung}[Blockdeterminanten]
    Für Blockmatrizen $A = \bigl(\begin{smallmatrix} B & C \\ 0 & D \end{smallmatrix}\bigr)$
    mit quadratischen Blöcken $B \in K^{r \times r}$ und $D \in K^{s \times s}$ gilt:
    \[
        \det A = \det B \cdot \det D.
    \]
    Allgemeiner gilt für invertierbare $B$:
    $\det\bigl(\begin{smallmatrix}B & C \\ E & F\end{smallmatrix}\bigr)
    = \det B \cdot \det(F - EB^{-1}C)$, wobei $F - EB^{-1}C$ das
    \textbf{Schur-Komplement} von $B$ ist.

    Blockdeterminanten sind ein wichtiges Rechenwerkzeug: Statt eine
    große Determinante direkt zu berechnen, nutzt man die Struktur
    aus, um auf kleinere Determinanten zu reduzieren.
\end{beobachtung}

\begin{beobachtung}[Determinante und lineare Unabhängigkeit]
    Vektoren $v_1, \ldots, v_n \in K^n$ sind genau dann linear unabhängig,
    wenn $\det(v_1 \mid \cdots \mid v_n) \neq 0$. Dies ist eine
    direkte Konsequenz des Multiplikationssatzes: Die Vektoren bilden
    einen Träger von $K^n$ genau dann, wenn die Matrix $A = (v_1 \mid \cdots \mid v_n)$
    invertierbar ist, d.\,h.\ $\det A \neq 0$.

    In der Praxis: Ob $n$ gegebene Vektoren linear unabhängig sind,
    kann durch Gauß-Elimination in $\Theta(n^3)$ getestet werden --
    oder durch Berechnung der Determinante, die ebenfalls $\Theta(n^3)$
    kostet (über LU-Zerlegung, die den Gaußschen Algorithmus verwendet).
\end{beobachtung}


\section{Algorithmen zur Determinantenberechnung}

\begin{beobachtung}[LU-Zerlegung und Determinante]
    Die effizienteste Methode zur Determinantenberechnung ist die
    \textbf{LU-Zerlegung} via Gauß-Elimination: Man bringt $A$
    in die Form $PA = LU$ mit Permutationsmatrix $P$, unterer
    Dreiecksmatrix $L$ (mit Diagonaleinträgen $1$) und oberer
    Dreiecksmatrix $U$. Dann gilt:
    \[
        \det A = \mathrm{sgn}(P) \cdot \prod_{i=1}^n u_{ii},
    \]
    wobei $\mathrm{sgn}(P) = \pm 1$ das Signum der Permutation $P$ ist.
    Die Diagonaleinträge $u_{ii}$ sind die Pivot-Elemente der
    Gauß-Elimination. Der Gesamtaufwand beträgt $\Theta(n^3)$.
\end{beobachtung}

\begin{beobachtung}[Determinante und numerische Stabilität]
    Bei der numerischen Berechnung von Determinanten großer Matrizen
    treten zwei Probleme auf:
    \begin{itemize}
        \item[-] \textbf{Über\-lauf}: Für $n = 100$ und Einträge der
              Größen\-ordnung $10$, hat $\det A$ die Größen\-ordnung $10^{100}$
              -- jenseits der Darstellbarkeit
              in doppelter Genauigkeit ($\approx 10^{308}$ maximal).
              ($\approx 10^{308}$ maximal).
        \item[-] \textbf{Auslöschung}: Nahezu singuläre Matrizen
              (kleine singuläre Werte) führen zu massivem Rundungsfehler.
    \end{itemize}
    In der Praxis wird daher nicht $\det A$ direkt berechnet, sondern
    $\log|\det A| = \sum_i \log|u_{ii}|$ -- die logarithmische
    Determinante, die numerisch stabil ist und keine Überläufe erzeugt.
    Alternativ berechnet man $\det A$ via $\prod_j \sigma_j^2$
    (Produkt der quadrierten Singulärwerte, stets nichtnegativ).
\end{beobachtung}

\chapter{Skalarprodukte und Minkowskiraum}
% ============================================================

\section{Allgemeine Skalarprodukte}

\begin{definition}[Reguläres Skalarprodukt]
    Sei $V$ ein $K$-Vektorraum. Eine \textbf{Bilinearform} auf $V$ ist
    eine Abbildung $(\cdot, \cdot) : V \times V \to K$, linear in beiden
    Argumenten. Sie heißt \textbf{symmetrisch}, falls $(u,v) = (v,u)$,
    und \textbf{regulär}, falls aus $(v, w) = 0$ für alle $w$ stets
    $v = 0$ folgt.

    Der \textbf{Index} einer symmetrischen Bilinearform ist die maximale
    Dimension eines Unterraums, auf dem die Form negativ definit ist.
\end{definition}

\begin{beobachtung}[Signatur und Sylvesterscher Trägheitssatz]
    Für reelle symmetrische Bilinearformen besagt der \textbf{Sylvestersche
    Trägheitssatz}: In jeder Orthogonalträger hat die Form genau
    $p$ positive und $q$ negative Diagonaleinträge. Das Paar $(p,q)$
    heißt \textbf{Signatur} und ist eine Invariante der Form -- unabhängig
    von der Wahl der Orthogonalträger.

    Eine positiv definite Form hat Signatur $(n, 0)$; das ist das
    euklidische Skalarprodukt. Die Signatur $(3, 1)$ -- drei positive,
    ein negativer Eintrag -- ist die Signatur des Minkowskiraums,
    dem geometrischen Rahmen der speziellen Relativitätstheorie.
\end{beobachtung}

% ============================================================
\section{Der Minkowskiraum}
% ============================================================

\begin{definition}[Minkowskiraum und Lorentzgruppe]
    Sei $V = \mathbb{R}^4$ mit dem \textbf{Minkowski-Skalarprodukt}:
    \[
        ((x_j), (y_j)) = x_1 y_1 + x_2 y_2 + x_3 y_3 - c^2 x_4 y_4,
    \]
    wobei $c$ die Lichtgeschwindigkeit bezeichnet. Der Vektorraum $V$
    mit diesem Skalarprodukt heißt \textbf{Minkowskiraum}.

    Die Vektoren zerfallen in drei Klassen:
    \begin{itemize}
        \item[-] \textbf{Raumartig}: $(v, v) > 0$
        \item[-] \textbf{Zeitartig}: $(v, v) < 0$
        \item[-] \textbf{Lichtartig} (Lichtkegel $L$): $(v, v) = 0$
    \end{itemize}
    Die Isometrien des Minkowskiraums heißen \textbf{Lorentz-Transformationen};
    ihre Gruppe ist die \textbf{Lorentzgruppe} $\mathcal{L}$.
\end{definition}

\begin{satz}[Struktur der Lorentzgruppe]
    Die Lorentzgruppe $\mathcal{L}$ enthält als ausgezeichneten Normalteiler
    die \textbf{eigentliche Lorentzgruppe} $\mathcal{L}^+ = \{G \in \mathcal{L}
    \mid G \mathcal{Z}^+ = \mathcal{Z}^+\}$, die den Zukunftslichtkegel
    $\mathcal{Z}^+ = \{v \mid (v,v) < 0,\, v_4 > 0\}$ erhält.
    Mit $\mathcal{S} = \{G \in \mathcal{L} \mid \det G = 1\}$ gilt
    $|\mathcal{L}/\mathcal{S}^+| = 4$.
\end{satz}
\begin{proof}[Beweisskizze]
    Sei $G = \mathcal{L}$ die Lorentz\-gruppe. Die Determinante ist ein
    Homomorphismus $\det : G \to \{+1, -1\}$, ebenso das Vorzeichen
    der vierten Komponente (Zeit-Richtung). Die vier Zusammenhangskomponenten
    entstehen durch die vier Kombinationen
    $(\det, \mathrm{sgn}(G_{44})) \in \{+1,-1\}^2$.
    Die eigentliche Lorentzgruppe $\mathcal{L}^+ = \ker(\det) \cap \{G_{44} > 0\}$
    ist ein Normalteiler, da das Vorzeichen des Lichtkegels eine invariante
    Bedingung ist. Der Quotient $\mathcal{L}/\mathcal{S}^+$ hat genau
    4 Elemente.
\end{proof}


\begin{beobachtung}[Was der Minkowskiraum geometrisch bedeutet]
    Der Minkowskiraum ist kein euklidischer Raum -- das negative Vorzeichen
    vor dem Zeitterm erzeugt eine fundamental andere Geometrie:
    \begin{itemize}
        \item[-] Der Abstand zweier Ereignisse $\Delta s^2 = \Delta x^2 +
              \Delta y^2 + \Delta z^2 - c^2 \Delta t^2$ ist das
              \emph{invariante Raumzeitintervall} -- unabhängig vom
              Bezugssystem.
        \item[-] Der Lichtkegel $(v,v) = 0$ trennt kausal erreichbare
              Ereignisse (zeitartig, $(v,v) < 0$) von kausal nicht
              erreichbaren (raumartig, $(v,v) > 0$).
    \end{itemize}
    Das Vorzeichen $-c^2$ ist keine algebraische Willkür, sondern
    Ausdruck des physikalischen Prinzips: Kein Signal propagiert
    schneller als Licht. Die Geometrie des Minkowskiraums \emph{kodiert}
    die Kausalstruktur der Raumzeit.
\end{beobachtung}

% ============================================================
\section{Lorentz-Transformationen und Einsteins Relativitätstheorie}
% ============================================================

\begin{satz}[Lorentz-Translation und Einsteinsches Additionsgesetz]
    Für $|u| < c$ ist die \textbf{Lorentz-Translation} $L(v)$ mit
    $v = u e_3$ die Isometrie, die in $e_3$-Richtung den Boost beschreibt:
    \[
        L(v) e_3 = \frac{1}{d}\!\left(e_3 + \frac{u}{c^2} e_4\right),
        \qquad
        L(v) e_4 = \frac{1}{d}(u e_3 + e_4),
        \qquad d = \sqrt{1 - \tfrac{u^2}{c^2}} > 0.
    \]
    Für zwei kollineare Geschwindigkeiten gilt das
    \textbf{Einsteinsche Additionsgesetz}:
    \[
        L(u_1 e_3)\, L(u_2 e_3) = L\!\left(\frac{u_1 + u_2}{1 + u_1 u_2 / c^2}\, e_3\right).
    \]
\end{satz}
\begin{proof}
    Man verifiziert direkt, dass $L(v)$ die Minkowski-Form erhält:
    $(L(v)e_3, L(v)e_3) = \tfrac{1-u^2/c^2}{d^2} \cdot c^2
    = c^2 \cdot \tfrac{1-u^2/c^2}{1-u^2/c^2} = \ldots$
    Für das Additionsgesetz gilt $L(u_1 e_3) L(u_2 e_3)$: Man multipliziert
    die Matrizen aus und liest die resultierende Geschwindigkeit ab.
    Die Nichtüberschreitung der Lichtgeschwindigkeit: Für $|u_1|, |u_2| < c$ gilt
    $\left|\frac{u_1+u_2}{1+u_1u_2/c^2}\right| < c$, da
    $(c^2 - u_1^2)(c^2 - u_2^2) > 0$ impliziert $(c \pm (u_1+u_2))/(1+u_1u_2/c^2)
    = c(c \pm u_1)(c \pm u_2)/(c^2 + u_1 u_2) > 0$.
\end{proof}


\begin{beobachtung}[Lorentzkontraktion und Zeitdilatation]
    Die Lorentz-Translation $L(v)$ kodiert unmittelbar zwei
    experimentell bestätigte Phänomene:
    \begin{itemize}
        \item[-] \textbf{Zeitdilatation}: Ein bewegtes Objekt (Geschwindigkeit $u$)
              misst eine um den Faktor $d = \sqrt{1 - u^2/c^2}$ \emph{kürzere}
              Eigenzeit -- bewegte Uhren gehen langsamer.
        \item[-] \textbf{Lorentzkontraktion}: In Bewegungsrichtung erscheinen
              Längen um den Faktor $d$ \emph{verkürzt} -- bewegte Objekte
              sind kürzer.
    \end{itemize}
    Das Einsteinsche Additionsgesetz zeigt: Auch bei Addition von
    Geschwindigkeiten nahe $c$ bleibt das Ergebnis stets kleiner als $c$
    -- die Lichtgeschwindigkeit ist eine universelle Schranke,
    die aus der Gruppenstruktur der Lorentz-Translationen folgt,
    nicht aus einer externen Annahme.

    Die algebraische Grundlage dieser Physik ist ausschließlich
    Lineare Algebra: Bilinearformen, Isometrien, Gruppenstruktur.
\end{beobachtung}

\begin{beobachtung}[Die lineare Algebra als Grammatik der Physik]
    Der Minkowskiraum ist der deutlichste Beleg, dass Lineare Algebra
    keine abstrakte Spielerei ist, sondern die \emph{Grammatik der
    Naturgesetze}:
    \begin{itemize}
        \item[-] Die \textbf{Quantenmechanik} ist in Hilberträumen formuliert
              -- unendlichdimensionalen euklidischen Räumen. Observablen
              sind hermitesche Operatoren, Messungen sind Eigenwertprobleme.
        \item[-] Die \textbf{Relativitätstheorie} ist Geometrie im
              Minkowskiraum -- einer Bilinearform mit Signatur $(3,1)$.
              Lorentz-Transformationen sind ihre Isometrien.
        \item[-] Die \textbf{Elektrodynamik} (Maxwell-Gleichungen) ist
              in der Sprache von Vierervektoren und antisymmetrischen
              Tensoren formuliert -- lineare Algebra über $\mathbb{R}^4$.
    \end{itemize}
    All diese Theorien sind in ihrem mathematischen Kern Instanzen der
    Theorie linearer Abbildungen auf Vektorräumen mit Skalarprodukt.
    Der Vektorraumaxiomatik -- aufgebaut aus Assoziativität, Distributivität
    und Skalarmultiplikation -- entspringt die gesamte mathematische
    Struktur der modernen Physik.
\end{beobachtung}

% ============================================================
% ============================================================
\chapter{Determinanten als Volumenfunktionen}
% ============================================================

\section{Die Leibniz-Formel und ihre Eigenschaften}

Die Determinante wurde bisher als Werkzeug eingeführt. Hier wird ihr
tiefer geometrischer Charakter als orientiertes Volumen sichtbar.

\begin{definition}[Volumenfunktion]
    Sei $W = K^n$ oder ein $n$-dimensionaler $K$-Vektorraum. Eine Abbildung
    $\mathrm{Vol} : W^n \to K$ heißt \textbf{Volumenfunktion}, falls:
    \begin{itemize}
        \item[-] \textbf{Multilinearität}: Für alle $i$ ist
              $v \mapsto \mathrm{Vol}(z_1, \dots, v, \dots, z_n)$ linear.
        \item[-] \textbf{Alternierend}: Ist $z_i = z_j$ für $i \neq j$,
              so gilt $\mathrm{Vol}(z_1, \dots, z_n) = 0$.
    \end{itemize}
\end{definition}

\begin{satz}[Charakterisierung der Determinante]
    Alle Volumenfunktionen auf $K^n$ haben die Gestalt
    $\mathrm{Vol}(z_1, \dots, z_n) = c \cdot \det A$, wobei $A$ die Matrix
    mit den Zeilen $z_i$ ist und $c \in K$ beliebig.

    Insbesondere gilt für jede Volumenfunktion $\mathrm{Vol}$ und alle
    $\pi \in S_n$:
    \[
        \mathrm{Vol}(z_{\pi 1}, \dots, z_{\pi n})
        = \mathrm{sgn}(\pi) \cdot \mathrm{Vol}(z_1, \dots, z_n).
    \]
    Sind ferner $w_i = \sum_j a_{ij} z_j$, so gilt
    $\mathrm{Vol}(w_1, \dots, w_n) = \det(a_{ij}) \cdot \mathrm{Vol}(z_1, \dots, z_n)$.
\end{satz}
\begin{proof}
    Sei $F$ eine multilineare alternierende Funktion. Da $[e_1, \ldots, e_n]$ einen Träger ist,
    lässt sich jede Zeile $z_i = \sum_j a_{ij} e_j$ schreiben. Durch $n$-fache
    Multilinearität erhält man:
    \[
        F(z_1, \ldots, z_n) = \sum_{\pi \in S_n} a_{1,\pi(1)} \cdots a_{n,\pi(n)}
        F(e_{\pi(1)}, \ldots, e_{\pi(n)}).
    \]
    $F(e_{\pi(1)}, \ldots, e_{\pi(n)})
    = \mathrm{sgn}(\pi)\, F(e_1, \ldots, e_n)$.
    Somit $F = F(e_1,\ldots,e_n) \cdot \det$.
\end{proof}

\begin{satz}[Multiplikationssatz]
    Für $A, B \in K^{n \times n}$ gilt:
    \[
        \det(AB) = \det A \cdot \det B.
    \]
    Insbesondere: $A$ ist genau dann invertierbar, wenn $\det A \neq 0$.
    Für invertierbare $A$ gilt $\det A^{-1} = (\det A)^{-1}$,
    und ähnliche Matrizen haben dieselbe Determinante.
\end{satz}
\begin{proof}
    Sei $F_B(z_1, \ldots, z_n) := \det(z_1 B, \ldots, z_n B)$ mit festen Zeilenvektoren
    $z_i$ und Matrix $B$. Dann ist $F_B$ multilinear und alternierend, also
    $F_B = c \cdot \det$ mit $c = F_B(e_1, \ldots, e_n) = \det B$.
    Damit $\det(AB) = (\det B) \cdot \det A = \det A \cdot \det B$.
    Ist $\det A \neq 0$, so ist $A$ invertierbar, und $\det(A^{-1}) = (\det A)^{-1}$
    folgt aus $\det(A \cdot A^{-1}) = \det E = 1$.
\end{proof}

\begin{beobachtung}[Determinante als orientiertes Volumen]
    Der Satz von der Volumenfunktion sagt: Die Determinante ist
    \emph{die einzige} multilineare alternierende Funktion auf den
    Spaltenvektoren, die auf der Einheitsmatrix den Wert $1$ annimmt.
    Geometrisch: $|\det A|$ ist das Volumen des Parallelepipeds, das
    von den Spalten von $A$ aufgespannt wird.

    Das Vorzeichen -- das Signum -- misst die Orientierung: Ob die
    Spalten eine Rechts-Träger (positiv) oder Links-Träger (negativ) bilden.
    In der Ebene $\mathbb{R}^2$ ist $\det(z_1, z_2)$ die vorzeichenbehaftete
    Fläche des Parallelogramms: positiv bei Linksdreh, negativ bei Rechtsdreh.

    Der Multiplikationssatz $\det(AB) = \det A \cdot \det B$ sagt:
    Hintereinanderschalten von Abbildungen multipliziert die Volumenskalierungen.
    Dies ist der gruppentheoretische Kern: $\det$ ist ein Homomorphismus
    von $GL(n, K)$ in $K^*$.
\end{beobachtung}

% ============================================================
\section{Die Adjunkte und die Cramersche Regel}
% ============================================================

\begin{definition}[Adjunkte]
    Sei $A = (a_{ij}) \in K^{n \times n}$. Der \textbf{Kofaktor} $A_{ij}$ ist
    \[
        A_{ij} = (-1)^{i+j} M_{ij},
    \]
    wobei $M_{ij}$ die Determinante der $(n-1) \times (n-1)$-Matrix ist,
    die durch Streichen der $i$-ten Zeile und $j$-ten Spalte entsteht.
    Die \textbf{Adjunkte} (klassische Adjungierte) ist die transponierte
    Kofaktormatrix: $\tilde{A} = (A_{ji})$.
\end{definition}

\begin{satz}[Entwicklungssatz und Inverse über Adjunkte]
    Für $A \in K^{n \times n}$ gilt:
    \[
        A \tilde{A} = \tilde{A} A = (\det A) \cdot E.
    \]
    Der \textbf{Laplacescher Entwicklungssatz} nach Zeile $i$:
    \[
        \det A = \sum_{j=1}^n a_{ij} A_{ij} = \sum_{j=1}^n a_{ij} (-1)^{i+j} M_{ij}.
    \]
    Ist $\det A \neq 0$, so gilt $A^{-1} = (\det A)^{-1} \tilde{A}$.
\end{satz}
\begin{proof}
    Das $(i,k)$-Element von $A\tilde{A}$ ist $\sum_j a_{ij} A_{kj}$.
    Für $i = k$ ist dies die Entwicklung von $\det A$ nach Zeile $i$: $= \det A$.
    Für $i \neq k$ entwickelt man $\det A'$, wobei $A'$ aus $A$ entsteht,
    indem man Zeile $k$ durch Zeile $i$ ersetzt. Da $A'$ zwei gleiche
    Zeilen hat, gilt $\det A' = 0$. Also $A\tilde{A} = (\det A) E$.
\end{proof}

\begin{satz}[Cramersche Regel]
    Sei $A \in K^{n \times n}$ mit $\det A \neq 0$ und $b \in K^n$.
    Dann hat das Gleichungssystem $Ax = b$ die eindeutige Lösung $x = (x_i)$
    mit
    \[
        x_i = \frac{\det A_i}{\det A},
    \]
    wobei $A_i$ die Matrix entsteht, indem die $i$-te Spalte von $A$
    durch den Vektor $b$ ersetzt wird.
\end{satz}
\begin{proof}
    Mit $x = A^{-1}b = (\det A)^{-1}\tilde{A} b$ gilt
    $x_i = (\det A)^{-1} \sum_j A_{ji} b_j$.
    Nach dem Entwicklungssatz nach Spalte $i$ (mit $b$ als Spalte $i$)
    ist $\sum_j A_{ji} b_j = \det A_i$.
\end{proof}

\begin{beobachtung}[Cramersche Regel: Theorie vs. Praxis]
    Die Cramersche Regel ist ein elegantes theoretisches Resultat:
    Sie zeigt, dass die Lösung rationaler Gleichungssysteme (mit
    $a_{ij}, b_j \in \mathbb{Q}$) wieder rational ist.

    Für die \emph{numerische} Praxis ist sie jedoch ungeeignet:
    Die Berechnung von $n+1$ Determinanten der Größe $n$ erfordert
    $O(n \cdot n!) $ Operationen -- exponentiell in $n$. Der Gaußsche
    Algorithmus benötigt nur $O(n^3)$.
    Die Cramersche Regel wird daher in der Praxis nur für $n \leq 3$
    verwendet, sowie in algebraischen Beweisen über Definitionsbereiche
    (Ganzzahligkeit von Lösungen).
\end{beobachtung}

\begin{beobachtung}[Die Graßmann-Algebra als Verallgemeinerung]
    Die Alternierungseigenschaft ist das Herzstück der
    \textbf{Gra\ss{}mann-Algebra} (äußere Algebra): Das äußere Produkt
    $z_1 \wedge z_2 \wedge \cdots \wedge z_n$ ist eine solche
    alternierende Multilinearform.

    Die Gra\ss{}mann-Algebra $\bigwedge V$ enthält als höchste Stufe
    $\bigwedge^n V$ -- einen eindimensionalen Raum, auf dem
    jede lineare Abbildung $A : V \to V$ durch einen einzigen Skalar
    wirkt: eben $\det A$. Die Determinante ist damit das
    geometrisch natürliche Objekt der äußeren Algebra -- nicht
    eine ad-hoc-Formel, sondern die unvermeidliche Konsequenz der
    Alternierungsstruktur.
\end{beobachtung}

% ============================================================
\chapter{Normierte Vektorräume und Matrixexponential}
% ============================================================

\section{Normierte Vektorräume}

\begin{definition}[Norm und normierter Vektorraum]
    Sei $V$ ein $K$-Vektorraum ($K = \mathbb{R}$ oder $\mathbb{C}$).
    Eine Abbildung $\|\cdot\| : V \to \mathbb{R}_{\geq 0}$ heißt
    \textbf{Norm}, falls für alle $v, w \in V$ und $a \in K$ gilt:
    \begin{itemize}
        \item[-] $\|v\| = 0 \Leftrightarrow v = 0$ \quad (Definitheit)
        \item[-] $\|av\| = |a| \cdot \|v\|$ \quad (Homogenität)
        \item[-] $\|v + w\| \leq \|v\| + \|w\|$ \quad (Dreiecksungleichung)
    \end{itemize}
    $V$ mit einer Norm heißt \textbf{normierter Vektorraum}.
\end{definition}

\begin{satz}[Äquivalenz aller Normen in endlicher Dimension]
    Sei $V$ ein endlichdimensionaler $K$-Vektorraum. Dann sind je zwei
    Normen auf $V$ äquivalent: Zu $\|\cdot\|$ und $\|\cdot\|'$
    gibt es $c, d > 0$ mit
    \[
        d\, \|v\| \leq \|v\|' \leq c\, \|v\| \qquad \forall v \in V.
    \]
    Insbesondere definieren alle Normen dieselben konvergenten Folgen
    und dieselben stetigen Abbildungen.
\end{satz}
\begin{proof}
    Es genügt zu zeigen, dass jede Norm $\|\cdot\|$ zur Euklidischen Norm
    $\|\cdot\|_2$ äquivalent ist. Sei $[e_1, \ldots, n]$ einen Träger.
    \textbf{Obere Schranke:} $\|v\| = \|\sum a_i e_i\| \leq \sum |a_i| \|e_i\|
    \leq C \|v\|_2$ mit $C = \sqrt{n} \max_i \|e_i\|$ (nach Cauchy-Schwarz).
    \textbf{Untere Schranke:} Die Einheitssphäre $S = \{v \mid \|v\|_2 = 1\}$
    ist kompakt. Die Funktion $v \mapsto \|v\|$ ist stetig bezüglich $\|\cdot\|_2$
    (wegen der oberen Schranke) und positiv auf $S$. Daher nimmt sie auf $S$
    ein Minimum $d > 0$ an: $\|v\| \geq d \|v\|_2$ für alle $v$.
\end{proof}


\begin{satz}[Stetige lineare Abbildungen]
    Seien $V$ und $W$ normierte Vektorräume. Eine lineare Abbildung
    $A : V \to W$ ist genau dann \textbf{stetig}, wenn sie
    \textbf{beschränkt} ist: $\|Av\| \leq M\|v\|$ für ein $M > 0$
    und alle $v \in V$. Ist $\dim V < \infty$, so ist jede lineare
    Abbildung stetig.
\end{satz}
\begin{proof}
    \textbf{Beschränkt $\Rightarrow$ stetig:} Aus $\|Av\| \leq M\|v\|$ folgt
    $\|Av - Aw\| = \|A(v-w)\| \leq M\|v-w\| \to 0$.
    \textbf{Stetig $\Rightarrow$ beschränkt:} Ist $A$ stetig in $0$, so gibt es
    $\delta > 0$ mit $\|Av\| \leq 1$ für $\|v\| \leq \delta$.
    Für $v \neq 0$ gilt $\|A(\delta v / \|v\|)\| \leq 1$, also $\|Av\| \leq \frac{1}{\delta}\|v\|$.
    \textbf{Für $\dim V < \infty$:} Alle Normen sind äquivalent; da $A$ bezüglich
    der Euklidischen Norm beschränkt ist (wie oben), ist $A$ stetig.
\end{proof}


\begin{beobachtung}[Warum Äquivalenz aller Normen so wichtig ist]
    In endlicher Dimension gibt es keine sinnvolle Unterscheidung zwischen
    verschiedenen Normen -- alle erzeugen dieselbe Topologie. Dies ist
    in unendlicher Dimension grundlegend anders: $L^1$ und $L^2$ sind
    nicht äquivalent, und die Wahl der Norm bestimmt die Physik.

    In endlicher Dimension bedeutet der Satz: Konvergenz, Stetigkeit,
    Kompaktheit -- all diese Begriffe sind koordinatenunabhängig.
    Eine Folge konvergiert genau dann, wenn sie komponentenweise konvergiert,
    unabhängig davon, mit welcher Norm man misst. Dies ist die
    mathematische Rechtfertigung dafür, dass Koordinaten in der
    Linearen Algebra frei gewählt werden können.
\end{beobachtung}

% ============================================================
\section{Die Matrixexponentialfunktion}
% ============================================================

\begin{definition}[Matrixexponential]
    Da die Potenzreihe $e^z = \sum_{k=0}^\infty \frac{z^k}{k!}$ den
    Konvergenzradius $\infty$ hat, konvergiert für jede Matrix
    $A \in \mathbb{C}^{n \times n}$ die Reihe
    \[
        e^A = \sum_{k=0}^\infty \frac{1}{k!} A^k.
    \]
    Man nennt $e^A$ die \textbf{Matrixexponentialfunktion}.
\end{definition}

\begin{satz}[Eigenschaften des Matrixexponentials]
    Sei $A \in \mathbb{C}^{n \times n}$.
    \begin{itemize}
        \item[-] \textbf{Eigenwerte}: Hat $A$ die Eigenwerte $a_j$, so hat
              $e^A$ die Eigenwerte $e^{a_j}$. Insbesondere:
              $\det e^A = e^{\mathrm{Sp}\, A}$.
        \item[-] \textbf{Kommutativität}: Gilt $AB = BA$, so folgt
              $e^{A+B} = e^A e^B$. Insbesondere ist $(e^A)^{-1} = e^{-A}$.
        \item[-] \textbf{Ableitung}: Für $F(t) = e^{tA}$ gilt
              $F'(t) = A e^{tA} = e^{tA} A$.
        \item[-] \textbf{Differentialgleichung}: Ist $y'(t) = Ay(t)$,
              so gilt $y(t) = e^{tA} y(0)$.
    \end{itemize}
\end{satz}
\begin{proof}
    \textbf{Konvergenz:} $\sum_{k=0}^\infty \frac{\|A\|^k}{k!} = e^{\|A\|} < \infty$,
    also konvergiert $e^A$ absolut.
    \textbf{Eigenwerte:} Ist $Av = av$, so $e^A v = \sum_k \frac{A^k}{k!} v
    = \sum_k \frac{a^k}{k!} v = e^a v$.
    \textbf{Kommutativität:} Das Cauchy-Produkt der absolut konvergenten Reihen
    $e^A e^B = (\sum_j A^j/j!)(\sum_k B^k/k!)$ ergibt $\sum_n \frac{1}{n!}(A+B)^n = e^{A+B}$,
    sofern $AB = BA$ (damit der Binomische Lehrsatz gilt).
    \textbf{Ableitung:} $\frac{d}{dt} e^{tA} = \sum_{k=1}^\infty \frac{k t^{k-1}}{k!} A^k
    = A \sum_{k=0}^\infty \frac{t^k}{k!} A^k = A e^{tA}$.
\end{proof}


\begin{beobachtung}[Matrixexponential und lineare DGL-Systeme]
    Das Matrixexponential ist die vollständige Antwort auf die Frage nach
    der Lösung linearer Differentialgleichungssysteme $\dot{x} = Ax$.
    Die Lösung $x(t) = e^{tA} x(0)$ gilt für jede Matrix $A$ --
    auch wenn $A$ nicht diagonalisierbar ist.

    In Verbindung mit der Jordanschen Normalform ergibt sich die
    die vollständige Lösungsstruktur: Für $J_m(\lambda)$
    ist
    \[
        e^{tJ_m(\lambda)} = e^{\lambda t}
        \begin{pmatrix}
            1 & t & \frac{t^2}{2!} & \cdots \\
              & 1 & t & \cdots \\
              & & \ddots & \ddots \\
              & & & 1
        \end{pmatrix}.
    \]
    Mehrfache Eigenwerte erzeugen polynomiales Wachstum $t^k e^{\lambda t}$,
    einfache nur exponentielles $e^{\lambda t}$. Die algebraische Struktur
    der Jordanform bestimmt vollständig das Langzeitverhalten dynamischer
    Systeme.
\end{beobachtung}

\begin{satz}[Eulersche Formel und komplexe Exponentialfunktion]
    Für $z = x + iy \in \mathbb{C}$ gilt:
    \[
        e^z = e^x (\cos y + i \sin y), \qquad |e^z| = e^x.
    \]
    Insbesondere ist $e^{iy} = \cos y + i \sin y$ (Eulersche Formel) und
    $\cos y = \frac{1}{2}(e^{iy} + e^{-iy})$,
    $\sin y = \frac{1}{2i}(e^{iy} - e^{-iy})$.
    Für jedes $0 \neq c \in \mathbb{C}$ gibt es ein $w \in \mathbb{C}$ mit
    $c = e^w$.
\end{satz}
\begin{proof}
    Für $z = iy$ ($y \in \mathbb{R}$) gilt:
    $e^{iy} = \sum_{k=0}^\infty \frac{(iy)^k}{k!}
    = \sum_{j=0}^\infty \frac{(-1)^j y^{2j}}{(2j)!}
    + i \sum_{j=0}^\infty \frac{(-1)^j y^{2j+1}}{(2j+1)!}
    = \cos y + i \sin y$.
    Dies sind gerade die Taylor-Reihen von $\cos$ und $\sin$.
    Für allgemeines $z = x + iy$: $e^z = e^x e^{iy} = e^x(\cos y + i \sin y)$,
    also $|e^z| = e^x$. Die Surjektivität von $e^{\cdot}$ auf $\mathbb{C}^*$
    folgt: Für $c = |c| e^{i\arg c}$ gilt $c = e^{\ln|c| + i\arg c}$.
\end{proof}


\begin{beobachtung}[Die Eulersche Formel als Brücke zwischen Analysis und Algebra]
    Die Eulersche Formel $e^{i\pi} + 1 = 0$ verknüpft die fünf
    fundamentalsten Konstanten der Mathematik. Mathematisch tiefer
    ist ihre strukturelle Bedeutung: Die Abbildung $y \mapsto e^{iy}$
    ist ein Gruppenhomomorphismus von $(\mathbb{R}, +)$ auf die
    Einheitskreisgruppe $\{z \in \mathbb{C} \mid |z| = 1\} \cong S^1$.

    In der Matrixwelt: Eine schiefsymmetrische Matrix $A$ (d.\,h. $A + A^t = 0$)
    erzeugt über $e^{tA}$ eine orthogonale Abbildung für alle $t$ --
    die Exponentialfunktion ist die Brücke von der Lie-Algebra zur
    Lie-Gruppe, von den infinitesimalen Drehungen zu den Drehungen selbst.
\end{beobachtung}

% ============================================================
\chapter{Hilberträume und adjungierte Abbildungen}
% ============================================================

\section{Endlichdimensionale Hilberträume}

\begin{definition}[Hilbertraum]
    Sei $V$ ein Vektorraum über $\mathbb{R}$ oder $\mathbb{C}$ mit einem
    \textbf{definiten Skalarprodukt} $(\cdot, \cdot)$: symmetrisch (bzw.\
    hermitesch), bilinear (bzw.\ sesquilinear) und $(v,v) > 0$ für $v \neq 0$.
    Dann definiert $\|v\| = \sqrt{(v,v)}$ eine Norm.

    Ist $V$ vollständig bezüglich dieser Norm, so heißt $V$
    \textbf{Hilbertraum}. Jeder endlichdimensionale Vektorraum mit
    definitem Skalarprodukt ist automatisch ein Hilbertraum.
\end{definition}

\begin{satz}[Orthogonalisierungsverfahren von E. Schmidt]
    Sei $V$ ein Hilbertraum und $v_1, \dots, v_m$
    linear unabhängige Vektoren.
    $\langle v_1, \dots, v_m \rangle$ mit $(w_j, w_k) = \delta_{jk}$,
    konstruiert durch:
    \[
        w_j = c_j \!\left(v_j - \sum_{k=1}^{j-1} (v_j, w_k) w_k\right),
        \qquad c_j > 0 \text{ Normierungsfaktor.}
    \]
    Insbesondere besitzt jeder endlichdimensionale Hilbertraum eine
    Orthonormalträger.
\end{satz}
\begin{proof}
    Wir konstruieren $w_j$ induktiv. Setze $\tilde{w}_1 = v_1 \neq 0$ und
    $w_1 = \tilde{w}_1 / \|\tilde{w}_1\|$.
    Für $j > 1$ setze $\tilde{w}_j = v_j - \sum_{k=1}^{j-1} (v_j, w_k) w_k$.
    Dann $(w_l, \tilde{w}_j) = (w_l, v_j) - \sum_k (v_j, w_k)(w_l, w_k)
    = (w_l, v_j) - (v_j, w_l) = 0$ für $l < j$.
    Da $[v_1, \ldots, v_j]$ linear unabhängig ist, gilt $\tilde{w}_j \neq 0$,
    und $w_j = \tilde{w}_j / \|\tilde{w}_j\|$ erfüllt $(w_j, w_l) = \delta_{jl}$.
    Per Induktion gilt $\langle w_1, \ldots, w_j \rangle = \langle v_1, \ldots, v_j \rangle$.
\end{proof}


\begin{algorithm}[H]
\DontPrintSemicolon
\caption{Gram-Schmidt-Orthogonalisierung}\label{alg:schmidt}
\KwIn{Linear unabhängige Vektoren $v_1, \ldots, v_m \in V$}
\KwOut{Orthonormalträger $w_1, \ldots, w_m$}
\For{$j \leftarrow 1$ \KwTo $m$}{
    $\tilde{w}_j \leftarrow v_j$\;
    \For{$k \leftarrow 1$ \KwTo $j-1$}{
        $\tilde{w}_j \leftarrow \tilde{w}_j - (v_j, w_k)\cdot w_k$\;
    }
    $w_j \leftarrow \tilde{w}_j / \|\tilde{w}_j\|$\;
}
\Return{$w_1, \ldots, w_m$}\;
\BlankLine
\textbf{Aufwand:} $\Theta(m^2 \cdot n)$ innere Produkte ($n = \dim V$)
\end{algorithm}


\begin{beobachtung}[Parallelogrammgleichung als Charakterisierung]
    Nicht jede Norm stammt von einem Skalarprodukt -- aber man kann es
    testen: Eine Norm $\|\cdot\|$ kommt genau dann von einem Skalarprodukt,
    wenn die \textbf{Parallelogrammgleichung} gilt:
    \[
        \|v_1 + v_2\|^2 + \|v_1 - v_2\|^2 = 2(\|v_1\|^2 + \|v_2\|^2).
    \]
    (J.\,von Neumann.) In diesem Fall kann das Skalarprodukt durch die
    Polarisierungsidentität zurückgewonnen werden.

    Dies unterscheidet Hilberträume von allgemeinen Banachräumen:
    Die geometrische Eigenschaft der Parallelogrammgleichung ist äquivalent
    zur algebraischen Existenz eines Skalarprodukts.
\end{beobachtung}

% ============================================================
\section{Adjungierte Abbildungen}
% ============================================================

\begin{definition}[Adjungierte Abbildung]
    Sei $V$ ein Hilbertraum und $A \in \mathrm{End}(V)$. Die
    \textbf{adjungierte Abbildung} $A^*$ ist die eindeutig bestimmte
    Abbildung mit
    \[
        (Av, w) = (v, A^*w) \qquad \forall v, w \in V.
    \]
    Bezüglich einer Orthonormalträger gilt: Die Matrix zu $A^*$ ist die
    transponiert-konjugiert-komplexe Matrix $\overline{A}^t$.

    $A$ heißt \textbf{normal}, falls $AA^* = A^*A$;
    \textbf{selbstadjungiert} (hermitesch), falls $A = A^*$;
    \textbf{unitär}, falls $AA^* = A^*A = E$.
\end{definition}

\begin{satz}[Spektralsatz für normale Abbildungen]
    Sei $V$ ein endlichdimensionaler Hilbertraum über $\mathbb{C}$ und
    $A \in \mathrm{End}(V)$ normal. Dann gilt:
    \begin{itemize}
        \item[-] Zu verschiedenen Eigenwerten gehören orthogonale Eigenräume.
        \item[-] $V$ hat eine Orthonormalträger aus Eigenvektoren von $A$.
        \item[-] $A^*v = \bar{a}v$ für jeden Eigenvektor $Av = av$.
        \item[-] Für selbstadjungiertes $A$: alle Eigenwerte sind reell.
        \item[-] Für unitäres $A$: alle Eigenwerte liegen auf dem Einheitskreis.
    \end{itemize}
\end{satz}
\begin{proof}
    \textbf{Orthogonalität verschiedener Eigenräume:} Sei $Av = av$ und $Aw = bw$
    mit $a \neq b$. Da $AA^* = A^*A$, gilt für einen normalen Operator
    $\|Av\|^2 = (Av, Av) = (v, A^*Av) = (v, AA^*v) = \|A^*v\|^2$,
    also $\|A^*v\| = \|Av\|$. Daraus folgt $A^*v = \bar{a}v$.
    Dann $(v, w) \cdot a = (av, w) = (Av, w) = (v, A^*w) = (v, \bar{b}w) = b(v,w)$.
    Da $a \neq b$, folgt $(v, w) = 0$.
    \textbf{Existenz einer ONB:} Ãœber $\mathbb{C}$ hat $A$ stets mindestens
    einen Eigenwert $a$ (Fundamentalsatz der Algebra). Der Eigenraum $E_a$ ist
    $A$-invariant, und $A|_{E_a^\perp}$ ist wieder normal; durch Induktion
    über $\dim V$ erhält man eine vollständige ONB aus Eigenvektoren.
\end{proof}


\begin{beobachtung}[Der Spektralsatz als Fundament der Quantenmechanik]
    Der Spektralsatz für selbstadjungierte Abbildungen ist das mathematische
    Fundament der Quantenmechanik: Observablen sind selbstadjungierte
    Operatoren auf einem Hilbertraum. Ihre Eigenwerte sind die möglichen
    Messwerte -- und wegen Selbstadjungiertheit sind diese stets reell,
    also physikalisch sinnvoll.

    Unitäre Abbildungen beschreiben die Zeitentwicklung quantenmechanischer
    Systeme ($U = e^{-iHt/\hbar}$ mit selbstadjungiertem $H$): Da
    unitäre Abbildungen normerhaltend sind, bleibt die
    Gesamtwahrscheinlichkeit erhalten -- die Physik der Wahrscheinlichkeitserhaltung
    ist die Unitarität des Zeitentwicklungsoperators.

    Normale Abbildungen sind der gemeinsame Rahmen: diagonalisierbar
    bezüglich einer Orthonormalträger genau dann, wenn normal.
\end{beobachtung}

\begin{satz}[Eigenwertabschätzungen nach Courant--Fischer]
    Sei $V$ ein Hilbertraum der Dimension $n$ und $A = A^*$ mit
    Eigenwerten $a_1 \geq \dots \geq a_n$. Dann gilt das
    \textbf{Minimax-Prinzip}:
    \[
        a_k = \min_{\substack{W \leq V \\ \dim W = n-k+1}}
        \max_{\substack{w \in W \\ \|w\| = 1}} (Aw, w).
    \]
    Wird $A$ auf einen Unterraum $U$ der Kodimension $1$ eingeschränkt
    mit Eigenwerten $b_1 \geq \dots \geq b_{n-1}$, so gilt der
    \textbf{Verschachtelungssatz}:
    \[
        a_1 \geq b_1 \geq a_2 \geq b_2 \geq \cdots \geq b_{n-1} \geq a_n.
    \]
\end{satz}
\begin{proof}
    Sei $[v_1, \ldots, v_n]$ eine ONB aus Eigenvektoren mit $Av_j = a_j v_j$.
    \textbf{Minimax-Prinzip:} Für jeden $(n-k+1)$-dimensionalen Unterraum $W$
    enthält $W$ einen Vektor $w = \sum_{j \geq k} c_j v_j$ (orthogonal zu
    $v_1, \ldots, v_{k-1}$); dann $(Aw, w) = \sum_{j \geq k} a_j |c_j|^2
    / \sum |c_j|^2 \geq a_k$. Das Minimum wird für $W = \langle v_k, \ldots, v_n \rangle$
    mit $w = v_k$ angenommen.
    \textbf{Verschachtelung:} Sei $U = \{v \mid (v, u) = 0\}$ für ein festes $u$.
    Für jedes $k$ enthält $E_k := \langle v_k, \ldots, v_n \rangle \cap U$
    einen Vektor (da $\dim E_k = n-k+1 > n - (k+1) + 1 - 1$), und das
    Minimax-Prinzip angewandt auf die Einschränkung liefert $b_k \leq a_k \leq b_{k-1}$.
\end{proof}


\begin{beobachtung}[Verschachtelungssatz als Stabilitätsprinzip]
    Der Verschachtelungssatz sagt: Entfernt man eine Dimension,
    verschieben sich die Eigenwerte -- aber nur innerhalb enger Grenzen.
    Die Eigenwerte sind stabil gegenüber Einschränkungen der Abbildung.

    In der Ingenieurpraxis: Die Eigenfrequenzen eines schwingenden Systems
    können durch Hinzufügen einer Einspannung (Reduktion der Freiheitsgrade)
    nur innerhalb des durch den Verschachtelungssatz vorgegebenen Rahmens
    verändert werden. Die Mathematik garantiert damit die physikalische
    Stabilität der Eigenfrequenzen.
\end{beobachtung}

% ============================================================

% ── Erweiterung ──────────────────────────────────────────────────────────
\section{Anwendung: Kleinste Quadrate und Pseudoinverse}

\begin{definition}[Kleinste-Quadrate-Problem]
    Sei $A \in \mathbb{R}^{m \times n}$ mit $m > n$ (überbestimmtes System)
    und $b \in \mathbb{R}^m$. Das \textbf{Kleinste-Quadrate-Problem} sucht
    \[
        x^* = \arg\min_{x \in \mathbb{R}^n} \|Ax - b\|^2.
    \]
    Die Lösung $x^*$ minimiert die Summe der quadratischen Residuen --
    die ``kleinsten Quadrate'' der Abweichungen.
\end{definition}

\begin{satz}[Normalengleichungen]
    Das Kleinste-Quadrate-Problem hat die Lösung(en) $x^*$ genau durch
    die \textbf{Normalengleichungen}:
    \[
        A^t A x^* = A^t b.
    \]
    Ist $\mathrm{Rang}\,A = n$, so ist $A^t A \in \mathbb{R}^{n \times n}$
    positiv definit und invertierbar, und die eindeutige Lösung ist
    $x^* = (A^t A)^{-1} A^t b =: A^+ b$, wobei $A^+ = (A^t A)^{-1} A^t$
    die \textbf{Moore-Penrose-Pseudoinverse} von $A$ heißt.
\end{satz}
\begin{proof}
    Die Residuumsfunktion $f(x) = \|Ax - b\|^2 = (Ax-b)^t(Ax-b)$
    ist konvex und differenzierbar. Das Minimum erfüllt
    $\nabla_x f = 2A^t(Ax - b) = 0$, also $A^t A x = A^t b$.
    Geometrisch: $Ax^*$ ist die Orthogonalprojektion von $b$ auf
    $\mathrm{Im}\,A$, denn $b - Ax^* \perp \mathrm{Im}\,A$
    ist äquivalent zu $A^t(b - Ax^*) = 0$.
\end{proof}

\begin{beobachtung}[Pseudoinverse und SVD]
    Die Moore-Penrose-Pseudoinverse lässt sich elegant über die
    singuläre Wertzerlegung $A = U\Sigma V^t$ ausdrücken:
    $A^+ = V \Sigma^+ U^t$, wobei $\Sigma^+$ aus $\Sigma$ entsteht,
    indem man jeden Nichtnull-Diagonaleintrag $\sigma_j$ durch
    $1/\sigma_j$ ersetzt.

    Dies ist die stabilste numerische Methode zur Lösung von
    Kleinste-Quadrate-Problemen: Kleine singuläre Werte, die
    numerische Instabilität signalisieren, können durch
    \textbf{Regularisierung} (Setzen auf $0$ ab einem Schwellwert)
    kontrolliert werden.
\end{beobachtung}
\section{Selbstadjungierte Operatoren}

\begin{beobachtung}[Der harmonische Oszillator in der Quantenmechanik]
    Der \textbf{quantenmechanische harmonische Oszillator} ist der
    prototypische selbstadjungierte Operator. Der Hamilton-Operator
    \[
        H = -\frac{\hbar^2}{2m}\frac{d^2}{dx^2} + \frac{1}{2}m\omega^2 x^2
    \]
    wirkt auf dem Hilbertraum $L^2(\mathbb{R})$. Seine Eigenwerte
    (die Energieniveaus) sind $E_n = \hbar\omega(n + \frac{1}{2})$
    für $n = 0, 1, 2, \ldots$ -- äquidistant, ausschließlich reell
    (da $H = H^*$), und nach unten beschränkt.

    Die zugehörigen Eigenfunktionen sind die \textbf{Hermiteschen Polynome},
    \[
        \psi_n(x) = H_n(x) e^{-x^2/2},
    \]
    Jeder Zustand $\psi \in L^2(\mathbb{R})$ entwickelt sich eindeutig
    in dieser Trägerbasis: $\psi = \sum_n c_n \psi_n$ -- der
    Spektralsatz für selbstadjungierte Operatoren in
    unendlicher Dimension.
\end{beobachtung}

\chapter{Quaternionen und die Lie-Algebra}
% ============================================================

\section{Die Hamiltonschen Quaternionen}

\begin{satz}[Der Schiefkörper der Quaternionen]
    Im Matrixring $(\mathbb{C})_2$ ist
    \[
        \mathbb{H} = \left\{ \begin{pmatrix} a & -\overline{b} \\ b & \overline{a} \end{pmatrix}
        \;\middle|\; a, b \in \mathbb{C} \right\}
    \]
    ein \textbf{Schiefkörper} -- jedes $0 \neq q \in \mathbb{H}$ ist
    invertierbar mit $q^{-1} = \frac{1}{|a|^2+|b|^2}\bigl(\begin{smallmatrix}
    \overline{a} & \overline{b} \\ -b & a \end{smallmatrix}\bigr)$.

    Die $\mathbb{R}$-Träger von $\mathbb{H}$ bilden die Elemente
    $e_0 = E$, $e_1$, $e_2$, $e_3$ mit den Relationen
    $e_j^2 = -e_0$ für $j = 1,2,3$ und $e_1 e_2 = -e_2 e_1 = e_3$,
    $e_2 e_3 = -e_3 e_2 = e_1$, $e_3 e_1 = -e_1 e_3 = e_2$.
\end{satz}
\begin{proof}
    Die Menge $\mathbb{H}$ ist abgeschlossen unter Addition und Multiplikation
    (durch direktes Nachrechnen der Matrizenmultiplikation). Die Einheitsmatrix
    $E \in \mathbb{H}$ ist das multiplikative Einselement.
    Für $q = \bigl(\begin{smallmatrix} a & -\bar{b} \\ b & \bar{a} \end{smallmatrix}\bigr)
    \neq 0$ gilt $\det q = |a|^2 + |b|^2 > 0$, also ist $q$ invertierbar
    mit $q^{-1} = \frac{1}{|a|^2+|b|^2}\bigl(\begin{smallmatrix}\bar{a} & \bar{b}
    \\ -b & a \end{smallmatrix}\bigr) \in \mathbb{H}$.
    Die Nichtkommutativität: $e_1 e_2 = e_3 \neq -e_3 = e_2 e_1$.
    Die $\mathbb{R}$-Trägerrelationen $e_j^2 = -e_0$ verifiziert man durch
    direktes Einsetzen der Matrizen.
\end{proof}


\begin{beobachtung}[Quaternionen und Hamiltons Entdeckung]
    William Rowan Hamilton suchte jahrelang nach einer dreidimensionalen
    Erweiterung der komplexen Zahlen -- und erkannte 1843, dass dies nur
    in vier Dimensionen möglich ist, wenn man die Kommutativität aufgibt.

    Die Relationen $e_1 e_2 = e_3$, $e_2 e_3 = e_1$, $e_3 e_1 = e_2$
    sind exakt die Regeln des Vektorprodukts im $\mathbb{R}^3$.
    Hamiltons Quaternionen kodieren die Geometrie des dreidimensionalen
    Raums in einer algebraischen Struktur -- lange vor der Einführung
    der modernen Vektorrechnung.
\end{beobachtung}

\begin{satz}[Quaternionennorm und Vierfachquadratsatz]
    Für $q = \sum_{j=0}^3 x_j e_j \in \mathbb{H}$ ist die Norm
    $N(q) = \sum_{j=0}^3 x_j^2 = \det q$. Es gilt $qq^* = q^*q = N(q)e_0$
    und $N(q_1 q_2) = N(q_1) N(q_2)$. Daraus folgt:
    \[
        \left(\sum_{j=0}^3 x_j^2\right)\!\left(\sum_{j=0}^3 y_j^2\right)
        = \sum_{j=0}^3 z_j^2,
    \]
    mit $z_0 = x_0 y_0 - x_1 y_1 - x_2 y_2 - x_3 y_3$
    und $z_1 = x_0 y_1 + x_1 y_0 + x_2 y_3 - x_3 y_2$ usw.
    (\textbf{Vierfach\-quadrat\-satz}: Das Produkt zweier
    Summen von vier Quadraten ist wieder eine Summe von vier Quadraten.)
    Quadraten ist wieder eine Summe von vier Quadraten.)
\end{satz}
\begin{proof}
    Für $q = \bigl(\begin{smallmatrix} a & -\bar{b} \\ b & \bar{a} \end{smallmatrix}\bigr)$
    gilt $q^* = \bigl(\begin{smallmatrix} \bar{a} & \bar{b} \\ -b & a \end{smallmatrix}\bigr)$
    und $qq^* = (|a|^2 + |b|^2) E = N(q) E$.
    Die Multiplikativität $N(q_1 q_2) = N(q_1) N(q_2)$ folgt aus
    $\det(q_1 q_2) = \det(q_1) \det(q_2)$. Für $q_j = x_{j,0} e_0 + \cdots
    + x_{j,3} e_3$ mit reellen $x_{j,k}$ liefert das Ausrechnen von
    $N(q_1 q_2) = N(q_1) N(q_2)$ komponentenweise die Vierfachquadrat-Identität.
\end{proof}


% ============================================================
\section{Orthogonale Abbildungen und die Lie-Algebra}
% ============================================================

\begin{definition}[Orthogonale Gruppe und Lie-Algebra]
    Eine lineare Abbildung $A \in \mathrm{End}(V)$ eines euklidischen
    Vektorraums heißt \textbf{orthogonal}, falls $A^t A = E$.
    Die Gruppe aller orthogonalen Abbildungen heißt $O(V)$, die
    Untergruppe mit $\det A = 1$ heißt $SO(V)$ (Drehgruppe).

    Die \textbf{Lie-Algebra} $\mathfrak{so}(V)$ ist die Menge der
    schiefsymmetrischen Endomorphismen $\{X \mid X + X^t = 0\}$
    mit der Lie-Klammer $[X, Y] = XY - YX$.
    Für $X \in \mathfrak{so}(V)$ ist $e^{tX} \in SO(V)$ für alle $t \in \mathbb{R}$.
\end{definition}

\begin{beobachtung}[Die Lie-Algebra als infinitesimale Drehungen]
    Die Lie-Algebra $\mathfrak{so}(3)$ ist dreidimensional. Die Standardträger
    $L_x, L_y, L_z$ erzeugt infinitesimale Drehungen um die drei Achsen,
    und ihre Lie-Klammern $[L_x, L_y] = L_z$, $[L_y, L_z] = L_x$,
    $[L_z, L_x] = L_y$ sind genau die Relationen des Vektorprodukts.

    In der Physik sind $L_x, L_y, L_z$ die \emph{Drehimpulsoperatoren}
    der Quantenmechanik. Ihre Kommutatorrelationen bestimmen die möglichen
    Quantenzahlen des Drehimpulses. Die Verbindung zwischen Quaternionen,
    Lie-Algebra und Quantenmechanik ist keine Metapher -- sie ist
    strukturelle Identität.
\end{beobachtung}

\begin{satz}[Quaternionen und \texorpdfstring{$SO(3)$}{SO(3)}]
    Sei $S^3 = \{q \in \mathbb{H} \mid N(q) = 1\}$ die Einheitssphäre
    in $\mathbb{H}$. Für $q \in S^3$ ist die Konjugation
    $\varphi_q(v) = q v q^{-1}$ auf dem Unterraum der imaginären
    Quaternionen $\cong \mathbb{R}^3$ eine Drehung. Es gibt einen
    surjektiven Homomorphismus
    \[
        \Phi : S^3 \longrightarrow SO(3) \quad \text{mit} \quad
        \ker \Phi = \{e_0, -e_0\} \cong \mathbb{Z}/2\mathbb{Z}.
    \]
\end{satz}
\begin{proof}[Beweisskizze]
    Der imaginäre Unterraum $V_0 = \{q \in \mathbb{H} \mid S(q) = 0\}
    = \langle e_1, e_2, e_3 \rangle \cong \mathbb{R}^3$ ist invariant unter
    $\varphi_q : v \mapsto qvq^{-1}$ (da $S(qvq^{-1}) = S(v) = 0$).
    Isometrie: $N(qvq^{-1}) = N(q)N(v)N(q)^{-1} = N(v)$, also erhält
    $\varphi_q$ das Skalarprodukt auf $V_0$. Determinante $1$: Der Weg
    $t \mapsto \varphi_{\cos t + \sin t \, e_1}$ verbindet $\mathrm{id}$ mit
    $\varphi_q$ stetig in $O(V_0)$; da $\det = \pm 1$ und der Ausdruck
    stetig ist, folgt $\det \varphi_q = 1$.
    Kern: $\varphi_q = \mathrm{id}$ gdw.\ $qv = vq$ für alle $v \in V_0$,
    was $q \in \mathbb{R} e_0$ impliziert; mit $N(q) = 1$ folgt $q = \pm e_0$.
\end{proof}


\begin{beobachtung}[Warum \texorpdfstring{$SO(3)$}{SO(3)} doppelt überlagert wird]
    Der Kern $\ker \Phi = \{e_0, -e_0\}$ bedeutet: Genau zwei Quaternionen
    $q$ und $-q$ erzeugen dieselbe Drehung im $\mathbb{R}^3$. Dies ist
    ein fundamentales Phänomen -- die doppelte Überlagerung $S^3 \to SO(3)$.

    In der Quantenmechanik erklärt dies den \emph{Spin}: Ein Elektron
    muss um $720^\circ$ gedreht werden, um in seinen Ausgangszustand
    zurückzukehren -- weil der Zustandsraum nicht $SO(3)$ ist, sondern
    $S^3 \cong SU(2)$, die universelle Ãœberlagerung von $SO(3)$.
    Die Algebra der Quaternionen beschreibt buchstäblich den Spin
    von Elementarteilchen.
\end{beobachtung}

% ============================================================

% ── Erweiterung ──────────────────────────────────────────────────────────
\section{Exponentialabbildung und Lie-Gruppen}

\begin{definition}[Lie-Gruppe und Tangentialraum]
    Eine \textbf{Lie-Gruppe} ist eine Gruppe $G$, die gleichzeitig eine
    glatte Mannigfaltigkeit ist, so dass Gruppenoperation und Inversenbildung
    glatte Abbildungen sind. Der Tangentialraum $T_e G$ im Einselement
    heißt \textbf{Lie-Algebra} $\mathfrak{g}$.

    Die \textbf{Exponentialabbildung} $\exp : \mathfrak{g} \to G$,
    $X \mapsto e^X$, verbindet Lie-Algebra und Lie-Gruppe. Für
    $G = SO(n)$ ist $\mathfrak{g} = \mathfrak{so}(n)$ (schiefsymmetrische
    Matrizen) und $\exp(X) = e^X \in SO(n)$ für alle $X \in \mathfrak{so}(n)$.
\end{definition}

\begin{beispiel}[Drehungen in $\mathbb{R}^3$ durch Exponentialabbildung]
    Jede Drehung $R \in SO(3)$ um die Achse $\hat{n} = (n_1, n_2, n_3)^t$
    (Einheitsvektor) um den Winkel $\theta$ lässt sich als
    Matrixexponential schreiben. Sei
    $[\hat{n}]_\times = \bigl(\begin{smallmatrix}
    0 & -n_3 & n_2 \\ n_3 & 0 & -n_1 \\ -n_2 & n_1 & 0
    \end{smallmatrix}\bigr) \in \mathfrak{so}(3)$
    die schiefsymmetrische Matrix zu $\hat{n}$. Dann gilt die
    \textbf{Rodrigues-Formel}:
    \[
        R = e^{\theta [\hat{n}]_\times}
        = E + \sin\theta \cdot [\hat{n}]_\times
        + (1 - \cos\theta) \cdot [\hat{n}]_\times^2.
    \]
    Der Beweis folgt aus $[\hat{n}]_\times^3 = -[\hat{n}]_\times$
    (für Einheitsvektoren), was die Potenzreihe auf drei Terme reduziert.
\end{beispiel}

\begin{beobachtung}[Quaternionen als Doppelüberlagerung und Spin]
    Die Abbildung $\Phi : S^3 \to SO(3)$ hat Kern $\{e_0, -e_0\}$.
    Dies bedeutet geometrisch: Ein Punkt auf $SO(3)$ (eine Drehung)
    hat genau zwei Urbilder auf $S^3$. Man sagt, $S^3 \cong SU(2)$
    ist die \textbf{universelle Ãœberlagerung} von $SO(3)$.

    In der Quantenmechanik: Spinor-Felder (Elektronen) transformieren
    unter $SU(2)$, nicht unter $SO(3)$. Eine $2\pi$-Drehung ändert
    das Vorzeichen des Spinors (Rotation um $-1$ in $SU(2)$), während
    erst eine $4\pi$-Drehung den Spinor invariant lässt. Dieses
    ``Spin-Halbzahligkeit''-Phänomen ist die physikalische Konsequenz
    der doppelten Ãœberlagerung -- ein direktes Resultat der
    Quaternionenstruktur.
\end{beobachtung}

\chapter{Arithmetik in Integritätsbereichen}
% ============================================================

\section{Teilbarkeit und euklidische Ringe}

Die Arithmetik der ganzen Zahlen und die der Polynome folgen demselben
Muster. Beide sind Instanzen einer einheitlichen algebraischen Theorie:
der Theorie der Integritätsbereiche.

\begin{definition}[Integritätsbereich und Teilbarkeit]
    Ein kommutativer Ring $R$ mit Einselement heißt
    \textbf{Integritätsbereich}, falls aus $ab = 0$ stets $a = 0$
    oder $b = 0$ folgt.

    Für $a, b \in R$ schreibt man $a \mid b$ (``$a$ teilt $b$''),
    falls $b = ra$ für ein $r \in R$. Ein Element $e \in R$ heißt
    \textbf{Einheit}, falls $eR = R$, d.\,h.\ falls $e$ ein
    Inverses besitzt. Zwei Elemente $a, b$ heißen \textbf{assoziiert}
    ($a \sim b$), falls $a = eb$ mit einer Einheit $e$.

    Ein Element $p \in R$ heißt \textbf{irreduzibel}, falls $p \neq 0$,
    $p$ keine Einheit ist, und aus $p = ab$ stets $a$ oder $b$ eine
    Einheit ist. $p$ heißt \textbf{Primelement}, falls $p \mid ab$
    stets $p \mid a$ oder $p \mid b$ impliziert.
\end{definition}

\begin{beobachtung}[Prim vs. irreduzibel]
    In beliebigen Integritätsbereichen gilt: Jedes Primelement ist
    irreduzibel (denn aus $p = ab$ und $p \mid ab$ folgt o.\,B.\,d.\,A.\
    $p \mid a$, also $a = pc$, also $p = pcb$, also $cb = 1$ -- $b$ ist Einheit).
    Die Umkehrung gilt im Allgemeinen \emph{nicht}: In
    $\mathbb{Z}[\sqrt{-5}]$ ist $2$ irreduzibel aber nicht prim, da
    $2 \mid 6 = 2 \cdot 3 = (1 + \sqrt{-5})(1 - \sqrt{-5})$, aber
    $2 \nmid (1 \pm \sqrt{-5})$.

    In Hauptidealringen aber gilt: irreduzibel $\Leftrightarrow$ prim.
    Dies ist der algebraische Kern der eindeutigen Primfaktorzerlegung.
\end{beobachtung}

\begin{definition}[Hauptidealring und euklidischer Ring]
    Ein Integritätsbereich heißt \textbf{Hauptidealring} (HIR), falls jedes
    Ideal ein Hauptideal $Ra = \{ra \mid r \in R\}$ ist.

    Er heißt \textbf{euklidischer Ring}, falls es eine Abbildung
    $\varphi : R \setminus \{0\} \to \mathbb{N}_0$ gibt mit der
    Eigenschaft: Zu jedem Paar $a, b$ mit $a \neq 0$ gibt es $q, r \in R$
    mit $b = qa + r$ und $r = 0$ oder $\varphi(r) < \varphi(a)$
    (\textbf{Division mit Rest}).

    Jeder euklidische Ring ist ein Hauptidealring.
\end{definition}

\begin{satz}[Größter gemeinsamer Teiler in Hauptidealringen]
    Sei $R$ ein Hauptidealring und $r_1, \dots, r_n \in R$.
    Dann existieren $\mathrm{ggT}(r_1, \dots, r_n)$ und
    $\mathrm{kgV}(r_1, \dots, r_n)$. Es gilt:
    \[
        R r_1 + \cdots + R r_n = R \cdot \mathrm{ggT}(r_1, \dots, r_n),
    \]
    und es gibt $s_j \in R$ mit
    $\mathrm{ggT}(r_1, \dots, r_n) = \sum_{j=1}^n s_j r_j$
    (\textbf{Bézout-Darstellung}).
\end{satz}
\begin{proof}
    Das Ideal $I = Rr_1 + \cdots + Rr_n$ ist ein Ideal in $R$.
    Da $R$ ein Hauptidealring ist, gilt $I = Rd$ für ein $d \in R$.
    Aus $d \in I$ folgt $d = \sum_j s_j r_j$ (Bézout-Darstellung).
    Da $r_j \in I = Rd$, gilt $d \mid r_j$ für alle $j$.
    Ist $c$ ein gemeinsamer Teiler aller $r_j$, so $c \mid d$ aus der
    Bézout-Darstellung. Also ist $d = \mathrm{ggT}(r_1, \ldots, r_n)$.
\end{proof}

\begin{beobachtung}[Ganzzahlen und Polynomringe als euklidische Ringe]
    Die ganzen Zahlen $\mathbb{Z}$ mit $\varphi(a) = |a|$ und der
    Polynomring $K[x]$ mit $\varphi(f) = \deg f$ sind euklidische Ringe.
    Ihre Einheitengruppen sind $E(\mathbb{Z}) = \{1, -1\}$ und
    $E(K[x]) = K^*$. Die irreduziblen Elemente in $\mathbb{Z}$ sind
    $\pm p$ für Primzahlen $p$; in $K[x]$ die irreduziblen Polynome.

    Der Euklidische Algorithmus zur Berechnung des $\mathrm{ggT}$
    funktioniert in beiden Ringen vollständig analog: wiederholte
    Division mit Rest, bis der Rest $0$ wird. Dies ist kein Zufall --
    es ist die gemeinsame algebraische Struktur des euklidischen Rings.
\end{beobachtung}

\begin{satz}[Eindeutige Primfaktorzerlegung]
    In jedem Hauptidealring $R$ besitzt jedes Element $r \neq 0$,
    das keine Einheit ist, eine \textbf{Primfaktorzerlegung}
    $r = p_1 \cdots p_k$ in irreduzible Elemente $p_j$. Diese
    Zerlegung ist eindeutig bis auf Reihenfolge und Assoziiertheit
    (d.\,h.\ bis auf Einheitenfaktoren).

    Insbesondere: In $\mathbb{Z}$ ist dies der Fundamentalsatz der
    Arithmetik. In $K[x]$ ist es die eindeutige Zerlegung in
    irreduzible Polynome.
\end{satz}
\begin{proof}
    \textbf{Existenz:} Sei $r$ kein Primelement und keine Einheit. Dann
    $r = ab$ mit $a, b$ Nicht-Einheiten. Da $|R/Rr| < |R/Ra|$ und $|R/Rr| < |R/Rb|$
    (in endlichen Ringen) bzw.\ durch Gradargument in $K[x]$, liefert Induktion
    Primfaktorzerlegungen von $a$ und $b$ und damit von $r$.
    \textbf{Eindeutigkeit:} In Hauptidealringen gilt: irreduzibel $\Rightarrow$ prim
    (da $p \mid ab$ und $p \nmid a$ impliziert $\mathrm{ggT}(p, a) = 1$, also
    $1 = sp + ta$, also $b = spb + tab$ und $p \mid tab$ mit $p \mid tab$).
    Aus der Primeigenschaft folgt die Eindeutigkeit der Zerlegung wie beim klassischen
    Beweis in $\mathbb{Z}$.
\end{proof}

% ============================================================
\section{Irreduzible Polynome}
% ============================================================

\begin{definition}[Irreduzibles Polynom]
    Ein Polynom $f \in K[x]$ mit $\deg f \geq 1$ heißt \textbf{irreduzibel},
    falls es keine Zerlegung $f = gh$ mit $1 \leq \deg g, \deg h < \deg f$ gibt.

    Ãœber $\mathbb{C}$ sind die irreduziblen Polynome genau die Linearfaktoren
    $(x - a)$ (Fundamentalsatz der Algebra). Ãœber $\mathbb{R}$ kommen
    zusätzlich irreduzible quadratische Polynome $x^2 + bx + c$ mit
    $b^2 - 4c < 0$ vor. Ãœber $\mathbb{Q}$ gibt es irreduzible Polynome
    beliebigen Grades.
\end{definition}

\begin{satz}[Erweiterungskörper durch irreduzible Polynome]
    Sei $K$ ein Körper und $f \in K[x]$ irreduzibel. Dann ist
    $L = K[x]/(f)$ ein Körper (der \textbf{Erweiterungskörper})
    mit $[L : K] = \deg f$ und einer Nullstelle $\alpha = x + (f)$ von $f$.
\end{satz}
\begin{proof}
    Da $f$ irreduzibel ist, erzeugt $(f)$ ein maximales Ideal in $K[x]$
    (denn: $(f) \subseteq (g)$ impliziert $g \mid f$, also $g \sim 1$ oder $g \sim f$).
    Ein Faktorring modulo einem maximalen Ideal ist ein Körper.
    Die Dimension $[L:K] = \deg f$ folgt daraus, dass $[1, \alpha, \ldots, \alpha^{n-1}]$
    (mit $n = \deg f$) eine $K$-Träger von $L$ ist.
\end{proof}

\begin{beobachtung}[Ideale, Primideale und der Zusammenhang zur Algebraischen Geometrie]
    Das Ideal $(f)$ in $K[x]$ ist genau dann prim, wenn $f$ irreduzibel ist.
    Der Faktorring $K[x]/(f)$ ist dann ein Integritätsbereich.
    Diese Konstruktion ist das algebraische Gegenstück zur
    \textbf{algebraischen Geometrie}: Die Nullstellen eines Polynoms $f$
    entsprechen dem Primideal $(f)$, und der Erweiterungskörper $K[x]/(f)$
    ist der ``Koordinatenring'' der Nullstellenmenge von $f$.

    Für $f = x^2 + 1$ über $\mathbb{R}$: $L = \mathbb{R}[x]/(x^2+1) \cong \mathbb{C}$.
    Die Konstruktion $\mathbb{C} = \mathbb{R}[x]/(x^2+1)$ ist die präzise
    algebraische Rechtfertigung der komplexen Zahlen: Sie entstehen genau
    als der kleinste Erweiterungskörper, in dem $x^2 + 1$ eine Nullstelle hat.
\end{beobachtung}


\section{Euklidischer Algorithmus und Bézout-Identität}

\begin{beispiel}[Euklidischer Algorithmus in Ganzzahlen und Polynomringen]
    \textbf{In $\mathbb{Z}$:} Berechne $\mathrm{ggT}(252, 105)$:
    \begin{align*}
        252 &= 2 \cdot 105 + 42 \\
        105 &= 2 \cdot 42 + 21 \\
        42  &= 2 \cdot 21 + 0
    \end{align*}
    Also $\mathrm{ggT}(252, 105) = 21$. Rückwärtseinsetzen liefert die
    Bézout-Darstellung: $21 = 105 - 2 \cdot 42 = 105 - 2(252 - 2 \cdot 105)
    = 5 \cdot 105 - 2 \cdot 252$.

    \textbf{In $K[x]$:} $\mathrm{ggT}(x^3 - 1,\; x^2 - 1)$:
    \begin{align*}
        x^3 - 1 &= (x^2 - 1) \cdot x + (x - 1) \\
        x^2 - 1 &= (x-1)(x+1) + 0
    \end{align*}
    Also $\mathrm{ggT}(x^3-1, x^2-1) = x - 1$.
\end{beispiel}


\begin{algorithm}[H]
\DontPrintSemicolon
\caption{Erweiterter Euklidischer Algorithmus}\label{alg:euklid}
\KwIn{$a, b \in R$ mit $a \neq 0$}
\KwOut{$d = \mathrm{ggT}(a,b)$ und $s, t \in R$ mit $d = sa + tb$}
\eIf{$b = 0$}{
    \Return{$(a,\; 1,\; 0)$}\;
}{
    $(q, r) \leftarrow$ Division mit Rest: $a = qb + r$\;
    $(d, s', t') \leftarrow \mathrm{EuklidExt}(b, r)$\;
    \Return{$(d,\; t',\; s' - qt')$}\;
}
\BlankLine
\textbf{Laufzeit:} $O(\log\min(a,b))$ Divisionen \quad
\textbf{Korrektheit:} $d = t'b + s'r = t'b + s'(a - qb) = s'a + (t'-qs')b$
\end{algorithm}

\begin{beobachtung}[RSA und Euklidischer Algorithmus]
    Das RSA-Verfahren (Kapitel 2) nutzt den Euklidischen Algorithmus
    zentral: Um den Entschlüsselungsexponenten $d$ mit $de \equiv 1
    \pmod{\phi(n)}$ zu finden, wendet man den \textbf{erweiterten
    Euklidischen Algorithmus} auf $e$ und $\phi(n)$ an. Die Bézout-Darstellung
    $1 = se + t\phi(n)$ liefert direkt $d = s \bmod \phi(n)$.

    Die Laufzeit des Euklidischen Algorithmus ist $O(\log(\min(a,b)))$
    Divisionsschritte -- weit effizienter als die naive Suche nach
    gemeinsamen Teilern. Dies macht RSA auch für große Zahlen
    ($n \approx 2^{2048}$) praktisch berechenbar.
\end{beobachtung}

\chapter{Hermitesche Abbildungen und Polarzerlegung}
% ============================================================

\section{Hermitesche Abbildungen und Hauptachsen}


\begin{satz}[Spektralsatz für hermitesche Abbildungen]
    Sei $V$ ein Hilbertraum und $A = A^* \in \mathrm{End}(V)$
    hermitesch. Dann gilt:
    \begin{itemize}
        \item[-] Es gibt einen reellen Eigenwert $a$ mit $|a| = \|A\|$.
        \item[-] Alle Eigenwerte von $A$ sind reell.
        \item[-] $V$ besitzt eine Orthonormalträger aus Eigenvektoren von $A$.
        \item[-] Für alle $v \in V$ ist $(Av, v) \in \mathbb{R}$.
    \end{itemize}
    Eine hermitesche Matrix $A \in \mathbb{C}^{n \times n}$ ist unitär
    diagonalisierbar: Es gibt unitäres $U$ mit $U^{-1}AU = \mathrm{diag}(a_1, \dots, a_n)$,
    $a_j \in \mathbb{R}$.
\end{satz}
\begin{proof}
    \textbf{Reelle Eigenwerte:} Ist $Av = av$, so $a\|v\|^2 = (Av, v) = (v, Av) = \bar{a}\|v\|^2$,
    also $a = \bar{a} \in \mathbb{R}$.
    \textbf{Existenz eines reellen Eigenwerts mit $|a| = \|A\|$:}
    Das Maximum von $(Av, v)$ auf der Einheitssphäre wird angenommen
    (Kompaktheit); durch Variation $v \mapsto v + \varepsilon w$ zeigt man,
    dass das Maximum ein Eigenwert ist (Lagrange-Methode).
    \textbf{Orthogonaldiagonalisierbarkeit:} Ist $Av_1 = a_1 v_1$, so ist
    $v_1^\perp$ invariant unter $A$ (da $A = A^*$: $(Aw, v_1) = (w, Av_1) = a_1(w,v_1) = 0$).
    Induktion über $\dim V$ liefert eine ONB aus Eigenvektoren.
\end{proof}


\begin{definition}[Positiv (semi-)definite Abbildungen]
    Eine hermitesche Abbildung $A = A^*$ heißt \textbf{positiv semidefinit}
    ($A \geq 0$), falls $(Av, v) \geq 0$ für alle $v$;
    \textbf{positiv definit} ($A > 0$), falls $(Av, v) > 0$ für alle
    $v \neq 0$.

    Äquivalent: Alle Eigenwerte von $A$ sind $\geq 0$ (bzw. $> 0$).
    Für jedes $A$ ist $A^*A \geq 0$.
\end{definition}

\begin{beobachtung}[Hauptachsentransformation quadratischer Formen]
    Eine quadratische Form $q : \mathbb{R}^n \to \mathbb{R}$,
    $q(x) = x^t A x$ mit symmetrischer Matrix $A$, wird durch
    eine Orthonormalträger aus Eigenvektoren von $A$ auf
    Diagonalgestalt gebracht:
    \[
        q(x) = a_1 y_1^2 + a_2 y_2^2 + \cdots + a_n y_n^2.
    \]
    Dies ist die \textbf{Hauptachsentransformation}: Die Eigenvektoren
    sind die Hauptachsen der Quadrik $\{q(x) = 1\}$, die Eigenwerte
    die reziproken quadratischen Halbachsenlängen.

    Ist $A$ positiv definit, so ist die Quadrik ein Ellipsoid.
    Die Hauptachsentransformation klassifiziert alle reellen Quadriken
    durch ihre Eigenwertstruktur.
\end{beobachtung}

\begin{satz}[Sylvestersches Kriterium]
    Eine hermitesche Matrix $A \in \mathbb{C}^{n \times n}$ ist positiv
    definit genau dann, wenn alle \textbf{Hauptminoren}
    $\delta_r = \det(a_{ij})_{i,j=1}^r > 0$ für $r = 1, \dots, n$.
\end{satz}
\begin{proof}
    \textbf{Positiv definit $\Rightarrow$ alle $\delta_r > 0$:} Die
    Einschränkung $A|_{\langle e_1,\ldots,e_r\rangle}$ ist ebenfalls positiv definit,
    also invertierbar, also $\delta_r = \det(A|_{\langle e_1,\ldots,e_r\rangle}) > 0$.
    \textbf{Alle $\delta_r > 0 \Rightarrow$ positiv definit:} Durch Induktion:
    Für $r = 1$ ist $\delta_1 = a_{11} > 0$. Im Schritt von $r-1$ auf $r$ zeigt man
    mittels der Block-Cholesky-Zerlegung $A = L D L^*$ (mit positiven Diagonaleinträgen
    $D$ aus den Quotienten $\delta_r / \delta_{r-1}$), dass $A$ positiv definit ist.
\end{proof}


% ============================================================
\section{Polarzerlegung und singuläre Wertzerlegung}
% ============================================================

\begin{satz}[Polarzerlegung]
    Sei $V$ ein Hilbertraum und $A \in \mathrm{End}(V)$. Dann gibt es
    ein eindeutiges hermitesches $H \geq 0$ und ein unitäres $U$ mit
    \[
        A = UH, \qquad H^2 = A^*A, \qquad \|A\| = \|H\|.
    \]
    Ist $A$ regulär, so sind $H > 0$ und $U$ beide eindeutig bestimmt.
\end{satz}
\begin{proof}
    Da $A^*A \geq 0$ hermitesch ist, existiert die eindeutige positive Wurzel
    $H = \sqrt{A^*A} \geq 0$. Definiere $U$ auf $\mathrm{Im}\, H$ durch $U(Hv) = Av$;
    dies ist wohldefiniert und isometrisch:
    $\|Av\|^2 = (A^*Av, v) = (H^2 v, v) = (Hv, Hv) = \|Hv\|^2$.
    $U$ erweitert man isometrisch auf $(\mathrm{Im}\, H)^\perp = \ker H = \ker A$.
    Dann ist $U$ unitär und $A = UH$ per Konstruktion.
    Eindeutigkeit für reguläres $A$: Aus $A = UH = U'H'$ folgt
    $A^*A = H^2 = H'^2$, also $H = H'$ (Eindeutigkeit der positiven Wurzel),
    und $U = AH^{-1} = U'$.
\end{proof}


\begin{beobachtung}[Analogie zur komplexen Polardarstellung]
    Die Polarzerlegung $A = UH$ ist das exakte Matrixanalogon zur
    Polardarstellung $z = |z| e^{i\varphi}$ einer komplexen Zahl:
    $H = \sqrt{A^*A}$ ist der ``Betrag'' -- die positive
    Wurzel aus $A^*A$ --, und $U$ ist der ``Phasenfaktor'' --
    eine unitäre Abbildung.

    Geometrisch: Jede invertierbare lineare Abbildung lässt sich
    zerlegen in eine reine Streckung/Stauchung (durch $H$, die
    Richtungen der Eigenvektoren) gefolgt von einer reinen Drehung
    (durch $U$). Dies ist die lineare Algebra der Bewegungskomponenten.
\end{beobachtung}

\begin{satz}[Singuläre Wertzerlegung]
    Sei $A \in \mathbb{C}^{m \times n}$. Dann gibt es unitäre Matrizen
    $U \in \mathbb{C}^{m \times m}$, $V \in \mathbb{C}^{n \times n}$ und
    eine ``Diagonalmatrix'' $\Sigma$ mit nichtnegativen Diagonaleinträgen
    $\sigma_1 \geq \sigma_2 \geq \cdots \geq \sigma_r > 0 = \sigma_{r+1} = \cdots$
    (\textbf{singuläre Werte}) mit
    \[
        A = U \Sigma V^*.
    \]
    Die singulären Werte sind die Quadratwurzeln der Eigenwerte von $A^*A$.
    Es gilt $\mathrm{Rang}\, A = r$ (Anzahl positiver singulärer Werte).
\end{satz}
\begin{proof}
    Da $A^*A \geq 0$ hermitesch ist, existiert eine ONB
    aus Eigenvektoren: $A^*A v_j = \sigma_j^2 v_j$ mit $\sigma_j \geq 0$.
    Für $\sigma_j > 0$ setze $u_j = \frac{1}{\sigma_j} A v_j$; dann gilt
    $(u_i, u_j) = \tfrac{1}{\sigma_i \sigma_j}(Av_i, Av_j) = \tfrac{1}{\sigma_i \sigma_j}
    (v_i, A^*Av_j) = \frac{\sigma_j^2}{\sigma_i \sigma_j} \delta_{ij} = \delta_{ij}$.
    Ergänze $[u_1, \ldots, u_r]$ zu einer ONB $[u_1, \ldots, u_m]$ von $K^m$.
    Dann gilt $A v_j = \sigma_j u_j$ für $j \leq r$ und $A v_j = 0$ für $j > r$.
    Mit $U = (u_1 \cdots u_m)$, $V = (v_1 \cdots v_n)$ und $\Sigma$ wie definiert
    ergibt sich $A = U \Sigma V^*$.
\end{proof}


\begin{beobachtung}[Die singuläre Wertzerlegung als universelles Werkzeug]
    Die singuläre Wertzerlegung (SVD) ist das mächtigste Werkzeug der
    numerischen linearen Algebra. Sie liefert:
    \begin{itemize}
        \item[-] \textbf{Rang und Kondition}: $\mathrm{Rang}\, A = r$,
              Konditionszahl $\kappa(A) = \sigma_1 / \sigma_r$.
        \item[-] \textbf{Beste Approximation}: Unter allen Matrizen
              vom Rang $\leq k$ ist $A_k = U \Sigma_k V^*$ (mit den
              $k$ größten singulären Werten) die beste Approximation
              an $A$ im Sinne der Operatornorm (\textbf{Eckart-Young-Theorem}).
        \item[-] \textbf{Pseudoinverse}: $A^+ = V \Sigma^+ U^*$ löst
              lineare Gleichungssysteme im Sinne der kleinsten Quadrate.
    \end{itemize}
    In der Datenanalyse (PCA), Bildkompression, maschinellem Lernen und
    der numerischen Simulation ist die SVD allgegenwärtig. Sie ist die
    rechnerische Verwirklichung des Spektralsatzes für beliebige
    (nicht-quadratische) Matrizen.
\end{beobachtung}

% ============================================================

\section{Algorithmus zur Singulärwertzerlegung}

\begin{beobachtung}[Bidiagonalisierung und SVD-Berechnung]
    Die numerische Berechnung der SVD erfolgt in zwei Phasen:
    \begin{enumerate}
        \item \textbf{Bidiagonalisierung}: Durch orthogonale
              Transformationen (Householder-Spiege\-lungen) wird $A$
              auf Bidiagonalform $B$ gebracht: $A = U_1 B V_1^t$.
              Aufwand: $\Theta(mn^2)$ für eine $m \times n$-Matrix.
        \item \textbf{SVD der Bidiagonalmatrix}: Auf $B$ wird der
              \textbf{Golub-Reinsch-Algorithmus} angewandt, der
              iterativ die Bidiagonalmatrix durch orthogonale
              Transformationen in Diagonalform überführt.
              Aufwand: $O(n^2)$ pro Iteration, typischerweise
              $O(n)$ Iterationen nötig.
    \end{enumerate}
    Gesamtaufwand: $\Theta(mn^2)$ für $m \geq n$, was für die
    Praxis (z.\,B.\ $m = 10^6$, $n = 10^3$ in maschinellen Lernen)
    sehr effizient ist.
\end{beobachtung}

\begin{beobachtung}[Rangapproximation durch abgeschnittene SVD]
    Die \textbf{Eckart-Young-Theorem} besagt: Die beste Rang-$k$-Approximation
    an $A$ im Sinne der Frobenius-Norm oder der Spektralnorm ist
    $A_k = \sum_{j=1}^k \sigma_j u_j v_j^t$ (abgeschnittene SVD).
    Der Approximationsfehler ist $\|A - A_k\|_2 = \sigma_{k+1}$.

    Anwendungsbeispiel \textbf{Bildkompression}: Ein Graustufenbild
    der Größe $m \times n$ wird als Matrix $A$ gespeichert.
    Die Rang-$k$-Approximation $A_k$ speichert statt $mn$ Zahlen
    nur $k(m + n + 1)$ Zahlen (Singulärwerte und -vektoren).
    Für $k = 50$, $m = n = 1000$: Kompressionsrate $\approx 10:1$
    bei guter visueller Qualität.
\end{beobachtung}

\chapter{Der Satz von Perron--Frobenius}
% ============================================================

\section{Nichtnegative Matrizen}

Nichtnegative Matrizen treten überall auf, wo Übergänge, Flüsse oder
Häufigkeiten modelliert werden: Markov-Ketten, Populationsdynamik,
Netzwerkanalyse. Der Satz von Perron--Frobenius ist das zentrale
Strukturresultat dieser Theorie.

\begin{definition}[Nichtnegative und irreduzible Matrizen]
    Für Matrizen $A = (a_{ij}), B = (b_{ij}) \in \mathbb{R}^{n \times n}$
    schreibt man $A \leq B$ (bzw.\ $A < B$), falls $a_{ij} \leq b_{ij}$
    (bzw.\ $< b_{ij}$) für alle $i, j$.

    Eine Matrix $A \geq 0$ heißt \textbf{irreduzibel}, falls der
    zugehörige gerichtete Graph $\Gamma(A)$ (mit Kante $i \to j$ für
    $a_{ij} > 0$) \textbf{stark zusammenhängend} ist: Je zwei Knoten
    $i, j$ sind durch einen gerichteten Pfad verbunden.

    Äquivalent: Es gibt kein $k$ mit $1 \leq k \leq n-1$ und eine Permutationsmatrix
    $P$, so dass $P^t A P$ Blockdreiecksgestalt
    $\bigl(\begin{smallmatrix} B_{11} & B_{12} \\ 0 & B_{22} \end{smallmatrix}\bigr)$
    hat -- $A$ lässt sich nicht durch Umordnung der Indizes in eine
    obere Blockdreiecksform bringen.
\end{definition}

\begin{definition}[Spektralradius]
    Der \textbf{Spektralradius} einer Matrix $A \in \mathbb{C}^{n \times n}$ ist
    \[
        \rho(A) = \max\{|\lambda| \mid \lambda \text{ Eigenwert von } A\}.
    \]
    Für nichtnegative Matrizen gilt stets $\rho(A) \geq 0$.
\end{definition}

\begin{satz}[Perron--Frobenius]
    Sei $A \in \mathbb{R}^{n \times n}$ irreduzibel und $A \geq 0$. Dann gilt:
    \begin{itemize}
        \item[-] Der \textbf{Spek\-tral\-radius} $\rho(A)$ ist ein Eigenwert
              von $A$ -- der \textbf{Perron--Fro\-benius-Eigen\-wert}.
        \item[-] Es gibt einen Eigenvektor $v > 0$ mit $Av = \rho(A) v$.
              Dieser ist bis auf positive Skalierung eindeutig.
        \item[-] $\rho(A)$ ist ein einfacher Eigenwert.
        \item[-] Ist $Aw = bw$ mit $0 \neq w \geq 0$, so ist $b = \rho(A)$.
    \end{itemize}
\end{satz}
\begin{proof}
    Sei $K = \{x \in \mathbb{R}^n \mid x \geq 0,\, \|x\|_1 = 1\}$ der
    Einheitssimplex. Die Abbildung $T : K \to K$, $T(x) = Ax/\|Ax\|_1$,
    ist stetig und $K$ kompakt-konvex. Nach dem \textbf{Brouwerschen Fixpunktsatz}
    gibt es $v \in K$ mit $T(v) = v$, also $Av = \rho(A) v$ mit $\rho(A) = \|Av\|_1$.
    \textbf{Positivität von $v$:} Da $A$ irreduzibel ist, gibt es für jedes $i$
    einen Pfad $i \to j$ in $\Gamma(A)$, also $(A^k)_{ji} > 0$ für ein $k$.
    Aus $v_i > 0$ für ein $i$ folgt durch Iteration $v > 0$.
    \textbf{Eindeutigkeit:} Sei $Aw = \rho(A) w$ mit $w \geq 0$.
    Dann gilt $w = cv$ für $c = \min_j(w_j/v_j)$; der Vektor $w - cv \geq 0$
    ist Eigenvektor, hat aber eine Null-Komponente, was mit $w - cv > 0$
    (aus Irreduzibilität) nur für $w - cv = 0$ vereinbar ist.
    \textbf{Einfachheit} folgt aus der Eindeutigkeit des positiven Eigenvektors.
\end{proof}

\begin{beobachtung}[Perron--Frobenius und der PageRank-Algorithmus]
    Der Satz von Perron--Frobenius liegt dem \textbf{PageRank}-Algorithmus
    von Google zugrunde: Sei $A$ die (geeignet normierte) Inzidenzmatrix des
    Webs. Der Wichtigkeitsvektor der Webseiten ist der positive Eigenvektor
    $v > 0$ zum Eigenwert $\rho(A) = 1$.

    Der Perron--Frobenius-Satz garantiert \emph{Existenz} und
    \emph{Eindeutigkeit} eines solchen Vektors -- und damit die
    Wohldefiniertheit des Rankings. Ohne diesen Satz gäbe es keine
    mathematische Garantie für ein konsistentes Suchmaschinen-Ranking.
\end{beobachtung}

% ============================================================
\section{Stochastische Matrizen}
% ============================================================

\begin{definition}[Stochastische Matrix]
    Eine Matrix $A \geq 0$ heißt \textbf{(zeilenstochastisch)}, falls
    alle Zeilensummen gleich $1$ sind: $\sum_j a_{ij} = 1$ für alle $i$.
    Äquivalent: Der Einheitsvektor $(1, \ldots, 1)^t$ ist Eigenvektor
    zum Eigenwert $1$.

    Eine stochastische Matrix ist irreduzibel genau dann, wenn der zugehörige
    Markov-Graph stark zusammenhängend ist: Von jedem Zustand aus
    ist jeder andere Zustand erreichbar.
\end{definition}

\begin{satz}[Konvergenz irreduzibeler Markov-Ketten]
    Sei $A$ eine irreduzible stochastische Matrix. Dann gilt:
    \begin{itemize}
        \item[-] $\rho(A) = 1$ ist einfacher Eigenwert.
        \item[-] Es gibt eine eindeutige \textbf{stationäre Verteilung}
              $\pi > 0$ mit $A^t \pi = \pi$ und $\sum_i \pi_i = 1$.
        \item[-] Für jede Startverteilung $x_0 \geq 0$ mit $\sum_i (x_0)_i = 1$
              gilt $\lim_{k \to \infty} (A^t)^k x_0 = \pi$.
    \end{itemize}
\end{satz}
\begin{proof}
    Da $\rho(A) = 1$ (alle Zeilensummen $=1$), liefert Perron--Frobenius
    die eindeutige positive stationäre Verteilung $\pi$.
    Die Konvergenz $\lim_k (A^t)^k x_0 = \pi$ folgt daraus, dass alle anderen
    Eigenwerte von $A$ Betrag $< 1$ haben (aus der Irreduzibilität und der
    Einfachheit von $\rho = 1$): Die entsprechenden Eigenraumanteile von
    $x_0$ klingen mit geometrischer Rate ab.
\end{proof}

\begin{beobachtung}[Perron--Frobenius und stochastische Prozesse]
    Die stationäre Verteilung $\pi$ ist der ``Gleich\-ge\-wichts\-zustand''
    der Markov-Kette: Unabhängig vom Startzustand konvergiert das System
    gegen $\pi$.

    Konkrete Anwendungen: Das \textbf{MCMC-Verfahren} (Markov Chain Monte Carlo)
    nutzt genau diese Konvergenz zum Sampling aus komplizierten Verteilungen.
    In der \textbf{Populationsgenetik} beschreiben irreduzible Markov-Ketten
    die Genotypenhäufigkeiten in großen Populationen.
    In der \textbf{Physik} entsprechen stationäre Verteilungen dem
    thermodynamischen Gleichgewicht detailbalancierter Systeme (Boltzmann-Verteilung).

    Die mathematische Tiefe liegt in der Kombination zweier Prinzipien:
    Perron--Frobenius garantiert die Eindeutigkeit des Fixpunkts,
    und die Eigenwertabschätzung $|\lambda_j| < 1$ für $j \geq 2$
    garantiert die Konvergenz -- mit exponentieller Rate $|\lambda_2|^k$.
\end{beobachtung}

% ============================================================

% ── Erweiterung ──────────────────────────────────────────────────────────
\section{Spektralradius und Normabschätzungen}

\begin{satz}[Gelfand-Formel]
    Für jede Matrix $A \in \mathbb{C}^{n \times n}$ gilt:
    \[
        \rho(A) = \lim_{k \to \infty} \|A^k\|^{1/k}
    \]
    für jede submultiplikative Matrixnorm $\|\cdot\|$. Insbesondere
    ist $\rho(A) \leq \|A\|$ für alle solchen Normen.
\end{satz}
\begin{proof}
    Da alle Eigenwerte $|\lambda_j| \leq \|A\|$, gilt $\rho(A) \leq \|A\|$.
    Die Jordan-Normalform liefert für jedes $\varepsilon > 0$:
    $\|A^k\| \leq C (\rho(A) + \varepsilon)^k$ für eine Konstante $C$.
    Damit $\limsup_k \|A^k\|^{1/k} \leq \rho(A) + \varepsilon$.
    Da $\varepsilon > 0$ beliebig und $\|A^k\|^{1/k} \geq \rho(A)$
    (aus dem Spektralhalbkreis), folgt $\lim_k \|A^k\|^{1/k} = \rho(A)$.
\end{proof}

\begin{beobachtung}[Stabilität linearer Systeme]
    Der Spektralradius entscheidet über die Langzeitstabilität eines
    linearen dynamischen Systems $x_{k+1} = Ax_k$:
    \begin{itemize}
        \item[-] $\rho(A) < 1$: Das System ist \textbf{stabil} --
              $\|A^k x_0\| \to 0$ für $k \to \infty$.
        \item[-] $\rho(A) = 1$: \textbf{Grenzstabilität} --
              Abhängig von der Jordanstruktur zu den Eigenwerten
              auf dem Einheitskreis.
        \item[-] $\rho(A) > 1$: Das System ist \textbf{instabil} --
              $\|A^k x_0\| \to \infty$ für generische Anfangsbedingungen.
    \end{itemize}
    Für kontinuierliche Systeme $\dot{x} = Ax$: Stabilität genau
    dann, wenn alle Eigenwerte negativen Realteil haben.
    Der Satz von Perron--Frobenius garantiert für irreduzible
    nichtnegative Matrizen, dass $\rho(A)$ ein einfacher Eigenwert ist --
    das stochastische System ist daher genau grenzstabil.
\end{beobachtung}

\chapter{Endliche Untergruppen von $SO(3)$}
% ============================================================

\section{Endliche Drehgruppen}

Die Frage, welche endlichen Gruppen als Untergruppen von $SO(3)$ auftreten,
beantwortet ein vollständiger Klassifikationssatz. Die Antwort ist
die algebraische Erklärung eines zweieinhalbtausend Jahre alten geometrischen Phänomens.

\begin{definition}[Drehungen und Drehachsen]
    Jedes $1 \neq \sigma \in SO(3)$ ist eine Drehung um einen eindeutig
    bestimmten Einheitsvektor (die \textbf{Drehachse}) um einen Winkel
    $\varphi \in (0, 2\pi)$. Eine Achse heißt \textbf{$n$-zählig},
    falls die zugehörigen Drehungen um die Winkel $\frac{2\pi k}{n}$
    ($k = 1, \ldots, n$) alle in $G$ liegen.

    Sei $G \leq SO(3)$ endlich mit $|G| = g > 1$.
    Die Menge $\Omega \subset S^2$ der Fixpunkte aller Drehachsen
    (Schnittpunkte mit der Einheitssphäre) trägt eine natürliche
    $G$-Wirkung; ihre Bahnen sind die \textbf{Pol-Bahnen}.
\end{definition}

\begin{satz}[Klassifikation nach Bahnen und Ordnungen]
    Sei $G$ eine endliche Untergruppe von $SO(3)$ mit $|G| = g > 1$.
    Seien $K_1, \dots, K_k$ die Bahnen konjugierter Achsen von $G$,
    die Achsen aus $K_j$ seien $n_j$-zählig ($n_j \geq 2$).
    Dann gilt die \textbf{Bahnformel}:
    \[
        2(g - 1) = \sum_{j=1}^k g\!\left(1 - \tfrac{1}{n_j}\right).
    \]
    Diese Gleichung hat -- für $k \leq 3$ und $2 \leq n_1 \leq n_2 \leq n_3$ --
    genau fünf Lösungstypen:
    \begin{enumerate}
        \item[(1)] $k = 2$, $n_1 = n_2 = g$: zyklische Gruppe $\mathbb{Z}_g$
        \item[(2)] $k = 3$, $n_1 = n_2 = 2$, $n_3 = g/2$: Diedergruppe $D_g$
        \item[(3)] $k = 3$, $n_1 = 2$, $n_2 = n_3 = 3$, $g = 12$: $A_4$
        \item[(4)] $k = 3$, $n_1 = 2$, $n_2 = 3$, $n_3 = 4$, $g = 24$: $S_4$
        \item[(5)] $k = 3$, $n_1 = 2$, $n_2 = 3$, $n_3 = 5$, $g = 60$: $A_5$
    \end{enumerate}
\end{satz}
\begin{proof}
    Jedes $1 \neq \sigma \in G$ ist eine Drehung mit einer Achse $\langle a \rangle$.
    Sei $\Omega$ die Menge der Fixpunkte aller solchen Achsen auf $S^2$.
    Zähle die Paare $(\sigma, p)$ mit $\sigma \neq 1$ und $\sigma(p) = p$:
    Jedes Nicht-Eins-Element hat genau zwei Fixpunkte auf $S^2$; jeder Fixpunkt
    $p$ hat $|\mathrm{Stab}(p)| - 1$ solche Elemente. Damit:
    \[
        2(g-1) = \sum_{p \in \Omega} (|\mathrm{Stab}(p)| - 1) = \sum_j |B_j|(n_j - 1)
        = \sum_j g(1 - 1/n_j),
    \]
    wobei $|B_j| = g/n_j$ die Bahnlänge der $j$-ten Achsen-Bahn ist.
    Die Auflösung dieser diophantischen Gleichung mit $n_j \geq 2$ liefert
    genau fünf Typen.
\end{proof}

% ============================================================
\section{Die fünf platonischen Körper}
% ============================================================

\begin{satz}[Platonische Körper und Symmetriegruppen]
    Die fünf Typen endlicher Untergruppen von $SO(3)$ entsprechen
    genau den Symmetriegruppen der \textbf{platonischen Körper}:
    \begin{itemize}
        \item[-] $\mathbb{Z}_n$: Symmetriegruppe der regelmäßigen $n$-seitigen Pyramide
        \item[-] $D_n$: Symmetriegruppe des regelmäßigen $n$-seitigen Prismas
        \item[-] $A_4 \cong$ Symmetriegruppe des \textbf{regulären Tetraeders} ($|G| = 12$)
        \item[-] $S_4 \cong$ Symmetriegruppe des \textbf{Würfels/Oktaeders} ($|G| = 24$)
        \item[-] $A_5 \cong$ Symmetriegruppe des \textbf{Dodekaeders/Ikosaeders} ($|G| = 60$)
    \end{itemize}
\end{satz}
\begin{proof}
    Der Beweis folgt aus der Klassifikation nach Bahnen und Ordnungen.
    \textbf{$A_4$ und Tetraeder:} Ein reguläres Tetraeder hat 4 Ecken, 6 Kanten, 4 Flächen.
    Die Ecken-Drehachsen sind 3-zählig ($n=3$), die Kanten-Mitte-Achsen 2-zählig ($n=2$).
    Mit $k=3$, $n_1=2$, $n_2=n_3=3$ liefert die Bahnformel $g=12$, und die
    Symmetriegruppe ist $A_4$ (die geraden Permutationen der 4 Ecken).
    \textbf{$S_4$ und Würfel/Oktaeder:} Würfel und Oktaeder sind dual; ihre
    Symmetriegruppe permutiert die 4 Raumdiagonalen des Würfels, also $g=24$, Typ $S_4$.
    \textbf{$A_5$ und Dodekaeder/Ikosaeder:} Dual zueinander; die 5-zähligen
    Achsen und $g=60$ erzwingen den Typ $A_5$. Die bijektive Zuordnung
    folgt direkt aus den Parametern $(k, n_1, n_2, n_3)$ der Bahnformel.
\end{proof}

\begin{beobachtung}[Die Eulerschen Charakteristik der platonischen Körper]
    Für jeden konvexen Polyeder gilt die \textbf{Eulersche Formel}:
    \[
        E - K + F = 2,
    \]
    wobei $E$ die Anzahl der Ecken, $K$ die der Kanten und $F$ die
    der Flächen ist. Für die fünf platonischen Körper:

    \begin{center}
    \begin{tabular}{lrrrrr}
        \hline
        Körper & $E$ & $K$ & $F$ & $E - K + F$ & Symmetriegruppe \\
        \hline
        Tetraeder  & 4  & 6  & 4  & 2 & $A_4$, $|G|=12$ \\
        Würfel     & 8  & 12 & 6  & 2 & $S_4$, $|G|=24$ \\
        Oktaeder   & 6  & 12 & 8  & 2 & $S_4$, $|G|=24$ \\
        Dodekaeder & 20 & 30 & 12 & 2 & $A_5$, $|G|=60$ \\
        Ikosaeder  & 12 & 30 & 20 & 2 & $A_5$, $|G|=60$ \\
        \hline
    \end{tabular}
    \end{center}

    Würfel und Oktaeder sowie Dodekaeder und Ikosaeder sind
    \textbf{duale Polyeder}: Man erhält das Duale, indem man Mittelpunkte der
    Flächen als neue Ecken nimmt. Die Dualität erklärt, warum duale Paare
    dieselbe Symmetriegruppe haben.
\end{beobachtung}

\begin{beobachtung}[Warum es genau fünf platonische Körper gibt]
    Der Klassifikationssatz beantwortet eine $2500$ Jahre alte Frage:
    Warum gibt es genau fünf platonische Körper?

    Die Antwort ist rein algebraisch: Es gibt genau fünf Typen endlicher
    Untergruppen von $SO(3)$ -- und jeder Typ entspricht dem
    Symmetrieprinzip eines platonischen Körpers. Die Unmöglichkeit
    weiterer platonischer Körper ist keine geometrische Tatsache,
    sondern eine Konsequenz der Gruppenstruktur: Die Bahnformel
    $2(g-1) = \sum_j g(1 - 1/n_j)$ hat genau fünf Lösungstypen.

    Euklid wusste um die fünf Körper; dass es keine weiteren gibt,
    erforderte die Algebra des 19.\ Jahrhunderts.
\end{beobachtung}

\begin{beobachtung}[Die alternierende Gruppe $A_5$ und das Ikosaeder]
    Die Gruppe $A_5$ der geraden Permutationen auf fünf Elementen --
    mit $|A_5| = 60$ -- ist die Symmetriegruppe des Ikosaeders.
    Sie ist die kleinste nichtabelsche einfache Gruppe (d.\,h.\ sie hat
    keine echten Normalteiler außer $\{e\}$ und sich selbst).

    $A_5$ taucht in zwei scheinbar völlig verschiedenen Kontexten auf:
    als Symmetriegruppe eines geometrischen Körpers und als Galoisgruppe
    der allgemeinen Gleichung fünften Grades. Abel und Galois zeigten,
    dass Gleichungen fünften Grades nicht durch Radikale lösbar sind --
    genau weil $A_5$ einfach ist. Die Geometrie des Ikosaeders und die
    Unlösbarkeit der Quintic sind tief verbunden durch dieselbe Gruppe.
\end{beobachtung}

% ============================================================

% ── Erweiterung ──────────────────────────────────────────────────────────
\section{Duale Körper und Kristallographie}

\begin{beobachtung}[Kristallographische Einschränkung]
    In der Kristallographie fragt man: Welche Drehsymmetrien sind mit
    einem periodischen Gitter verträglich? Die Antwort ist bekannt
    als die \textbf{kristallographische Einschränkung}: In zwei und
    drei Dimensionen sind nur $n$-zählige Drehachsen mit $n \in \{1,2,3,4,6\}$
    mit Gitterperiodizität verträglich.

    Beweis: Eine $n$-zählige Drehung $R$ muss Gitterpunkte auf
    Gitterpunkte abbilden. Die Spur $\mathrm{tr}(R) = 1 + 2\cos(2\pi/n)$
    muss eine ganze Zahl sein (da $R$ in der Trägerbasis von $\mathbb{Z}^2$
    dargestellt werden kann). Also $2\cos(2\pi/n) \in \mathbb{Z}$,
    was genau die Werte $n = 1, 2, 3, 4, 6$ erlaubt.

    Die $5$-zählige Symmetrie -- das Fünfeck -- ist mit periodischen
    Gittern unverträglich. Die Entdeckung von \textbf{Quasikristallen}
    mit $5$-zähliger Symmetrie (Shechtman, 1982, Nobelpreis 2011)
    revolutionierte die Kristallographie: Quasikristalle sind
    aperiodisch geordnete Strukturen mit verbotener Symmetrie.
\end{beobachtung}

\begin{beobachtung}[Platonische Körper in der Chemie und Physik]
    Die fünf platonischen Körper treten in der Natur und Wissenschaft auf:
    \begin{itemize}
        \item[-] \textbf{Tetraeder ($A_4$)}: Grundstruktur des
              Methan-Moleküls $\mathrm{CH}_4$; Kohlenstoff-Tetraeder
              in Diamant und Silizium.
        \item[-] \textbf{Würfel ($S_4$)}: Grundgitter des $\mathrm{NaCl}$-Kristalls
              (Natriumchlorid, Kochsalz).
        \item[-] \textbf{Oktaeder ($S_4$)}: Ligandanordnung in
              oktaedrischen Ãœbergangsmetallkomplexen.
        \item[-] \textbf{Ikosaeder ($A_5$)}: Struktur vieler Viren
              (z.\,B.\ Adenovirus); Buckminster-Fulleren $\mathrm{C}_{60}$
              hat ikosaedrische Symmetriegruppe (abgestumpftes Ikosaeder).
        \item[-] \textbf{Dodekaeder ($A_5$)}: In der kosmologischen
              Hypothese von Luminet et al.\ (2003) hat der Raum möglicherweise
              Dodekaeder-Topologie.
    \end{itemize}
\end{beobachtung}

\chapter{Moduln und abelsche Gruppen}
% ============================================================

\section{Moduln: Vektorräume über Ringen}

Der Modulbegriff vereinheitlicht Vektorräume und abelsche Gruppen
unter einem gemeinsamen Dach. Ãœber einem Hauptidealring gibt es
einen vollständigen Klassifikationssatz.

\begin{definition}[Modul und freier Modul]
    Sei $R$ ein Ring mit Einselement. Eine abelsche Gruppe $(M, +)$
    heißt \textbf{$R$-Modul} (Linksmodul), falls eine Skalarmultiplikation
    $R \times M \to M$ definiert ist, so dass für alle $r, r_j \in R$
    und $m, m_j \in M$ gilt:
    \[
        r(m_1 + m_2) = rm_1 + rm_2, \quad
        (r_1 + r_2)m = r_1 m + r_2 m, \quad
        (r_1 r_2)m = r_1(r_2 m), \quad
        1m = m.
    \]
    $M$ heißt \textbf{freier $R$-Modul} mit freien Erzeugenden
    $f_1, \dots, f_k$, falls $M = \bigoplus_{j=1}^k R f_j$ mit
    $Rf_j \cong R$. Dann heißt $k = \mathrm{Rg}\, M$ der \textbf{Rang}
    von $M$.
\end{definition}

\begin{beobachtung}[Moduln als gemeinsamer Rahmen]
    Der Modulbegriff vereinheitlicht mehrere fundamentale Strukturen:
    \begin{itemize}
        \item[-] Jeder $K$-Vektorraum ist ein $K$-Modul über einem Körper $K$.
        \item[-] Die $\mathbb{Z}$-Moduln sind genau die abelschen Gruppen.
        \item[-] Ist $V$ ein $K$-Vektorraum und $A \in \mathrm{End}(V)$,
              so wird $V$ durch $f(x) v = f(A)v$ zu einem $K[x]$-Modul.
              Die Untermoduln sind genau die $A$-invarianten Unterräume.
    \end{itemize}
    Das dritte Beispiel ist der Schlüssel zur Jordanschen Normalform:
    Ein Endomorphismus $A$ auf $V$ macht $V$ zu einem $K[x]$-Modul,
    und die Klassifikation der $K[x]$-Moduln -- der Hauptsatz über
    Moduln über Hauptidealringen -- liefert direkt die Jordanform.
\end{beobachtung}

\begin{definition}[Torsionsmodul und Torsionselement]
    Sei $R$ ein Integritätsbereich und $M$ ein $R$-Modul. Ein Element
    $m \in M$ heißt \textbf{Torsionselement}, falls es ein $0 \neq r \in R$
    gibt mit $rm = 0$. Die Menge $T(M)$ aller Torsionselemente ist ein
    Untermodul (der \textbf{Torsionsmodul}), und $M/T(M)$ ist
    torsionsfrei.

    Für $R = \mathbb{Z}$: Torsionselemente einer abelschen Gruppe sind
    die Elemente endlicher Ordnung. Für $R = K[x]$: Torsionselemente
    des $K[x]$-Moduls $V$ sind alle Vektoren (da $\chi_A(A) = 0$).
\end{definition}

\begin{satz}[Hauptsatz über endlich erzeugte Moduln über Hauptidealringen]
    Sei $R$ ein Hauptidealring und $M$ ein endlich erzeugbarer $R$-Modul.
    Dann gilt die \textbf{kanonische Zerlegung}:
    \[
        M \cong R^k \oplus \bigoplus_{j} R/R p_j^{e_j},
    \]
    wobei $k = \mathrm{Rg}(M/T(M))$ der Rang des torsionsfreien Anteils
    ist und $p_j$ Primelemente von $R$ sind. Der Rang $k$ und die
    Primpotenzen $p_j^{e_j}$ (die \textbf{Elementarteiler}) sind
    eindeutig bestimmt.
\end{satz}
\begin{proof}[Beweisskizze]
    \textbf{Existenz} der Zerlegung: Wähle Erzeuger $m_1, \ldots, m_k$ von $M$
    und betrachte den freien Modul $R^k$. Der Kern $K$ der Abbildung
    $R^k \to M$, $(r_i) \mapsto \sum r_i m_i$, ist ebenfalls endlich erzeugt
    (da $R$ Hauptidealring, also noethersch). Durch simultane elementare Zeilen-
    und Spaltenoperationen auf der Darstellungsmatrix von $K \hookrightarrow R^k$
    (Smith-Normalform) kann man einen Träger $(f_1, \ldots, f_k)$ von $R^k$ und
    $(e_1, \ldots, e_k)$ von $K$ finden mit $e_i = d_i f_i$ und $d_i \mid d_{i+1}$.
    Damit $M \cong \bigoplus_{i} R/Rd_i$.
    \textbf{Eindeutigkeit} der Elementarteiler: Die Anzahl der $d_i$ mit $d_i \neq 0$
    bestimmt den Rang; die Elementarteiler selbst sind als $\mathrm{ggT}$
    aller $k \times k$-Minoren der Darstellungsmatrix invariant.
\end{proof}

% ============================================================
\section{Abelsche Gruppen und Jordanform}
% ============================================================

\begin{satz}[Klassifikation endlicher abelscher Gruppen]
    Jede endlich erzeugte abelsche Gruppe $A$ besitzt eine direkte
    Zerlegung
    \[
        A = \langle a_1 \rangle \oplus \cdots \oplus \langle a_n \rangle
    \]
    in zyklische Gruppen, wobei $\mathrm{Ord}(a_i) \in \{p^e \mid p
    \text{ Primzahl}, e \geq 1\} \cup \{\infty\}$. Die Ordnungen sind
    bis auf Reihenfolge eindeutig bestimmt.
\end{satz}
\begin{proof}
    Jede endlich erzeugte abelsche Gruppe ist ein $\mathbb{Z}$-Modul.
    Der Hauptsatz über Moduln über Hauptidealringen (oben) mit $R = \mathbb{Z}$
    liefert die Zerlegung $A \cong \mathbb{Z}^k \oplus \bigoplus_j \mathbb{Z}/d_j\mathbb{Z}$
    mit $d_j \mid d_{j+1}$ und $d_j \geq 2$.
    Jede zyklische Gruppe $\mathbb{Z}/d\mathbb{Z}$ mit $d = \prod_i p_i^{e_i}$
    zerfällt nach dem Chinesischen Restsatz in $\bigoplus_i \mathbb{Z}/p_i^{e_i}\mathbb{Z}$.
    Dies liefert die Zerlegung in Primärkomponenten.
    \textbf{Eindeutigkeit:} Die Primärkomponenten $A_p = \{a \in A \mid p^k a = 0\}$
    sind durch $A$ eindeutig bestimmt (als charakteristische Untergruppen),
    und die Elementarteiler von $A_p$ sind die Invarianten des $\mathbb{Z}_p$-Moduls $A_p$.
\end{proof}

\begin{beobachtung}[Die Elementarteilertheorie als Klassifikationsmotor]
    Der Hauptsatz über Moduln über Hauptidealringen ist einer der
    mächtigsten Klassifikationssätze der Algebra -- er klassifiziert
    gleichzeitig:
    \begin{itemize}
        \item[-] \textbf{Endliche abelsche Gruppen} (Spezialfall $R = \mathbb{Z}$):
              Jede endliche abelsche Gruppe ist eindeutig bestimmt
              durch ihre Primärkomponenten. Zum Beispiel:
              Die abelschen Gruppen der Ordnung $12 = 2^2 \cdot 3$ sind
              $\mathbb{Z}_{12}$, $\mathbb{Z}_4 \times \mathbb{Z}_3$,
              $\mathbb{Z}_2 \times \mathbb{Z}_2 \times \mathbb{Z}_3$ -- genau drei.
        \item[-] \textbf{Lineare Abbildungen bis auf Ähnlichkeit}
              (Spezialfall $R = K[x]$): Die Elementarteiler von $xE - A$
              sind genau die Invarianzfaktoren, aus denen die Jordanform
              abgelesen wird.
        \item[-] \textbf{Gitterstrukturen} in Zahlentheorie und Geometrie
              (Spezialfall $R = \mathbb{Z}$, freier Modul): Jedes
              Gitter $\Lambda \leq \mathbb{Z}^n$ ist isomorph zu einem
              Standardgitter, charakterisiert durch seine Elementarteiler.
    \end{itemize}
    Ein einziger Satz -- drei fundamentale Klassifikationen.
    Das ist die Tiefe der Modultheorie.
\end{beobachtung}

\begin{beobachtung}[Smith-Normalform als algorithmisches Werkzeug]
    Die \textbf{Smith-Normalform} einer Matrix $A \in R^{m \times n}$
    über einem Hauptidealring $R$ ist die Diagonalmatrix
    $D = \mathrm{diag}(d_1, \ldots, d_r, 0, \ldots, 0)$ mit
    $d_j \mid d_{j+1}$, die durch elementare Zeilen- und Spaltenoperationen
    aus $A$ hervorgeht.

    Die Existenz der Smith-Normalform ist genau der Hauptsatz -- denn
    $A$ ist die Darstellungsmatrix des Moduls $M = R^n / \mathrm{Im}(A)$,
    und die Diagonaleinträge $d_j$ sind die Elementarteiler von $M$.

    Für $R = \mathbb{Z}$: Die Smith-Normalform ganzzahliger Matrizen
    ist ein grundlegendes Werkzeug der algebraischen Topologie
    (Berechnung von Homologiegruppen), der Zahlentheorie (Gittertheorie)
    und der Optimierung (ganzzahlige Programmierung).
\end{beobachtung}


% ── Erweiterung ──────────────────────────────────────────────────────────
\section{Noethersche Ringe und endliche Erzeugbarkeit}

\begin{definition}[Noetherscher Ring]
    Ein Ring $R$ heißt \textbf{noethersch}, falls jede aufsteigende
    Kette von Idealen $I_1 \subseteq I_2 \subseteq \cdots$ stationär
    wird: Es gibt $n_0$ mit $I_n = I_{n_0}$ für alle $n \geq n_0$.
\end{definition}

\begin{satz}[Hilbertscher Trägersatz]
    Ist $R$ ein noetherscher Ring, so ist auch der Polynomring $R[x]$
    noethersch. Insbesondere sind $\mathbb{Z}[x_1, \ldots, x_n]$ und
    $K[x_1, \ldots, x_n]$ (für jeden Körper $K$) noethersche Ringe.
\end{satz}
\begin{proof}[Beweisskizze]
    Sei $I \subseteq R[x]$ ein Ideal. Betrachte das Ideal
    $I_k = \{a \in R \mid \exists f \in I : \deg f = k, \text{ Leitkoeffizient } a\}
    \cup \{0\}$. Da $R$ noethersch, wird die aufsteigende Kette
    $I_0 \subseteq I_1 \subseteq \cdots$ stationär bei $I_{n_0}$.
    Dann erzeugen endlich viele Polynome aus $I$ (mit Graden $\leq n_0$)
    das gesamte Ideal $I$.
\end{proof}

\begin{beobachtung}[Verbindung zur Smith-Normalform]
    Die Smith-Normalform ist das algorithmische Herzstück des Hauptsatzes
    über Moduln. Über $R = \mathbb{Z}$ wird sie durch den
    erweiterten Euklidischen Algorithmus berechnet:
    \begin{enumerate}
        \item Wähle den Eintrag mit kleinstem $|\cdot|$ als Pivotelement.
        \item Eliminiere alle anderen Einträge in Zeile und Spalte
              durch elementare Operationen.
        \item Wiederhole auf der $(n-1)\times(m-1)$-Untermatrix.
    \end{enumerate}
    Die Laufzeit ist $O(n^3 \log(\max|a_{ij}|))$ für eine $n \times n$-Matrix.
    Die entstehenden Diagonaleinträge $d_1 \mid d_2 \mid \cdots \mid d_r$
    sind die \textbf{Invarianzfaktoren} des Moduls.
\end{beobachtung}

\chapter{Lineare Codes und der Golay-Code}
% ============================================================

\section{Grundbegriffe}

Die Codierungstheorie stellt die Frage: Wie können Nachrichten so kodiert
werden, dass Übertragungsfehler erkannt und korrigiert werden können?
Die Antwort gibt die Lineare Algebra über endlichen Körpern.

\begin{definition}[Linearer Code]
    Sei $K = \mathbb{F}_q$ ein endlicher Körper. Ein \textbf{linearer
    $[n, k]$-Code} über $K$ ist ein $k$-dimensionaler Unterraum
    $C \leq K^n$. Die Elemente heißen \textbf{Codeworte},
    $n$ die \textbf{Länge} und $k$ die \textbf{Dimension}.
    Die \textbf{Informationsrate} ist $R = k/n$.

    Der \textbf{Hamming-Abstand} zweier Codeworte ist
    $d(c, c') = |\{j \mid c_j \neq c'_j\}| = \mathrm{wt}(c - c')$,
    wobei $\mathrm{wt}(v)$ das \textbf{Hamming-Gewicht} (Anzahl der
    Nichtnull-Einträge) ist.
    Die \textbf{Minimaldistanz} ist $d = \min_{0 \neq c \in C} \mathrm{wt}(c)$.
    Man spricht von einem $[n, k, d]$-Code.
\end{definition}

\begin{definition}[Erzeuger- und Kontrollmatrix]
    Eine $k \times n$-Matrix $G$ mit $C = \{uG \mid u \in K^k\}$
    heißt \textbf{Generatormatrix} von $C$.

    Eine $(n-k) \times n$-Matrix $H$ mit $C = \ker H = \{x \in K^n \mid Hx^t = 0\}$
    heißt \textbf{Kontrollmatrix} (parity-check matrix).
    Ist $G$ in Standardform $G = (E_k \mid A)$, so ist $H = (-A^t \mid E_{n-k})$.
\end{definition}

\begin{beobachtung}[Fehlerkorrektur durch Minimalabstand]
    Ein $[n,k,d]$-Code kann
    \begin{itemize}
        \item[-] bis zu $d-1$ Fehler \textbf{erkennen}: Liegt das empfangene Wort
              nicht in $C$, so ist ein Fehler aufgetreten.
        \item[-] bis zu $t = \lfloor(d-1)/2\rfloor$ Fehler \textbf{korrigieren}:
              Das nächste Codewort (im Hamming-Abstand) ist eindeutig das gesendete.
    \end{itemize}
    Das \textbf{Syndrom} $s = Hy^t$ eines empfangenen Wortes $y = c + e$
    ist $He^t$ -- es hängt nur vom Fehlervektor $e$ ab, nicht vom
    gesendeten Codewort $c$. Der Dekodieralgorithmus sucht den
    wahrscheinlichsten Fehlervektor mit diesem Syndrom.
\end{beobachtung}

\begin{definition}[Dualcode]
    Zum Code $C \leq K^n$ mit dem Standard-Skalarprodukt $(c, c') =
    \sum_j c_j c'_j$ ist der \textbf{Dualcode}
    \[
        C^\perp = \{v \in K^n \mid (c, v) = 0\ \forall c \in C\}.
    \]
    Es gilt $\dim C^\perp = n - k$ und $(C^\perp)^\perp = C$.
    $C$ heißt \textbf{selbstdual}, falls $C = C^\perp$.
\end{definition}

\begin{satz}[Dualitätssatz von MacWilliams]
    Sei $C$ ein $[n, k]$-Code über $K$ mit $|K| = q$ und
    Gewichtspolynom $A(x) = \sum_{j=0}^n A_j x^j$. Dann hat
    $C^\perp$ das Gewichtspolynom
    \[
        B(x) = q^{-k}(1 + (q-1)x)^n\, A\!\left(\frac{1-x}{1+(q-1)x}\right).
    \]
\end{satz}
\begin{proof}[Beweisskizze]
    Definiere die \textbf{Hamming-Gewichtsenumerator} über Charaktersummen:
    \[
        B(x) = q^{-k} \sum_{v \in K^n} \hat{1}_C(v) x^{\mathrm{wt}(v)},
    \]
    wobei $\hat{1}_C(v) = \sum_{c \in C} \chi(v \cdot c)$ die Fouriertransformation
    über $(K^n, +)$ der Indikatorfunktion von $C$ ist. Anwenden der diskreten
    Fouriertransformation und Auswerten der Charaktersummen
    $\sum_{c \in C} \chi(\langle v, c \rangle)$ -- die für $v \in C^\perp$
    den Wert $|C|$ und sonst $0$ ergibt -- liefert die angegebene Formel.
\end{proof}

% ============================================================
\section{Schranken und Hamming-Code}
% ============================================================

\begin{satz}[Kugelpackungsschranke (Hamming-Schranke)]
    Ein $[n, k, d]$-Code über $\mathbb{F}_q$ mit $t = \lfloor(d-1)/2\rfloor$
    erfüllt:
    \[
        q^k \cdot \sum_{j=0}^t \binom{n}{j}(q-1)^j \leq q^n.
    \]
    Ein Code, für den Gleichheit gilt, heißt \textbf{perfekt}.
\end{satz}
\begin{proof}
    Die Hamming-Kugeln $B_t(c) = \{y \in K^n \mid d(y,c) \leq t\}$
    um je zwei verschiedene Codeworte $c, c'$ sind disjunkt (da $d \geq 2t+1$
    impliziert $B_t(c) \cap B_t(c') = \emptyset$). Jede Kugel enthält
    $\sum_{j=0}^t \binom{n}{j}(q-1)^j$ Elemente.
    Da alle Kugeln in $K^n$ enthalten sind, folgt die Schranke.
\end{proof}

\begin{beispiel}[Hamming-Code $\mathrm{Ham}(r, q)$]
    Der binäre \textbf{Hamming-Code} $\mathrm{Ham}(r, 2)$ hat Parameter
    $[n, k, d] = [2^r - 1,\, 2^r - r - 1,\, 3]$ und ist ein perfekter
    1-Fehler-korrigierender Code. Seine Kontrollmatrix $H$ hat als Spalten
    alle $2^r - 1$ Nichtnull-Vektoren aus $\mathbb{F}_2^r$.

    Für $r = 3$: $[7, 4, 3]$-Code. Das Syndrom $He^t$ eines
    Ein-Fehler-Musters $e = \delta_{ij}$ ist genau die $j$-te Spalte von $H$,
    also die Binärdarstellung von $j$ -- der Fehlerort wird direkt aus dem
    Syndrom abgelesen.
\end{beispiel}

% ============================================================
\section{Der Golay-Code}
% ============================================================

\begin{beispiel}[Der Golay-Code]
    Der \textbf{binäre Golay-Code} $\mathrm{Gol}(23)$ ist ein
    perfekter $[23, 12, 7]$-Code über $\mathbb{F}_2$: Er korrigiert
    bis zu $3$ Fehler. Die Kugelpackungsgleichung
    $2^{23} = 2^{12} \cdot \sum_{j=0}^3 \binom{23}{j}$ ist erfüllt.

    Der \textbf{erweiterte Golay-Code} $\mathrm{Gol}(24)$ ist ein
    selbstdualer $[24, 12, 8]$-Code, der durch Anhängen einer Prüfstelle
    entsteht. Seine Automorphismengruppe ist die Mathieu-Gruppe $M_{24}$
    mit $|M_{24}| = 244\,823\,040$ -- eine der 26 sporadischen einfachen
    Gruppen.
\end{beispiel}

\begin{beobachtung}[Codes zwischen Algebra und Anwendung]
    Lineare Codes verbinden drei mathematische Gebiete:
    \begin{itemize}
        \item[-] \textbf{Lineare Algebra über $\mathbb{F}_q$}: Unterräume,
              duale Räume, Trägertransformationen liefern die algebraische Struktur.
        \item[-] \textbf{Kombinatorik}: Die Hamming-Schranke und das
              Gewichtspolynom verbinden die algebraische mit der kombinatorischen Struktur.
        \item[-] \textbf{Gruppen\-theorie}: Die Auto\-morphismen\-gruppen perfekter Codes
              sind oft außer\-ge\-wöhn\-liche Gruppen
              (Mathieu-Gruppen, sporadische Gruppen).
    \end{itemize}

    Technische Anwendungen: Der Golay-Code wurde 1977 von der Raumsonde
    Voyager zur Ãœbertragung von Farbbildern des Jupiters verwendet.
    Reed-Solomon-Codes (eine Verallgemeinerung des Hamming-Codes über
    $\mathbb{F}_{2^m}$) werden in CDs, DVDs, QR-Codes und der
    Tief\-raum-Kommunikation eingesetzt.
\end{beobachtung}

% ============================================================

% ── Erweiterung ──────────────────────────────────────────────────────────
\section{Reed-Solomon-Codes}

\begin{definition}[Reed-Solomon-Code]
    Sei $q = p^m$ eine Primzahlpotenz, $K = \mathbb{F}_q$ und
    $\alpha_1, \ldots, \alpha_n \in K$ paarweise verschieden.
    Der \textbf{Reed-Solomon-Code} $\mathrm{RS}(n, k)$ ist:
    \[
        C = \{(f(\alpha_1), \ldots, f(\alpha_n)) \mid f \in K[x],\,\deg f < k\}.
    \]
    Er ist ein $[n, k, n-k+1]$-Code über $K$ -- er erreicht die
    \textbf{Singleton-Schranke} $d \leq n - k + 1$ mit Gleichheit
    (\textbf{MDS-Code}, Maximum Distance Separable).
\end{definition}

\begin{beobachtung}[Reed-Solomon in der Praxis]
    Reed-Solomon-Codes sind die am weitesten verbreiteten Fehlerkorrekturcodes:
    \begin{itemize}
        \item[-] \textbf{CDs und DVDs}: $\mathrm{RS}(255, 223)$ über
              $\mathbb{F}_{2^8}$ korrigiert bis zu 16 Bytfehler pro Block.
        \item[-] \textbf{QR-Codes}: Verschiedene Reed-Solomon-Level
              (7\,\%, 15\,\%, 25\,\%, 30\,\% Redundanz).
        \item[-] \textbf{RAID-6}: Zwei unabhängige Reed-Solomon-Prüfsummen
              erlauben Rekonstruktion bei Ausfall von zwei Festplatten.
        \item[-] \textbf{Deep-Space-Kommunikation}: Die Raumsonde
              Voyager nutzte einen verketteten Code aus einem
              Golay-Code und einem Reed-Solomon-Code.
    \end{itemize}
    Die algebraische Struktur -- Auswertung von Polynomen über
    endlichen Körpern -- macht Reed-Solomon-Codes effizient
    kodier- und dekodierbar (Berlekamp-Welch-Algorithmus in $O(n^2)$).
\end{beobachtung}

\chapter{Lineare Schwingungen}
% ============================================================

\section{Schwingungen ohne Reibung}

\begin{definition}[Schwingungsgleichung]
    Seien $n$ Massenpunkte mit Massen $m_1, \dots, m_n > 0$ durch
    Hooksche Federn (Federkonstante $c > 0$) miteinander verbunden.
    Die Auslenkungen $x_j(t)$ aus der Gleichgewichtslage genügen dem
    System
    \[
        M x''(t) = -A x(t),
    \]
    wobei $M = \mathrm{diag}(m_1, \dots, m_n)$ die \textbf{Massenmatrix}
    und $A = A^*$ die \textbf{Steifigkeitsmatrix} mit $A \geq 0$ ist.
\end{definition}

\begin{satz}[Eigenfrequenzen und Normalschwingungen]
    Sei $M > 0$ und $A \geq 0$. Da $M^{-1/2} A M^{-1/2}$ hermitesch
    und $\geq 0$ ist, hat $M^{-1}A$ reelle, nichtnegative Eigenwerte
    $0 \leq \omega_1^2 \leq \omega_2^2 \leq \cdots \leq \omega_n^2$.
    Die allgemeine Lösung ist:
    \[
        x(t) = \sum_{l=0}^{r} (c_{l1} e^{i\omega_l t} + c_{l2} e^{-i\omega_l t})
        v_l + (\text{pol. Anteil für } \omega_l = 0),
    \]
    wobei $v_l$ die Eigenvektoren zu den Eigenfrequenzen $\omega_l$ sind.
    Die speziellen Lösungen $x(t) = (\cos \omega_l t) v_l$ heißen
    \textbf{Normalschwingungen}.
\end{satz}
\begin{proof}
    Da $M > 0$, ist $M^{1/2}$ wohldefiniert und invertierbar. Der Ansatz
    $x(t) = e^{i\omega t} v$ in $Mx'' = -Ax$ ergibt $-\omega^2 M v = -Av$,
    also das \textbf{verallgemeinerte Eigenwertproblem} $Av = \omega^2 M v$.
    Durch Substitution $w = M^{1/2} v$ wird dies zu $(M^{-1/2} A M^{-1/2}) w = \omega^2 w$,
    einem gewöhnlichen hermiteschen Eigenwertproblem.
    Da $M^{-1/2} A M^{-1/2}$ hermitesch und $\geq 0$ ist, liefert der Spektralsatz
    reelle nichtnegative Eigenwerte $\omega_j^2 \geq 0$ und eine ONB aus
    Eigenvektoren $w_j$. Rücktransformation $v_j = M^{-1/2} w_j$ liefert
    die Normalschwingungen; die allgemeine Lösung ergibt sich durch Superposition.
\end{proof}


\begin{beispiel}[Gekoppelte Pendel auf der $x$-Achse]
    Für $n$ gleichschwere Massen $m$ mit gleichen Federn $c$ zwischen
    festen Endpunkten $0$ und $L$ ergibt sich die tridiagonale
    Steifigkeitsmatrix $A = 2E - B$ mit der Shift-Matrix $B$.
    Die Eigenfrequenzen und Eigenvektoren sind explizit berechenbar:
    \[
        \omega_j^2 = \frac{4c}{m} \sin^2\!\frac{j\pi}{2(n+1)},
        \qquad
        (v_j)_k = \sin\frac{jk\pi}{n+1} \quad (j, k = 1, \dots, n).
    \]
    Dies sind stehende Wellen: Die $j$-te Normalschwingung hat exakt
    $j$ Knoten. Das Diskretes-Sinus-Transform (DST) der Signalverarbeitung
    ist die Diagonalisierungstransformation dieser Matrix.
\end{beispiel}

\begin{beobachtung}[Entkopplung durch Eigenvektoren]
    Die Normalschwingungen sind entkoppelt: In den Eigenkoordinaten
    $y = T^{-1}x$ (wobei $T$ die Eigenvektoren enthält) zerfällt das
    System $M x'' = -Ax$ in $n$ unabhängige harmonische Oszillatoren
    $y_j'' = -\omega_j^2 y_j$.

    Physikalisch: Die Normalschwingungen sind die ``natürlichen''
    Schwingungsmoden des Systems. Jede beliebige Bewegung ist eine
    Superposition von Normalschwingungen -- das Prinzip der
    Modensuperposition. Die Eigenwertzerlegung ist die mathematische
    Formulierung der akustischen Resonanz.
\end{beobachtung}

% ============================================================
\section{Schwingungen mit Reibung und Resonanz}
% ============================================================

\begin{definition}[Gedämpfte Schwingungsgleichung]
    Wirkt zusätzlich eine geschwindigkeitsproportionale Dämpfung, so
    lautet das System:
    \[
        M x'' + D x' + A x = f(t),
    \]
    wobei $D \geq 0$ die \textbf{Dämpfungsmatrix} und $f(t)$ eine
    äußere Anregung ist.
\end{definition}

\begin{satz}[Struktursatz gedämpfter Schwingungen]
    Für proportionale Dämpfung $D = \alpha M + \beta A$ ($\alpha, \beta \geq 0$)
    lässt sich das System durch dieselbe Transformation, die $M^{-1}A$
    diagonalisiert, vollständig entkoppeln. Die $j$-te Komponente genügt:
    \[
        y_j'' + \delta_j y_j' + \omega_j^2 y_j = g_j(t),
        \qquad \delta_j = \alpha + \beta \omega_j^2.
    \]
    Für $\delta_j < 2\omega_j$ (\textbf{Unterdämpfung}) schwingt das
    System mit der \textbf{Eigenfrequenz} $\tilde{\omega}_j =
    \sqrt{\omega_j^2 - \delta_j^2/4}$ mit abklingender Amplitude.
    Für $\delta_j \geq 2\omega_j$ (\textbf{Überdämpfung}) klingt die
    Bewegung ohne Schwingung ab.
\end{satz}
\begin{proof}
    Sei $T$ die Matrix, die $M^{-1}A$ diagonalisiert: $T^{-1}(M^{-1}A)T = \mathrm{diag}(\omega_j^2)$.
    Mit $y = T^{-1} M^{-1/2} x$ und proportionaler Dämpfung $D = \alpha M + \beta A$
    gilt $M^{-1}D = \alpha E + \beta M^{-1}A$, also $T^{-1}(M^{-1}D)T = \mathrm{diag}(\delta_j)$
    mit $\delta_j = \alpha + \beta \omega_j^2$.
    Das transformierte System zerfällt in $n$ skalare Gleichungen
    $y_j'' + \delta_j y_j' + \omega_j^2 y_j = g_j(t)$.
    Für Unterdämpfung ($\delta_j < 2\omega_j$) hat das charakteristische Polynom
    $\lambda^2 + \delta_j \lambda + \omega_j^2$ die Wurzeln
    $\lambda_{1,2} = -\delta_j/2 \pm i\tilde{\omega}_j$ mit
    $\tilde{\omega}_j = \sqrt{\omega_j^2 - \delta_j^2/4} > 0$.
\end{proof}


\begin{beobachtung}[Resonanz und ihre Folgen]
    Für harmonische Anregung $f(t) = f_0 e^{i\omega t}$ wächst die
    Amplitude im Fall der Resonanz ($\omega = \omega_j$, undämpft)
    \emph{linear} mit der Zeit: $|y_j(t)| \sim t$. Dies ist das
    Phänomen der \textbf{Resonanzkatastrophe} -- die Brücke von Tacoma
    Narrows (1940) kollabierte, weil die Windfrequenz mit einer
    Eigenfrequenz der Brücke zusammenfiel.

    Bei positiver Dämpfung $\delta_j > 0$ bleibt die stationäre
    Amplitude beschränkt:
    \[
        |y_j^\mathrm{st}(\omega)| = \frac{|g_j|}{\sqrt{(\omega_j^2 - \omega^2)^2
        + \delta_j^2 \omega^2}}.
    \]
    Das Maximum wird bei $\omega = \sqrt{\omega_j^2 - \delta_j^2/2}$
    angenommen -- der \textbf{Resonanzgipfel} liegt unterhalb der
    Eigenfrequenz. Spektroskopie, Kernspintomographie (MRI) und
    Lasertechnik beruhen auf dem präzisen Ausnutzen von Resonanzgipfeln.
\end{beobachtung}

\begin{beobachtung}[Normalschwingungen als Schwerpunkt der Ingenieurdynamik]
    Die Eigenwertzerlegung der Schwingungsgleichung ist der Kern der
    Strukturdynamik: Brücken, Gebäude, Motoren, Flugzeugflügel --
    alle werden auf ihre Eigenfrequenzen analysiert, um Resonanz mit
    äußeren Anregungen zu vermeiden.

    Die Mathematik ist stets dieselbe: Das Eigenwertproblem $Av = \lambda Mv$
    (das \textbf{verallgemeinerte Eigenwertproblem}) bestimmt die
    Eigenfrequenzen. Die Entkopplung durch Eigenvektoren verwandelt
    ein System von $n$ gekoppelten Differentialgleichungen in $n$
    unabhängige -- die Grundleistung der linearen Algebra für die
    Dynamik kontinuierlicher und diskreter Systeme.
\end{beobachtung}

% ============================================================

% ── Erweiterung ──────────────────────────────────────────────────────────
\section{Der Schwerpunkt schwingender Systeme}

\begin{beobachtung}[Normalschwingungen als Träger des Schwingungsraums]
    Die Normalschwingungen $v_1, \ldots, v_n$ eines schwingenden
    Systems mit $n$ Freiheitsgraden bilden einen Träger des
    Konfigurationsraums $\mathbb{R}^n$ im Sinne dieser Arbeit:
    Sie sind linear unabhängig (da Eigenvektoren zu verschiedenen
    Eigenfrequenzen $\omega_j^2$) und erzeugen den gesamten Raum.

    Dies ist der \textbf{Schwerpunkt des Schwingungssystems}: In den
    Normalkoordinaten $y_j = v_j^t x$ entkoppeln die Bewegungsgleichungen
    vollständig. Jede Normalschwingung $y_j$ erfüllt den einfachen
    Harmonischen Oszillator $y_j'' + \omega_j^2 y_j = 0$, ohne Kopplung
    an die anderen Freiheitsgrade. Das System ist maximal effizient
    lösbar -- genau das Schwerpunkt-Prinzip aus Kapitel~6.
\end{beobachtung}

\begin{beispiel}[Zwei gekoppelte Pendel]
    Betrachte zwei Pendel der Masse $m$ und Länge $\ell$, verbunden
    durch eine Feder der Konstanten $k$. Die Bewegungsgleichungen sind:
    \[
        m\ddot{x}_1 = -\frac{mg}{\ell}x_1 - k(x_1 - x_2), \qquad
        m\ddot{x}_2 = -\frac{mg}{\ell}x_2 + k(x_1 - x_2).
    \]
    Die Steifigkeitsmatrix lautet
    $A = \bigl(\begin{smallmatrix} \omega_0^2 + \kappa & -\kappa \\
    -\kappa & \omega_0^2 + \kappa \end{smallmatrix}\bigr)$
    mit $\omega_0^2 = g/\ell$ und $\kappa = k/m$.

    Die Eigenfrequenzen sind $\omega_1^2 = \omega_0^2$ und
    $\omega_2^2 = \omega_0^2 + 2\kappa$, die zugehörigen
    Normalschwingungen:
    \begin{itemize}
        \item[-] $v_1 = (1,1)^t/\sqrt{2}$: Gleichphasige Schwingung
              (beide Pendel in Gleichschritt, Feder ohne Wirkung).
        \item[-] $v_2 = (1,-1)^t/\sqrt{2}$: Gegenphasige Schwingung
              (Pendel in Gegenphase, Feder maximal gedehnt).
    \end{itemize}
    Die allgemeine Lösung ist die Superposition:
    $x(t) = c_1 \cos(\omega_1 t + \phi_1) v_1 + c_2 \cos(\omega_2 t + \phi_2) v_2$.
    Das Phänomen der \textbf{Schwebung} tritt auf, wenn $\omega_1 \approx \omega_2$:
    Die Schwingungsenergie pendelt mit der Schwebungsfrequenz
    $(\omega_2 - \omega_1)/2$ periodisch zwischen den beiden Pendeln hin und her.
\end{beispiel}

\begin{beobachtung}[Resonanz und die Grenzen des Schwerpunkts]
    Der Schwerpunkt eines Schwingungssystems -- die Normalmodenzerlegung --
    ist nur im linearen, ungedämpften Fall exakt. In der Praxis gibt es
    zwei wesentliche Abweichungen:
    \begin{itemize}
        \item[-] \textbf{Nichtlinearität}: Große Auslenkungen erzeugen
              nichtlineare Terme in den Bewegungsgleichungen
              (z.\,B.\ $\sin\varphi \approx \varphi - \varphi^3/6$
              für ein Pendel). Die Eigenfrequenzen werden
              amplitudenabhängig -- der Schwerpunkt verschiebt sich.
        \item[-] \textbf{Dämpfung und Resonanz}: Proportionale Dämpfung
              erhält den Schwerpunkt (da $D = \alpha M + \beta A$
              dieselben Eigenvektoren hat wie $M^{-1}A$). Allgemeine
              Dämpfung koppelt die Normalmoden und zerstört den Schwerpunkt.
    \end{itemize}
    Die Resonanzkatastrophe (Tacoma Narrows, 1940; Millennium Bridge, 2000)
    ist ein Beispiel, wo das reale System durch Windlasten in eine
    Normalmode angeregt wurde, die im linearen Modell als Schwerpunkt
    erscheint -- die nichtlinearen Effekte zerstörten dann das System.
\end{beobachtung}

\chapter{Zusammenfassung und Ausblick}
% ============================================================

Diese Arbeit hat die Lineare Algebra von ihren mengentheoretischen
Grundlagen über die zentralen Struktursätze bis hin zu tiefen
Verbindungen mit Physik, Kryptographie, Codierungstheorie und
Gruppentheorie entwickelt. Im abschließenden Kapitel wird der
zurückgelegte Weg zusammengefasst und ein Ausblick auf weiterführende
Themen gegeben, die organisch aus dem Dargestellten erwachsen.

% ============================================================
\section{Der rote Faden der Arbeit}
% ============================================================

\subsection*{Von Mengen zu Strukturen}

Der Einstieg in Kapitel 1 mit Mengen, Abbildungen und dem Binomischen
Lehrsatz war bewusst abstrakt gehalten: Mengen und Abbildungen bilden
die Sprache, in der alle algebraischen Strukturen formuliert werden.
Die Grundbegriffe der Bijektivität, der Äquivalenzrelation und des
Inklusions-Exklusions-Prinzips kehrten in verfeinerter Form in
jedem späteren Kapitel wieder -- als Isomorphismen, Kongruenzklassen
und Siebformeln.

Kapitel 2 führte die algebraischen Hauptstrukturen ein: Gruppen,
Ringe und Körper. Der Vektorraum als $K$-Modul über einem Körper $K$
ist die zentrale Struktur der gesamten Arbeit. Die Definition des
Trägers -- einer linear unabhängigen Menge, die den gesamten
Vektorraum erzeugt -- ist der erste strukturelle Höhepunkt:
Der Steinitzsche Austauschsatz zeigt, dass alle Träger eines
endlichdimensionalen Vektorraums dieselbe Kardinalität haben,
und diese Kardinalität heißt Dimension.

\subsection*{Lineare Abbildungen als Herzstück}

Kapitel 3 legte den konzeptuellen Kern frei: Lineare Abbildungen
sind die strukturerhaltenden Abbildungen zwischen Vektorräumen,
und ihre vollständige Beschreibung durch Matrizen (nach Wahl eines
Trägers) verwandelt abstrakte Algebra in konkretes Rechnen.
Der \textbf{Rangsatz} $\dim\ker f + \dim\mathrm{Im}\,f = \dim V$
ist das zentrale Strukturresultat der linearen Algebra -- er sagt,
dass der ``Informationsverlust'' einer linearen Abbildung (ihr Kern)
und ihr ``Informationsgehalt'' (ihr Bild) sich stets zur Dimension
des Definitionsbereichs ergänzen.

Matrizen wurden als Koordinatendarstellungen linearer Abbildungen
verstanden. Determinanten erschienen zunächst als Volumenmaß,
dann als Gruppencharakter von $GL(n,K)$, dann als alternierende
Multilinearform -- drei Perspektiven, die sich im Multiplikationssatz
$\det(AB) = \det A \cdot \det B$ vereinen.

\subsection*{Eigenwerte und Normalformen}

Die Kapitel 4 und 9 entfalteten die Theorie der Normalformen:
Das charakteristische Polynom $\chi_A$ und das Minimalpolynom $m_A$
beschreiben die Struktur eines Endomorphismus vollständig.
Der Satz von Cayley--Hamilton -- jede Matrix ist Nullstelle ihres
eigenen charakteristischen Polynoms -- ist das tiefste Ergebnis
dieser Theorie und hat unmittelbare Konsequenzen für
Matrixinverse, Potenzberechnung und Differentialgleichungen.

Die \textbf{Jordansche Normalform} ist der vollständige
Struktursatz für Endomorphismen über algebraisch abgeschlossenen
Körpern: Jede Matrix ist ähnlich zu einer Blockdiagonalmatrix
aus Jordankästchen. Die Größen der Kästchen sind invariant und
aus den Rängen der Potenzen $(A - \lambda E)^k$ ablesbar.
Die Verbindung zur Primärkomponentenzerlegung und zum Chinesischen
Restsatz zeigt, dass die Jordanform nicht ein Hilfsmittel der
Linearen Algebra ist, sondern ein Spezialfall der
allgemeinen Modultheorie über Hauptidealringen.

\subsection*{Skalarprodukte und Spektraltheorie}

Die Einführung von Skalarprodukten in den Kapiteln 10 und 13
verwandelte Vektorräume in metrische Räume. Das
\textbf{Gram-Schmidt-Verfahren} zeigt, dass jeder
endlichdimensionale Hilbertraum einen Orthonormalträger hat;
die Orthogonalprojektion liefert die beste Approximation
eines Vektors aus einem Unterraum.

Der \textbf{Spektralsatz für normale Abbildungen} ist der
Höhepunkt der metrischen linearen Algebra: Normale Abbildungen
($AA^* = A^*A$) sind genau diejenigen, die bezüglich eines
Orthonormalträgers diagonalisierbar sind. Für selbstadjungierte
Abbildungen ($A = A^*$) sind alle Eigenwerte reell -- was sie
zu den Observablen der Quantenmechanik macht. Für unitäre
Abbildungen ($AA^* = E$) liegen alle Eigenwerte auf dem
Einheitskreis -- was sie zu Zeitentwicklungsoperatoren macht.

Die \textbf{Singuläre Wertzerlegung} verallgemeinerte den
Spektralsatz auf beliebige (nicht-quadratische) Matrizen:
$A = U\Sigma V^*$ mit unitären $U, V$ und diagonalem $\Sigma$.
Die singulären Werte $\sigma_j$ sind die Quadratwurzeln der
Eigenwerte von $A^*A$ und messen die ``Streckungsstärken'' der
Abbildung in den einzelnen Richtungen.

\subsection*{Algebra, Zahlentheorie und Geometrie}

Die Kapitel 6, 8, 15 und 18 vertieften den Zusammenhang zwischen
Linearer Algebra und abstrakten algebraischen Strukturen.
Endliche Körper $\mathbb{F}_{p^n}$ -- vollständig klassifiziert
durch Primzahlpotenzen -- sind die natürlichen Skalarkörper
der Codierungstheorie. Der Frobenius-Automorphismus $k \mapsto k^p$
ist das einfachste Beispiel eines nichtidentischen Feldautomorphismus
und das Vorbild aller Galois-Theorie.

Der \textbf{Satz von Jordan--Hölder} und die Klassifikation der
endlichen einfachen Gruppen zeigen, dass jede endliche Gruppe
eindeutig aus einfachen Gruppen ``aufgebaut'' ist -- analog zur
Primfaktorzerlegung. Die \textbf{Klassifikation der endlichen
Untergruppen von $SO(3)$} durch die Bahnformel
$2(g-1) = \sum_j g(1 - 1/n_j)$ erklärt das zweieinhalbtausend
Jahre alte Phänomen der genau fünf platonischen Körper mit den
Mitteln der Gruppentheorie des 19.~Jahrhunderts.

\subsection*{Anwendungen: Physik, Technik, Informatik}

Die Querverbindungen zur Physik und Technik durchzogen die gesamte
Arbeit. Der \textbf{Minkowskiraum} mit Signatur $(3,1)$ ist der
geometrische Rahmen der Speziellen Relativitätstheorie;
Lorentz-Transformationen sind seine Isometrien, und das
Einsteinsche Additionsgesetz ist die Gruppenstruktur der Boosts.
Die \textbf{linearen Schwingungen} führten auf das verallgemeinerte
Eigenwertproblem $Av = \omega^2 Mv$ und zeigten, dass die
Eigenfrequenzen physikalischer Systeme (Brücken, Motoren, Moleküle)
durch Diagonalisierung der Steifigkeits- und Massenmatrix bestimmt
werden -- Resonanzkatastrophen entstehen, wenn Anregungsfrequenzen
mit Eigenfrequenzen zusammenfallen.

Der \textbf{Satz von Perron--Frobenius} und irreduzible stochastische
Matrizen bilden das mathematische Fundament des PageRank-Algorithmus,
der MCMC-Verfahren und der Populationsdynamik. Die \textbf{Codierungstheorie}
verbindet Unterraumtheorie über $\mathbb{F}_q$, kombinatorische
Schranken und Gruppentheorie zu einem Werkzeug, das jede moderne
Datenkommunikation und -speicherung trägt.

% ============================================================
\section{Weiterführende Gebiete}
% ============================================================

Die in dieser Arbeit entwickelte Lineare Algebra ist ein
Ausgangspunkt, von dem aus sich zahlreiche mathematische
Horizonte öffnen. Im Folgenden werden die wichtigsten
Fortsetzungsrichtungen skizziert.

\subsection*{Unendlichdimensionale lineare Algebra: Funktionalanalysis}

Der natürliche Übergang von endlicher zu unendlicher Dimension
führt in die \textbf{Funktionalanalysis}. Hilberträume unendlicher
Dimension -- etwa der Raum $L^2([0,1])$ der quadratintegrierbaren
Funktionen -- sind der Rahmen der modernen Analysis, der
Quantenmechanik und der partiellen Differentialgleichungen.

Viele Sätze der endlichdimensionalen Theorie übertragen sich,
aber mit wesentlichen Unterschieden:
\begin{itemize}
    \item[-] Der Spek\-tral\-satz für selbst\-adjungierte Operatoren auf
          un\-end\-lich\-di\-men\-sio\-na\-len Hilbert\-räumen erfordert das
          \textbf{Spektralmaß} (von-Neumann-Spek\-tral\-satz): Das
          Spektrum ist keine endliche Menge von Eigenwerten mehr,
          sondern eine Teilmenge von $\mathbb{R}$, die auch
          kontinuierliche und Restspektrum-Anteile enthalten kann.
    \item[-] Die \textbf{Kompakten Operatoren} sind das un\-end\-lich\-di\-men\-sio\-nale
          Analogon der end\-lich\-di\-men\-sio\-na\-len Abbildungen: Für sie gilt
          der Spektralsatz mit abzählbarem Spektrum und $0$ als
          einzigem Häufungspunkt. Kompakte selbst\-adjungierte Operatoren
          sind das Herzstück der \textbf{Sturm--Liouville-Theorie}
          und der Fourier-Analysis.
    \item[-] \textbf{Banachräume} -- vollständige normierte Räume
          ohne Skalarprodukt -- bilden den weiteren Rahmen.
          Der \textbf{Hahn--Banach-Satz} (Fortsetzung linearer
          Funktionale), der \textbf{Satz vom abgeschlossenen Graphen}
          und der \textbf{Satz von der gleichmäßigen Beschränktheit}
          (Banach--Steinhaus) sind die drei Grundpfeiler der
          Funktionalanalysis, ohne Analogon in endlicher Dimension.
\end{itemize}

\subsection*{Multilineare Algebra und Tensorrechnung}

Die \textbf{äußere Algebra} $\bigwedge V$, in dieser Arbeit als
Erklärungsrahmen für die Determinante eingeführt, ist ein Spezialfall
der \textbf{Tensoralgebra} $T(V) = \bigoplus_{k \geq 0} V^{\otimes k}$.
Das Tensorprodukt $V \otimes W$ zweier Vektorräume hat die universelle
Eigenschaft, bilineare Abbildungen in lineare zu linearisieren.

Tensorrechnung ist die Sprache der modernen Physik:
\begin{itemize}
    \item[-] Die \textbf{allgemeine Relativitätstheorie} (Einstein, 1915)
          ist in der Sprache der \textbf{Riemannschen Geometrie}
          formuliert, deren Grundobjekt der metrische Tensor $g_{ij}$
          ist -- ein symmetrisches $(0,2)$-Tensorfeld auf einer
          Mannigfaltigkeit. Die Einstein-Feldgleichungen
          $R_{\mu\nu} - \frac{1}{2}Rg_{\mu\nu} + \Lambda g_{\mu\nu}
          = \frac{8\pi G}{c^4} T_{\mu\nu}$ sind Gleichungen zwischen
          Tensoren.
    \item[-] In der \textbf{Quantenfeldtheorie} beschreiben
          Tensoren die Zustände von Mehrteilchensystemen: Der
          Zustand eines Systems aus zwei Teilchen lebt im
          Tensorprodukt $\mathcal{H}_1 \otimes \mathcal{H}_2$
          der Einteilchen-Hilberträume. \textbf{Verschränkung}
          (Entanglement) ist das Phänomen, dass ein Zustand
          \emph{nicht} als Produkt $v_1 \otimes v_2$ geschrieben
          werden kann.
    \item[-] \textbf{Tensornetze} und \textbf{Tensorzerlegungen}
          (Tucker-Zerlegung, CP-Zerlegung) sind aktive
          Forschungsgebiete der angewandten Mathematik und des
          maschinellen Lernens.
\end{itemize}

\subsection*{Algebraische Geometrie}

Die in Kapitel 15 hergestellte Verbindung zwischen irreduziblen
Polynomen und Erweiterungskörpern ist der Einstieg in die
\textbf{algebraische Geometrie}. Ihr Grundobjekt ist die
\textbf{affine Varietät}
$V(f_1, \ldots, f_r) = \{x \in K^n \mid f_1(x) = \cdots = f_r(x) = 0\}$,
die gemeinsame Nullstellenmenge von Polynomen.

Der \textbf{Hilbertsche Nullstellensatz} (Hilbert, 1893) ist das
fundamentale Ergebnis: Ãœber algebraisch abgeschlossenem $K$ gibt es
eine vollständige Korrespondenz zwischen geometrischen Objekten
(Varietäten) und algebraischen Objekten (radikalen Idealen im
Polynomring). Dies ist die algebraische Geometrie im engeren Sinne.

Moderne algebraische Geometrie (Grothen\-dieck, 1960er) verall\-gemeinert

dies zu \textbf{Schemata}: Geometrie über beliebigen kommutativen
Ringen, nicht nur über algebraisch abgeschlossenen Körpern.
Die \textbf{étale Kohomologie} und die \textbf{Weil-Ver\-mutungen}
(bewiesen von Deligne, 1974) verbinden Geometrie über endlichen
Körpern mit komplexer Analysis -- ein Kreis, der sich zurück zu den
endlichen Körpern dieser Arbeit schließt.

\subsection*{Darstellungstheorie}

Die in Kapitel 8 eingeführten Gruppencharaktere und
Kompositionsfaktoren sind der Einstieg in die
\textbf{Darstellungstheorie endlicher Gruppen}.
Eine \textbf{Darstellung} einer Gruppe $G$ ist ein
Gruppenhomomorphismus $\rho : G \to GL(V)$ für einen
$K$-Vektorraum $V$.

Die Darstellungstheorie fragt: In welche irreduziblen
Bestandteile zerfällt eine Darstellung?
\begin{itemize}
    \item[-] Für \textbf{endliche Gruppen} über $\mathbb{C}$
          liefert die \textbf{Charaktertheorie} (Frobenius, Schur)
          eine vollständige Antwort: Die Anzahl der irreduziblen
          Darstellungen ist gleich der Anzahl der Konjugationsklassen.
          Die Charaktertafel einer endlichen Gruppe ist eine unitäre
          Matrix -- der erste Beweis, dass Gruppen eine eigene
          ``Fourier-Analysis'' haben.
    \item[-] Für \textbf{kompakte Lie-Gruppen} (wie $SO(3)$, $SU(2)$)
          liefert die \textbf{Theorie der höchsten Gewichte}
          (Cartan, Weyl) eine vollständige Klassifikation der
          irreduziblen Darstellungen. Die Darstellungen von $SU(2)$
          sind genau die Spin-$j$-Darstellungen der Quantenmechanik
          ($j = 0, \frac{1}{2}, 1, \frac{3}{2}, \ldots$).
    \item[-] Die \textbf{Langlands-Korrespondenz} (Langlands-Programm,
          ab 1967) ist das ambitionierteste offene Programm der
          Mathematik: Sie postuliert eine tiefe Entsprechung zwischen
          Darstellungen von Galoisgruppen (Zahlentheorie) und
          Darstellungen von Lie-Gruppen (harmonische Analysis).
          Die Beweise des \textbf{Fermatschen Letzten Satzes}
          (Wiles, 1995) und der \textbf{Sato--Tate-Vermutung}
          (Taylor et al., 2011) waren Spezialfälle des
          Langlands-Programms.
\end{itemize}

\subsection*{Numerische lineare Algebra}

Die theoretischen Resultate dieser Arbeit haben direkte numerische
Entsprechungen, die in der Praxis von zentraler Bedeutung sind:
\begin{itemize}
    \item[-] Die \textbf{LU-Zerlegung} (Gaußscher Algorithmus mit
          Pivotisierung) löst $Ax = b$ in $O(n^3)$ Operationen
          stabil. Die \textbf{Cholesky-Zerlegung} ist die
          effizientere Variante für positiv definite Matrizen.
    \item[-] Der \textbf{QR-Algorithmus} (Francis, 1961) berechnet
          alle Eigenwerte einer Matrix durch iterative Schur-Zerlegung
          in $O(n^3)$. Er ist der Standardalgorithmus für
          Eigenwertprobleme in LAPACK und MATLAB.
    \item[-] Die \textbf{Singulärwertzerlegung} (SVD) ist das
          wichtigste Werkzeug der Datenanalyse (PCA), der
          Bildkompression (JPEG2000) und des maschinellen Lernens
          (latente semantische Analyse, Empfehlungssysteme).
    \item[-] \textbf{Iterative Verfahren} (Krylov-Unterraummethoden,
          konjugierte Gradienten, GMRES) lösen große, dünn besetzte
          lineare Gleichungssysteme in $O(n)$ bis $O(n^2)$ und sind
          unentbehrlich für partielle Differentialgleichungen,
          Simulationen und maschinelles Lernen.
\end{itemize}

\subsection*{Quanteninformation und Quantencomputing}

Die in Kapitel 13 eingeführten Hilberträume, unitären Abbildungen
und der Spektralsatz bilden das vollständige mathematische Fundament
des \textbf{Quantencomputings}:
\begin{itemize}
    \item[-] Ein \textbf{Qubit} ist ein Einheitsvektor in $\mathbb{C}^2$;
          ein $n$-Qubit-System lebt im Tensorprodukt
          $(\mathbb{C}^2)^{\otimes n}$ der Dimension $2^n$.
    \item[-] \textbf{Quantengatter} sind unitäre Abbildungen auf
          diesem Raum. Die Hürde des Quantencomputings ist, dass
          unitäre Operationen durch Dekohärenz (ungewollte Kopplung
          an die Umgebung) gestört werden.
    \item[-] \textbf{Quantenfehlerkorrektur} (Shor, Steane, 1995--96)
          ist die Anwendung linearer Codes (Kapitel 20) auf Qubits:
          Der $[[7,1,3]]$-Steane-Code ist der erste Quanten-Fehler-
          korrektur-Code und ein direktes Quantenanalogon des
          klassischen Hamming-Codes.
    \item[-] Der \textbf{Shor-Algorithmus} (1994) faktorisiert
          $n$-stellige Zahlen in polynomieller Zeit auf einem
          Quantencomputer -- und würde RSA-Verschlüsselung
          (Kapitel 2) brechen. Er beruht auf der
          Quantenfouriertransformation, einer unitären Abbildung,
          die die diskrete Fouriertransformation auf dem
          Vektorraum $\mathbb{C}^{2^n}$ realisiert.
\end{itemize}

\subsection*{Maschinelles Lernen und tiefe neuronale Netze}

Die modernsten Anwendungen der Linearen Algebra finden sich im
\textbf{maschinellen Lernen}:
\begin{itemize}
    \item[-] \textbf{Lineare Regression und PCA} sind direkte
          Anwendungen der SVD und der Pseudoinversen (Kapitel 16).
          Hauptkomponentenanalyse reduziert hochdimensionale Daten
          auf die Richtungen größter Varianz -- die singulären
          Vektoren zum größten singulären Wert.
    \item[-] \textbf{Tiefe neuronale Netze} sind im Kern Kompositionen
          affin-linearer Abbildungen $x \mapsto \sigma(Wx + b)$
          (mit nichtlinearer Aktivierungsfunktion $\sigma$).
          Die Gewichtsmatrizen $W$ sind die Parameter, die durch
          Gradientenabstieg optimiert werden. Die
          \textbf{Backpropagation} ist algorithmisch eine Anwendung
          der Kettenregel auf Kompositionen von Abbildungen --
          linear algebraisch: eine sukzessive Multiplikation von
          Jacobi-Matrizen.
    \item[-] \textbf{Attention-Mechanismen} (Transformer, Vaswani et al.,
          2017) -- die Grundlage großer Sprachmodelle -- sind
          im Kern Skalarprodukt-Operationen zwischen Vektoren
          hoher Dimension: $\mathrm{Attention}(Q, K, V) =
          \mathrm{softmax}\!\left(\frac{QK^t}{\sqrt{d_k}}\right)V$,
          wobei $Q, K, V$ Matrizen sind. Das Skalarprodukt $QK^t$
          misst die ``Ähnlichkeit'' zwischen Query- und Key-Vektoren.
    \item[-] Die \textbf{Geometrie der Dar\-stellungs\-räu\-me}
          (Embedding-Geometrie) -- in der Wortvektoren so
          angeordnet sind, dass semantische Beziehungen
          Vektordifferenzen entsprechen (``König $-$ Mann $+$
          Frau $\approx$ Königin'') -- ist eine empirische
          Entdeckung, die auf tiefen geometrischen Strukturen
          beruht, die noch nicht vollständig verstanden sind.
\end{itemize}

\subsection*{Kryptographie und Post-Quanten-Sicherheit}

Die in Kapitel 2 eingeführten endlichen Körper und das
RSA-Verfahren sind der Einstieg in die moderne Kryptographie:
\begin{itemize}
    \item[-] \textbf{Elliptische Kurven-Kryptographie} (ECC)
          nutzt die Gruppenstruktur elliptischer Kurven über
          endlichen Körpern $\mathbb{F}_p$. Die Sicherheit beruht
          auf dem \textbf{diskreten Logarithmusproblem} auf der
          Kurve, das deutlich schwerer zu lösen ist als der
          gewöhnliche diskrete Logarithmus.
    \item[-] \textbf{Gitterbasierte Kryptographie} (NTRU, CRYSTALS-Kyber)
          ist der Kandidat für Post-Quanten-Kryptographie
          (standardisiert durch NIST, 2022--24). Die Sicherheit
          beruht auf der Schwierigkeit des \textbf{Shortest
          Vector Problem} (SVP) in hochdimensionalen Gittern --
          einem Problem der Geometrie der Zahlen, das direkt
          auf der Theorie der Moduln über $\mathbb{Z}$
          (Kapitel 19) aufbaut.
    \item[-] Der in Kapitel 6 beschriebene
          \textbf{Frobenius-Automorphismus} spielt eine
          zentrale Rolle in der Theorie der
          elliptischen Kurven: Die Anzahl der $\mathbb{F}_p$-Punkte
          einer elliptischen Kurve ist mit dem Frobenius-Endomorphismus
          über den \textbf{Satz von Hasse} verbunden.
\end{itemize}

% ============================================================
\section{Die Einheit der Mathematik}
% ============================================================

Diese Arbeit begann mit der Frage, was ein Träger eines Vektorraums ist: eine linear
unabhängige Menge, die den Raum aufspannt und auf die sich
jeder Vektor eindeutig projizieren lässt. Diese scheinbar
elementare Definition ist der Ursprung eines Begriffsnetzes,
das sich über die gesamte Mathematik und Physik erstreckt.

\begin{beobachtung}[Die Universalität des Vektorraum-Gedankens]
    Der Vektorraum ist die flexibelste algebraische Struktur, die
    Linearität präzise einfängt. Jede Theorie, die Superposition,
    Linearisierung oder lineare Approximation verwendet, ist im
    Kern lineare Algebra:
    \begin{itemize}
        \item[-] Die \textbf{Fourier-Analysis} ist die
              Spektraltheorie des Laplace-Operators auf $L^2$ --
              eine Diagonalisierung in einem unendlichdimensionalen
              Hilbertraum.
        \item[-] Die \textbf{Differentialgeometrie} linearisiert
              Mannigfaltigkeiten durch Tangentialräume -- lokale
              Vektorräume, die das infinitesimale Verhalten tragen.
        \item[-] Die \textbf{Galois-Theorie} studiert Körpererweiterungen
              als Vektorräume und ihre Automorphismengruppen -- eine
              direkte Synthese von Linearer Algebra und Gruppentheorie.
        \item[-] Die \textbf{homologische Algebra} betrachtet Komplexe
              linearer Abbildungen $\cdots \to V_{n+1} \to V_n \to
              V_{n-1} \to \cdots$ und misst in den Kohomologiegruppen,
              wie weit diese Abbildungen davon entfernt sind, exakt
              zu sein.
    \end{itemize}
    In allen Fällen ist das Kernwerkzeug dasselbe: ein Vektorraum
    mit einem ausgezeichneten Träger, und die Frage, welche
    lineare Struktur dieser Träger trägt.
\end{beobachtung}

\begin{beobachtung}[Projektion als Grundprinzip]
    Die Wahl des Begriffs ``Träger'' in dieser Arbeit
    hat einen konzeptuellen Kern: Ein Träger ist die Menge derjenigen
    Vektoren, auf die projiziert wird. Die Projektion
    $v \mapsto P_U v$ auf einen Unterraum $U$ ist das fundamentale
    Analysewerkzeug -- sie zerlegt jeden Vektor in seinen
    ``wesentlichen'' und seinen ``unwesentlichen'' Anteil
    (bezüglich $U$).

    Dieses Prinzip durchzieht die gesamte Arbeit:
    \begin{itemize}
        \item[-] Die \textbf{Koordinatendarstellung} eines Vektors
              bezüglich eines Trägers $[t_1, \ldots, t_n]$ ist die
              Projektion auf jeden einzelnen Trägervektor.
        \item[-] Der \textbf{Spektralsatz} zerlegt einen Operator
              in eine Summe von Projektionen auf die Eigenräume.
        \item[-] Die \textbf{Polarzerlegung} $A = UH$ projiziert
              eine lineare Abbildung in ihre unitäre und ihre
              ``streckende'' Komponente.
        \item[-] Die \textbf{singuläre Wertzerlegung} projiziert
              eine Matrix auf ihre $k$ wesentlichsten
              Richtungen (Niedrigrangapproximation).
        \item[-] Die \textbf{Fourier-Zerlegung} projiziert eine
              Funktion auf die Frequenzkomponenten des Trägers
              $[e^{in\theta}]_{n \in \mathbb{Z}}$ des Hilbertraums
              $L^2(S^1)$.
    \end{itemize}
    Das Konzept des Trägers -- als Menge der Projektionsziele --
    ist damit nicht nur ein terminologisches Element dieser Arbeit,
    sondern ein strukturelles Leitprinzip: Mathematik ist zu
    wesentlichen Teilen die Kunst, den richtigen Träger zu wählen,
    auf den projiziert wird.
\end{beobachtung}

\begin{beobachtung}[Ausblick: Offene Fragen]
    Trotz der enormen Entwicklung der Linearen Algebra und ihrer
    Nachbargebiete in den letzten zwei Jahrhunderten gibt es
    grundlegende offene Fragen:
    \begin{itemize}
        \item[-] \textbf{Komplexitätstheorie der Matrixmultiplikation}:
              Der schnellste bekannte Algorithmus für die
              Multiplikation zweier $n \times n$-Matrizen hat
              Laufzeit $O(n^\omega)$ mit $\omega < 2{,}372$ (Williams et al., 2024).
              Ob $\omega = 2$ möglich ist (d.\,h.\ Matrixmultiplikation
              in quasilinearer Zeit), ist offen -- eine der wichtigsten
              Fragen der theoretischen Informatik.
        \item[-] \textbf{Permanente und $\#P$-Voll\-stän\-dig\-keit}:
              Das Permanente einer $0$-$1$-Matrix -- definiert wie
              die Determinante, aber ohne Vorzeichen --
              ist $\#P$-vollständig zu berechnen (Valiant, 1979).
              Warum ist die Determinante (mit Vorzeichen) in $O(n^3)$
              berechenbar, das Permanente (ohne Vorzeichen) aber
              (vermutlich) nicht? Dieser Unterschied hängt eng mit
              der $P \neq NP$-Frage zusammen.
        \item[-] \textbf{Geodäten auf Lie-Gruppen und optimale
              Steuerung}: Die Frage nach kürzesten Wegen auf $SO(n)$
              und anderen Lie-Gruppen (im Sinne des Riemannschen
              Abstands) ist ein aktives Forschungsgebiet mit
              Anwendungen in der Robotik, der Bildverarbeitung
              und der Quantensteuerung.
        \item[-] \textbf{Spektrallücken und Zufallsmatrizen}:
              Die \textbf{Zufallsmatrix-Theorie} (Wigner, Dyson,
              Mehta) beschreibt das statistische Verhalten der
              Eigenwerte großer zufälliger Matrizen. Das
              \textbf{Wigner-Halbkreisgesetz} und die
              \textbf{GUE-Universalität} sind tiefe Resultate
              über die Verteilung der Eigenwertabstände --
              mit Verbindungen zur Zahlentheorie (Nullstellen
              der Riemannschen Zeta-Funktion) und zur Physik
              (Niveauabstände komplexer Quantensysteme).
    \end{itemize}
\end{beobachtung}

Die Lineare Algebra, die in dieser Arbeit entwickelt wurde, ist
kein abgeschlossenes Gebäude, sondern ein lebendiger Organismus:
Ihre Begriffe -- Vektorraum, Träger, lineare Abbildung,
Skalarprodukt, Eigenwert, Determinante -- bilden die Grundsprache
der gesamten modernen Mathematik und Physik. Wer diese Sprache
beherrscht, hat den Schlüssel zu Quantenmechanik und
Relativitätstheorie, zu Codierungstheorie und Kryptographie,
zu numerischer Simulation und maschinellem Lernen,
zu algebraischer Geometrie und Darstellungstheorie.

Der Träger dieser Sprache -- um in der Terminologie dieser Arbeit
zu bleiben -- sind die fundamentalen Sätze: der Rangsatz,
der Spektralsatz, der Homomorphiesatz, der Satz von Cayley--Hamilton,
der Satz von Jordan--Hölder, der Satz von Perron--Frobenius.
Auf diese Sätze projiziert sich die gesamte Mathematik --
und von ihnen aus öffnet sich der Ausblick auf alle noch
nicht beschrittenen Wege.

\begin{thebibliography}{9}

\bibitem{huppert}
B.~Huppert und W.~Willems,
\textit{Lineare Algebra},
B.\,G.~Teubner Verlag, Wiesbaden, 2006.

\bibitem{lang}
S.~Lang,
\textit{Linear Algebra},
3.\,Aufl., Springer, New York, 1987.

\bibitem{strang}
G.~Strang,
\textit{Introduction to Linear Algebra},
5.\,Aufl., Wellesley-Cambridge Press, 2016.

\bibitem{fischer}
G.~Fischer,
\textit{Lineare Algebra},
18.\,Aufl., Springer Spektrum, 2014.

\end{thebibliography}

\end{document}

